% Lesson 2 — Grammar: Modals: can/may (permission), can/could (possibility), ought to/ought not to, need(n't), don't have to, should/ought to, must + deduction — Anglais 1ère A — Cours signé OnBuch+
% Programme officiel MINESEC. Compiler avec : tectonic ENA02-grammar-modals-can-may-permission-can-could-possibility.tex
\newcommand{\DOCMATIERE}{Anglais}
\newcommand{\DOCNIVEAU}{1ère A}
\newcommand{\DOCMODULE}{Chapter 1 : Family and Social Life}
\newcommand{\DOCLECON}{ENA02}
\newcommand{\DOCTITRE}{Lesson 2 — Grammar: Modals: can/may (permission), can/could (possibility), ought to/ought not to, need(n't), don't have to, should/ought to, must + deduction}
% OnBuch+ — préambule « cours signés » ENRICHI (compilé avec tectonic / XeLaTeX)
% Système visuel « Pop » d'OnBuch+ : fond crème (jamais blanc), bordures encre
% épaisses, ombres « dures » décalées, titres Archivo Black en capitales.
% Version MATHÉMATIQUES : graphes (pgfplots/tikz), boîtes pédagogiques
% (définition, propriété, méthode, attention, le savais-tu, exercice résolu,
% exercice, TP, bilan), page de garde et signature OnBuch+ ancrée en bas.
%
% Macros attendues dans main.tex AVANT \input{preamble} :
%   \DOCMATIERE  \DOCNIVEAU  \DOCMODULE  \DOCLECON (numéro)  \DOCTITRE
\documentclass[11pt]{article}

\usepackage{fontspec}
\usepackage[english,french]{babel} % français = langue principale ; l'anglais via \eng{..} / textebox
\usepackage[a4paper,margin=2cm,top=3.4cm,bottom=2.8cm,headheight=1.4cm,footskip=1.3cm]{geometry}
\usepackage{amsmath,amssymb,amsfonts}
\usepackage{mathtools} % \xmapsto, \coloneqq... souvent utilisés par les modèles
\usepackage{cancel} % simplifications barrées (bilans de forces)
% \abs{z} (module, valeur absolue) et \norm{u} (norme), utilisés par les modèles
\providecommand{\abs}{}\let\abs\relax\DeclarePairedDelimiter{\abs}{\lvert}{\rvert}
\providecommand{\norm}{}\let\norm\relax\DeclarePairedDelimiter{\norm}{\lVert}{\rVert}
\usepackage[table]{xcolor} % [table] : fournit aussi \rowcolors (alternance de lignes)
\usepackage{pagecolor}
\usepackage{tikz}
\usetikzlibrary{arrows.meta,positioning,calc,shapes.geometric,shapes.misc,shapes.arrows,decorations.pathmorphing,decorations.pathreplacing,decorations.markings,patterns,shadows,mindmap,trees,fit,backgrounds,matrix,intersections,angles,quotes}
\usepackage{pgfplots}
\usepackage[version=4]{mhchem} % équations nucléaires \ce{^{235}_{92}U}
\usepackage[european,siunitx]{circuitikz} % schémas électriques
\pgfplotsset{compat=1.18}
\usepgfplotslibrary{fillbetween,groupplots} % aires entre courbes (calcul intégral)
\usepackage{enumitem}
\usepackage{fancyhdr}
% [most] charge toutes les bibliothèques tcolorbox (listings, minted, raster,
% documentation...) et gonfle inutilement la mémoire TeX sur un document long
% (~300 boîtes) : on ne charge que ce qui est réellement utilisé.
\usepackage[skins,breakable]{tcolorbox}
\usepackage{graphicx}
\usepackage{colortbl}
\usepackage{booktabs}
\usepackage{tabularx}
\usepackage{diagbox} % case d'en-tête diagonale des tableaux à double entrée
% Colonnes extensibles centrée / gauche / droite (C, L, R), souvent utilisées par les modèles
\newcolumntype{C}{>{\centering\arraybackslash}X}
\newcolumntype{L}{>{\raggedright\arraybackslash}X}
\newcolumntype{R}{>{\raggedleft\arraybackslash}X}
\usepackage{array}
\usepackage{multicol}
\usepackage{multirow}
\usepackage{siunitx}
% parse-numbers=false : accepte aussi \SI{2,5 \times 10^{-3}}{...}
\sisetup{locale=FR,output-decimal-marker={,},group-minimum-digits=4,parse-numbers=false}
\DeclareSIUnit\ohm{\ensuremath{\symup{\Omega}}} % U+2126 absent de la police
\DeclareSIUnit\division{div} % réglages d'oscilloscope (ms/div, V/div)
\DeclareSIUnit\tr{tr} % tours (vitesse de rotation, ex. \SI{1500}{\tr\per\minute})
\DeclareSIUnit\molar{M} % concentration molaire (ex. \SI{3}{\molar}, \SI{1}{\micro\molar})
\DeclareSIUnit\an{an} % année (le modèle confond parfois avec \year, un compteur TeX) — ex. \SI{5}{\kilo\meter\per\an}
\DeclareSIUnit\FCFA{FCFA} % franc CFA (ex. \SI{500}{\FCFA\per\kilo\gram})
\DeclareSIUnit\degreeBrix{\textdegree Brix} % teneur en sucre d'un sirop/jus (réfractométrie)
\DeclareSIUnit\month{mois} % le modèle utilise parfois \month, un compteur TeX, comme unité de durée
\usepackage{qrcode}
% \degree : commande fréquemment utilisée par les modèles pour un angle ou une
% température en dehors de \SI{}{} (non fournie par les paquets chargés).
\providecommand{\degree}{^{\circ}}
% \qed (fin de démonstration), utilisé par les modèles sans amsthm
\providecommand{\qed}{\hfill$\square$}
% \barre : double barre des valeurs interdites dans un tableau de variations (tkz-tab)
\providecommand{\barre}{\ensuremath{\|}}
% environnement proof (amsthm non chargé) utilisé par les modèles
\ifdefined\proof\else\newenvironment{proof}[1][Démonstration]{\par\noindent\textit{#1.}\ }{\qed\par}\fi
\usepackage{lastpage}
\usepackage{hyperref}
\hypersetup{hidelinks}

% ============================================================
% Palette « Pop » OnBuch+ — mêmes hex que la classe P (plus_ui.dart)
% ============================================================
\definecolor{poporange}{HTML}{F59321}
\definecolor{popdark}{HTML}{E07A0C}
\definecolor{popcream}{HTML}{FFF6E8}
\definecolor{popink}{HTML}{1C1714}
\definecolor{poppurple}{HTML}{7A5AE0}
\definecolor{popgreen}{HTML}{2E9E6A}
\definecolor{poppink}{HTML}{E0457B}
\definecolor{popblue}{HTML}{2F7DE1}
\definecolor{popgold}{HTML}{C98A12}
\definecolor{popmuted}{HTML}{8A7F76}
\definecolor{popline}{HTML}{EADFD2}
\definecolor{popgray}{HTML}{8A7F76}
\colorlet{popgrey}{popgray}
% Alias tolérants (le modèle nomme parfois les couleurs différemment)
\colorlet{popwhite}{white}
\colorlet{popblack}{popink}
\colorlet{popred}{poppink}
\colorlet{popyellow}{popgold}
% Teintes claires (fonds de tableaux/encadrés) — le modèle nomme parfois
% les couleurs avec un suffixe L (ex. popblueL) sans les avoir définies.
\colorlet{poporangeL}{poporange!20}
\colorlet{popdarkL}{popdark!20}
\colorlet{poppurpleL}{poppurple!20}
\colorlet{popgreenL}{popgreen!20}
\colorlet{poppinkL}{poppink!20}
\colorlet{popblueL}{popblue!20}
\colorlet{popgoldL}{popgold!20}
\colorlet{popredL}{popred!20}
\colorlet{popgrayL}{popgray!20}
% Teintes claires pour fonds de tableaux / zones
\colorlet{popblueL}{popblue!10!white}
\colorlet{poporangeL}{poporange!14!white}
\colorlet{popgreenL}{popgreen!12!white}
\colorlet{poppurpleL}{poppurple!12!white}
\colorlet{poppinkL}{poppink!12!white}
\colorlet{popgoldL}{popgold!14!white}

\pagecolor{popcream}

% ============================================================
% Typographie
% ============================================================
\defaultfontfeatures{Ligatures=TeX}
\setmainfont{PlusJakartaSans-Medium.ttf}[Path=../../../fonts/,BoldFont=PlusJakartaSans-Bold.ttf]
% Symboles, API et emojis absents de la police principale : repli sur DejaVu Sans (caractères actifs)
\newfontface\symfont{DejaVuSans.ttf}[Path=../../../fonts/]
\makeatletter
\newcommand\symactive[1]{\catcode#1=13 \begingroup\lccode`\~=#1 \lowercase{\endgroup\def~}{{\symfont\char#1}}}
\newcommand\symblank[1]{\catcode#1=13 \begingroup\lccode`\~=#1 \lowercase{\endgroup\def~}{}}
\newcommand\symhyph[1]{\catcode#1=13 \begingroup\lccode`\~=#1 \lowercase{\endgroup\def~}{\mbox{-}}}
\newcount\symc
\newcommand\symrange[2]{\symc=#1 \@whilenum\symc<#2 \do{\expandafter\symactive\expandafter{\the\symc}\advance\symc by 1 }}
\newcommand\symblankrange[2]{\symc=#1 \@whilenum\symc<#2 \do{\expandafter\symblank\expandafter{\the\symc}\advance\symc by 1 }}
\makeatother
\symrange{"0250}{"02FF}\symactive{"2713}\symactive{"2714}\symactive{"2717}\symactive{"2718}\symactive{"2610}\symactive{"2611}\symactive{"2612}
\symhyph{"2011}\symhyph{"2E1A}\symblankrange{"1F300}{"1FAFF}\symblank{"FE0F}
% Maths en sans-serif (Fira Math) avec chiffres et lettres droites de Plus
% Jakarta, pour que les formules se fondent dans le texte.
\usepackage[math-style=ISO,bold-style=ISO]{unicode-math}
\setmathfont{FiraMath-Regular.otf}[Path=../../../fonts/]
\setmathfont{PlusJakartaSans-Medium.ttf}[Path=../../../fonts/,range={up/num,up/{Latin,latin},\mathup}]
\usepackage{icomma} % virgule décimale sans espace en mode math
% Caractères mathématiques Unicode absents des polices de texte (Plus Jakarta,
% Archivo Black) : rendus par leur équivalent mathématique, en texte comme en
% formule — sinon ils s'affichent en carré vide (ex. « Arithmétique dans ℤ »).
\usepackage{newunicodechar}
\newunicodechar{ℝ}{\ensuremath{\mathbb{R}}}
\newunicodechar{ℤ}{\ensuremath{\mathbb{Z}}}
\newunicodechar{ℂ}{\ensuremath{\mathbb{C}}}
\newunicodechar{ℕ}{\ensuremath{\mathbb{N}}}
\newunicodechar{ℚ}{\ensuremath{\mathbb{Q}}}
\newunicodechar{𝔻}{\ensuremath{\mathbb{D}}}
\newunicodechar{α}{\ensuremath{\alpha}}
\newunicodechar{β}{\ensuremath{\beta}}
\newunicodechar{γ}{\ensuremath{\gamma}}
\newunicodechar{δ}{\ensuremath{\delta}}
\newunicodechar{ε}{\ensuremath{\varepsilon}}
\newunicodechar{θ}{\ensuremath{\theta}}
\newunicodechar{λ}{\ensuremath{\lambda}}
\newunicodechar{μ}{\ensuremath{\mu}}
\newunicodechar{σ}{\ensuremath{\sigma}}
\newunicodechar{τ}{\ensuremath{\tau}}
\newunicodechar{φ}{\ensuremath{\varphi}}
\newunicodechar{ψ}{\ensuremath{\psi}}
\newunicodechar{ω}{\ensuremath{\omega}}
\newunicodechar{Σ}{\ensuremath{\Sigma}}
% majuscules produites par \MakeUppercase dans les titres (α-aminés -> Α)
\newunicodechar{Α}{\ensuremath{\alpha}}
\newunicodechar{Β}{\ensuremath{\beta}}
\newunicodechar{Γ}{\ensuremath{\Gamma}}
\newunicodechar{Δ}{\ensuremath{\Delta}}
\newunicodechar{Ε}{\ensuremath{\varepsilon}}
\newunicodechar{Θ}{\ensuremath{\Theta}}
\newunicodechar{Λ}{\ensuremath{\Lambda}}
\newunicodechar{Μ}{\ensuremath{\mu}}
\newunicodechar{Τ}{\ensuremath{\tau}}
\newunicodechar{Φ}{\ensuremath{\Phi}}
\newunicodechar{Ψ}{\ensuremath{\Psi}}
\newunicodechar{Ω}{\ensuremath{\Omega}}
\newunicodechar{↦}{\ensuremath{\mapsto}}
\newunicodechar{∈}{\ensuremath{\in}}
\newunicodechar{∉}{\ensuremath{\notin}}
\newunicodechar{∧}{\ensuremath{\wedge}}
\newunicodechar{⇔}{\ensuremath{\Leftrightarrow}}
\newunicodechar{⟺}{\ensuremath{\Longleftrightarrow}}
\newunicodechar{⇒}{\ensuremath{\Rightarrow}}
\newunicodechar{⟹}{\ensuremath{\Longrightarrow}}
\newunicodechar{∘}{\ensuremath{\circ}}
\newunicodechar{∪}{\ensuremath{\cup}}
\newunicodechar{∩}{\ensuremath{\cap}}
\newunicodechar{≡}{\ensuremath{\equiv}}
\newunicodechar{⊂}{\ensuremath{\subset}}
\newunicodechar{⊥}{\ensuremath{\perp}}
\newunicodechar{∃}{\ensuremath{\exists}}
\newunicodechar{∀}{\ensuremath{\forall}}
\newunicodechar{∛}{\ensuremath{\sqrt[3]{\,}}}
\newunicodechar{∜}{\ensuremath{\sqrt[4]{\,}}}
\newunicodechar{⃗}{\ensuremath{{}^{\to}}}
\newunicodechar{ⁿ}{\ifmmode^{n}\else\textsuperscript{\ensuremath{n}}\fi}
\newunicodechar{ˣ}{\ifmmode^{x}\else\textsuperscript{\ensuremath{x}}\fi}
\newunicodechar{ᵇ}{\ifmmode^{b}\else\textsuperscript{\ensuremath{b}}\fi}
\newunicodechar{ʳ}{\ifmmode^{r}\else\textsuperscript{\ensuremath{r}}\fi}
\newunicodechar{ᵘ}{\ifmmode^{u}\else\textsuperscript{\ensuremath{u}}\fi}
\newunicodechar{ʸ}{\ifmmode^{y}\else\textsuperscript{\ensuremath{y}}\fi}
\newunicodechar{ᵃ}{\ifmmode^{a}\else\textsuperscript{\ensuremath{a}}\fi}
\newunicodechar{ᵉ}{\ifmmode^{e}\else\textsuperscript{\ensuremath{e}}\fi}
\newunicodechar{ᵏ}{\ifmmode^{k}\else\textsuperscript{\ensuremath{k}}\fi}
\newunicodechar{ᵖ}{\ifmmode^{p}\else\textsuperscript{\ensuremath{p}}\fi}
\newunicodechar{ᴺ}{\ifmmode^{N}\else\textsuperscript{\ensuremath{N}}\fi}
\newunicodechar{ᶜ}{\ifmmode^{c}\else\textsuperscript{\ensuremath{c}}\fi}
\newunicodechar{ʰ}{\ifmmode^{h}\else\textsuperscript{\ensuremath{h}}\fi}
\newunicodechar{ᵠ}{\ifmmode^{q}\else\textsuperscript{\ensuremath{q}}\fi}
\newunicodechar{ₙ}{\ifmmode_{n}\else\textsubscript{\ensuremath{n}}\fi}
\newunicodechar{ₐ}{\ifmmode_{a}\else\textsubscript{\ensuremath{a}}\fi}
\newunicodechar{ᵢ}{\ifmmode_{i}\else\textsubscript{\ensuremath{i}}\fi}
\newunicodechar{ᵣ}{\ifmmode_{r}\else\textsubscript{\ensuremath{r}}\fi}
\newunicodechar{ₖ}{\ifmmode_{k}\else\textsubscript{\ensuremath{k}}\fi}
\newunicodechar{ᵦ}{\ifmmode_{\beta}\else\textsubscript{\ensuremath{\beta}}\fi}
\newunicodechar{ₓ}{\ifmmode_{x}\else\textsubscript{\ensuremath{x}}\fi}
\newfontfamily\popdisplay{ArchivoBlack-Regular.ttf}[Path=../../../fonts/]
\newfontfamily\poplogo{SpaceGrotesk-Bold.ttf}[Path=../../../fonts/]
\newfontfamily\popsemi{PlusJakartaSans-SemiBold.ttf}[Path=../../../fonts/]
\newfontfamily\popxbold{PlusJakartaSans-ExtraBold.ttf}[Path=../../../fonts/]
\newfontfamily\popxboldls{PlusJakartaSans-ExtraBold.ttf}[Path=../../../fonts/,LetterSpace=3.0]

\setlist[enumerate]{leftmargin=1.7em,itemsep=5pt,topsep=4pt}
\setlist[itemize]{leftmargin=1.7em,itemsep=4pt,topsep=4pt,label={\color{poporange}\textbullet}}
\setlength{\parindent}{0pt}
\setlength{\parskip}{5pt}
\renewcommand{\arraystretch}{1.25}

% pgfplots : style Pop par défaut
\pgfplotsset{
  every axis/.append style={axis line style={popink,line width=1pt},tick style={popink},
    grid style={popline,line width=0.5pt},label style={font=\small},tick label style={font=\footnotesize}},
  popcurve/.style={poporange,line width=1.8pt,no markers,smooth},
}
% Même style disponible dans un \draw TikZ simple (hors pgfplots)
\tikzset{popcurve/.style={poporange,line width=1.8pt}}
\tikzset{popgrid/.style={popline,very thin}}
\colorlet{popcurve}{poporange} % le nom du style est parfois utilisé comme couleur

% ============================================================
% Pill / squiggle
% ============================================================
\newcommand{\poppill}[3]{%
  \tikz[baseline=-0.55ex]\node[draw=popink,line width=1.2pt,rounded corners=2.6pt,%
    fill=#2,inner xsep=6.5pt,inner ysep=3pt]{%
    \popxboldls\fontsize{7}{7}\selectfont\color{#3}\MakeUppercase{#1}};%
}
\newcommand{\popsquiggle}[1]{%
  \tikz[baseline=-0.4ex]{
    \draw[#1,line width=2.6pt,line cap=round]
      plot [smooth] coordinates {(0,0.09) (0.34,0.01) (0.68,0.21) (1.02,0.01) (1.36,0.10)};
  }%
}
% Mot-clé surligné (marqueur orange sous le texte)
\newcommand{\cle}[1]{{\popxbold\color{popdark}#1}}

% ============================================================
% En-tête / pied
% ============================================================
\pagestyle{fancy}
\fancyhf{}
\renewcommand{\headrulewidth}{0pt}
\renewcommand{\footrulewidth}{0pt}
\lhead{\raisebox{-0.15ex}{\poplogo\fontsize{13}{13}\selectfont\color{popink}OnBuch\color{poporange}+}}
\rhead{\poppill{\DOCMATIERE\ \textbullet\ \DOCNIVEAU\ \textbullet\ Leçon \DOCLECON}{white}{popink}}
\lfoot{\popxbold\fontsize{7.6}{7.6}\selectfont\color{popmuted}\MakeUppercase{Cours signé OnBuch+}\ \textbullet\ {\popsemi\DOCTITRE}}
\rfoot{\tikz[baseline=-0.6ex]\node[rounded corners=8.5pt,fill=popink,minimum height=17pt,minimum width=17pt,inner xsep=5pt,inner sep=0]{\popxbold\fontsize{8}{8}\selectfont\color{popcream}\thepage\,/\,\pageref*{LastPage}};}

% ============================================================
% Titres
% ============================================================
\newcommand{\chaptitre}[2]{%
  \begin{tcolorbox}[enhanced,colback=poporange,colframe=popink,boxrule=2.6pt,arc=11pt,
      drop shadow=popink,left=15pt,right=15pt,top=12pt,bottom=11pt,before skip=3pt,after skip=18pt]
    \poppill{#2}{popink}{popcream}\par\vspace{9pt}
    {\raggedright\hyphenpenalty=10000\exhyphenpenalty=10000\popdisplay\fontsize{22}{24}\selectfont\color{popcream}\MakeUppercase{#1}\par}
  \end{tcolorbox}
}
% Section numérotée : pastille orange avec le numéro + titre Archivo
\newcounter{popsec}
\newcommand{\coursec}[1]{%
  \stepcounter{popsec}\setcounter{popsub}{0}%
  \par\vspace{16pt}\noindent
  \begin{minipage}{\linewidth}
  \tikz[baseline=-0.7ex]\node[draw=popink,line width=1.6pt,fill=poporange,rounded corners=4pt,minimum size=22pt,inner sep=0]{\popdisplay\fontsize{12}{12}\selectfont\color{popcream}\thepopsec};%
  \hspace{8pt}{\popdisplay\fontsize{15}{17}\selectfont\color{popink}\MakeUppercase{#1}}\par
  \vspace{4pt}\noindent{\color{poporange}\rule{36pt}{3.6pt}}
  \end{minipage}\par\nopagebreak\vspace{8pt}
}
\newcounter{popsub}
\newcommand{\courssub}[1]{%
  \stepcounter{popsub}%
  \par\vspace{9pt}\noindent{\popxbold\fontsize{11.5}{13}\selectfont\color{popink}{\color{poporange}\thepopsec.\thepopsub}\hspace{6pt}#1}\par\nopagebreak\vspace{4pt}
}

% ============================================================
% Boîtes pédagogiques — même recette : fond blanc, bordure épaisse colorée,
% ombre dure encre, badge plein. Titre optionnel [#1].
% ============================================================
\tcbset{popbase/.style={breakable,enhanced,colback=white,boxrule=2.3pt,arc=9pt,
  shadow={3pt}{-3pt}{0pt}{popink},left=14pt,right=14pt,top=11pt,bottom=11pt,before skip=11pt,after skip=15pt,
  fontupper=\popsemi\fontsize{10.6}{14.5}\selectfont\color{popink},
  fontlower=\popsemi\fontsize{10.6}{14.5}\selectfont\color{popink}}}
% Coche dessinée (le glyphe ✓ n'existe pas dans les polices du document)
\newcommand{\popcheck}[1]{\tikz[baseline=-0.5ex]\draw[#1,line width=1.6pt,line cap=round,line join=round](-0.13,0.0)--(-0.04,-0.09)--(0.13,0.1);}
\providecommand{\coche}{\popcheck{popgreen}}
\providecommand{\resultat}[1]{\par\smallskip\noindent\fcolorbox{popgreen}{popgreenL}{\parbox{\dimexpr\linewidth-2\fboxsep-2\fboxrule\relax}{\textbf{Résultat.} #1}}\par\smallskip}
\newcommand{\popboxhead}[3]{\poppill{#1}{#2}{white}\ifx\relax#3\relax\else\hspace{6pt}{\popxbold\fontsize{10.6}{12}\selectfont\color{popink}#3}\fi\par\vspace{7pt}}

\newtcolorbox{aretenir}[1][]{popbase,colframe=popgreen,before upper={\popboxhead{À retenir}{popgreen}{#1}}}
% Texte en anglais (document de lecture, dialogue, extrait) : cadre
% d'étude. Un titre optionnel (source, auteur) s'affiche dans l'en-tête.
\newtcolorbox{textebox}[1][]{popbase,colframe=popblue,before upper={\popboxhead{Text}{popblue}{#1}\begin{otherlanguage}{english}},after upper={\end{otherlanguage}}}
% Marques ✓ / ✗ (forme correcte / fautive) souvent utilisées par les modèles
\providecommand{\cross}{\textcolor{poppink}{\ensuremath{\times}}}
\providecommand{\xmark}{\cross}\providecommand{\crossmark}{\cross}\providecommand{\cmark}{\ensuremath{\checkmark}}
% Ligne de réponse / justification à compléter (inventée par les modèles)
\providecommand{\justification}[1]{\par\noindent\textit{Justification~:} \dotfill\par}
% \textipa (package tipa non chargé) : la police ne couvre pas l'API -> simple passage
\providecommand{\textipa}[1]{#1}
% Soulignement (ulem non chargé) : \uline, \ul
\providecommand{\uline}[1]{\underline{#1}}\providecommand{\ul}[1]{\underline{#1}}
% \glossaire{\item ...} inventée par les modèles : liste de glossaire
\providecommand{\glossaire}[1]{\begin{itemize}#1\end{itemize}}
% \alert (beamer) inventée par les modèles : texte mis en évidence
\providecommand{\alert}[1]{\textbf{\textcolor{poppink}{#1}}}
% variantes ASCII d'ordinaux écrites en macro
\providecommand{\iere}{\textsuperscript{re}}\providecommand{\ieme}{\textsuperscript{e}}\providecommand{\ere}{\textsuperscript{re}}\providecommand{\eme}{\textsuperscript{e}}\providecommand{\er}{\textsuperscript{er}}\providecommand{\ie}{\textsuperscript{e}}
% Ordinaux français écrits en macro par les modèles : 1\ière, 3\ième
\newcommand{\ière}{\textsuperscript{re}}\newcommand{\ième}{\textsuperscript{e}}
% \exemple (inventée par les modèles) : étiquette « Exemple : »
\providecommand{\exemple}{\textbf{Exemple~:}\ }
% \square utilisable aussi hors mode math (cases à cocher)
\AtBeginDocument{\let\squareMath\square\renewcommand{\square}{\ensuremath{\squareMath}}}
% Mot ou phrase en anglais à l'intérieur d'un paragraphe français
\newcommand{\eng}[1]{\ifmmode\text{\textit{#1}}\else\begingroup\selectlanguage{english}\itshape #1\endgroup\fi}
\let\esp\eng % alias toléré
% flèches d'intonation inventées par les modèles
\providecommand{\textsearrow}{}\renewcommand{\textsearrow}{\ensuremath{\searrow}}\providecommand{\textnearrow}{}\renewcommand{\textnearrow}{\ensuremath{\nearrow}}
% \fr{..} : mot français (inventé par les modèles)
\providecommand{\fr}[1]{\textit{#1}}
\newtcolorbox{exemplebox}[1][]{popbase,colframe=poporange,before upper={\popboxhead{Exemple}{poporange}{#1}}}
\newtcolorbox{definition}[1][]{popbase,colframe=popblue,before upper={\popboxhead{Définition}{popblue}{#1}}}
\newtcolorbox{propriete}[1][]{popbase,colframe=poppurple,before upper={\popboxhead{Propriété}{poppurple}{#1}}}
\newtcolorbox{methode}[1][]{popbase,colframe=popgold,before upper={\popboxhead{Méthode}{popgold}{#1}}}
\newtcolorbox{attention}[1][]{popbase,colframe=poppink,before upper={\popboxhead{Attention}{poppink}{#1}}}
\newtcolorbox{experience}[1][]{popbase,colframe=popblue,colback=popblueL,before upper={\popboxhead{Expérience}{popblue}{#1}}}
% « Le savais-tu ? » : carte violette pleine (contexte, culture, Cameroun)
\newtcolorbox{savaistu}[1][]{popbase,colback=poppurple,colframe=popink,
  fontupper=\popsemi\fontsize{10.4}{14.2}\selectfont\color{white},
  before upper={\poppill{Le savais-tu\,?}{white}{poppurple}\ifx\relax#1\relax\else\hspace{6pt}{\popxbold\color{white}#1}\fi\par\vspace{7pt}}}
% Exercice résolu : énoncé en haut, \tcblower puis la solution sur fond crème
\newcounter{popexo}\newcounter{popexr}
\newtcolorbox{exoresolu}[1][]{popbase,colframe=poporange,colbacklower=popcream,
  segmentation style={popink,line width=1.2pt,dashed},
  before upper={\stepcounter{popexr}\popboxhead{Exercice résolu \thepopexr}{poporange}{#1}},
  before lower={\poppill{Solution}{popink}{popcream}\par\vspace{6pt}}}
% Exercice (non résolu en place) — la correction est dans la partie Corrigés
\newtcolorbox{exercice}[1][]{popbase,colframe=popink,
  before upper={\stepcounter{popexo}\popboxhead{Exercice \thepopexo}{popink}{#1}}}
% Corrigé
\newtcolorbox{corrige}[1][]{popbase,colframe=popgreen,colback=popgreenL,
  before upper={\popboxhead{Corrigé}{popgreen}{#1}}}
% Objectifs / prérequis / bilan
\newtcolorbox{objectifs}{popbase,colframe=popink,colback=poporangeL,
  before upper={\popboxhead{Objectifs de la leçon}{popink}{}}}
\newtcolorbox{prerequis}{popbase,colframe=popink,colback=popblueL,
  before upper={\popboxhead{Prérequis}{popblue}{}}}
\newtcolorbox{bilan}{popbase,colframe=popink,colback=white,boxrule=2.8pt,
  before upper={\popboxhead{Fiche bilan}{poporange}{L'essentiel à connaître pour l'examen}}}
% Figure encadrée avec légende
\newtcolorbox{popfigure}[1][]{enhanced,breakable=false,colback=white,colframe=popline,boxrule=1.4pt,arc=8pt,
  left=10pt,right=10pt,top=10pt,bottom=8pt,before skip=10pt,after skip=12pt,halign=center,
  lower separated=false,halign lower=center,
  fontlower=\popsemi\fontsize{9}{11.5}\selectfont\color{popmuted},
  #1}
\newcommand{\legende}[1]{\par\vspace{6pt}{\popsemi\fontsize{9}{11.5}\selectfont\color{popmuted}#1}}

% ============================================================
% Signature OnBuch+
% ============================================================
\newcommand{\poplogoinline}[1]{{\poplogo\fontsize{#1}{#1}\selectfont\color{popink}OnBuch\color{poporange}+}}
\newcommand{\signatureonbuch}{%
  \begin{tcolorbox}[enhanced,colback=popink,colframe=popink,boxrule=2.2pt,arc=10pt,
      drop shadow=poporange,left=14pt,right=14pt,top=10pt,bottom=10pt,before skip=0pt,after skip=0pt]
    \tikz[baseline=-0.7ex]\node[circle,fill=popgreen,draw=popcream,line width=1.3pt,minimum size=22pt,inner sep=0]{\popcheck{white}};%
    \hspace{9pt}\begin{minipage}[c]{0.62\linewidth}
      {\popxbold\fontsize{12}{13}\selectfont\color{popcream}Signé\ \textbullet\ {\poplogo OnBuch\color{poporange}+}}\par\vspace{2pt}
      {\raggedright\popsemi\fontsize{8}{10.5}\selectfont\color{popline}\MakeUppercase{Équipe pédagogique OnBuch+}\\\MakeUppercase{Conforme au programme officiel MINESEC}\par}
    \end{minipage}\hfill
    {\poplogo\fontsize{22}{22}\selectfont\color{popcream}OnBuch\color{poporange}+}
  \end{tcolorbox}%
}

% ============================================================
% Page de garde
% #1 titre · #2 matière · #3 niveau · #4 module · #5 n° leçon · #6 durée ·
% #7 description / objectifs (texte ou itemize)
% ============================================================
\newcommand{\pagegarde}[7]{%
  \thispagestyle{empty}
  \newgeometry{a4paper,margin=1.7cm,top=1.6cm,bottom=1.4cm}%
  \begin{tikzpicture}[remember picture,overlay]
    \fill[poporange!18] ([xshift=-3.2cm,yshift=-3.6cm]current page.north east) circle (4.6cm);
    \fill[poppurple!14] ([xshift=2.4cm,yshift=7.6cm]current page.south west) circle (3.8cm);
  \end{tikzpicture}%
  {\poplogo\fontsize{30}{30}\selectfont\color{popink}OnBuch\color{poporange}+}\hfill
  \raisebox{0.6ex}{\poppill{Cours signé \textbullet\ Premium}{popink}{popcream}}\par
  \vspace{1pt}\popsquiggle{poppurple}\par
  \vspace{8pt}
  \begin{tcolorbox}[enhanced,colback=poporange,colframe=popink,boxrule=2.8pt,arc=13pt,
      drop shadow=popink,left=18pt,right=18pt,top=12pt,bottom=13pt,after skip=10pt]
    \poppill{#2 \textbullet\ #3}{popink}{popcream}\hspace{5pt}\poppill{#4}{white}{popink}\par\vspace{8pt}
    {\popdisplay\fontsize{14}{14}\selectfont\color{popink}LEÇON #5}\par\vspace{4pt}
    {\raggedright\hyphenpenalty=10000\exhyphenpenalty=10000\popdisplay\fontsize{25}{27}\selectfont\color{popcream}\MakeUppercase{#1}\par}
    \vspace{8pt}
    {\popxbold\fontsize{9.5}{9.5}\selectfont\color{popink}\MakeUppercase{Durée conseillée : #6}}
  \end{tcolorbox}
  \begin{tcolorbox}[enhanced,colback=white,colframe=popink,boxrule=2.2pt,arc=10pt,
      drop shadow=popink,left=16pt,right=16pt,top=10pt,bottom=10pt,after skip=8pt]
    \poppill{Dans cette leçon}{poppurple}{white}\par\vspace{5pt}
    \popsemi\fontsize{9.3}{12.6}\selectfont\color{popink}#7
  \end{tcolorbox}
  \noindent\begin{minipage}[t]{0.487\linewidth}
    \begin{tcolorbox}[enhanced,colback=popgreen,colframe=popink,boxrule=2.2pt,arc=10pt,
        drop shadow=popink,left=11pt,right=11pt,top=10pt,bottom=10pt,height=3.1cm,valign=center]
      \tcbox[colback=white,colframe=popink,boxrule=1.3pt,arc=5pt,boxsep=0pt,nobeforeafter,
        left=3pt,right=3pt,top=3pt,bottom=3pt,valign=center]{\qrcode[height=1.9cm]{https://play.google.com/store/apps/details?id=com.luvvix.onbuch&hl=fr}}%
      \hspace{8pt}\begin{minipage}[c]{0.5\linewidth}
        {\popdisplay\fontsize{11}{12.5}\selectfont\color{white}\MakeUppercase{Télécharge OnBuch}}\par\vspace{4pt}
        {\raggedright\popsemi\fontsize{8.4}{11}\selectfont\color{white}Gratuit \textbullet\ Google Play\\Tuteur IA, annales, concours, résultats\par}
      \end{minipage}
    \end{tcolorbox}
  \end{minipage}\hfill
  \begin{minipage}[t]{0.487\linewidth}
    \begin{tcolorbox}[enhanced,colback=poppurple,colframe=popink,boxrule=2.2pt,arc=10pt,
        drop shadow=popink,left=11pt,right=11pt,top=10pt,bottom=10pt,height=3.1cm,valign=center]
      \tikz[baseline=-0.7ex]\node[circle,fill=white,draw=popink,line width=1.3pt,minimum size=1.6cm,inner sep=0]{\popdisplay\fontsize{24}{24}\selectfont\color{poppurple}+};%
      \hspace{8pt}\begin{minipage}[c]{0.6\linewidth}
        {\popdisplay\fontsize{11}{12.5}\selectfont\color{white}\MakeUppercase{Passe à OnBuch+}}\par\vspace{4pt}
        {\raggedright\popsemi\fontsize{8.4}{11}\selectfont\color{white}Tous les cours signés, Tuteur IA illimité, fiches PDF — onglet OnBuch+ de l'app\par}
      \end{minipage}
    \end{tcolorbox}
  \end{minipage}
  \vspace{8pt}
  \popcontenu
  \vfill
  \signatureonbuch
  \restoregeometry
  \newpage
}

% Rangée « Ce cours contient » (page de garde)
\newcommand{\poptile}[3]{% #1 couleur #2 gros texte #3 légende
  \begin{tcolorbox}[enhanced,colback=#1,colframe=popink,boxrule=2pt,arc=9pt,shadow={2.5pt}{-2.5pt}{0pt}{popink},
      left=6pt,right=6pt,top=9pt,bottom=9pt,halign=center,valign=center,height=1.75cm,width=\linewidth,nobeforeafter]
    {\popdisplay\fontsize{12.5}{13}\selectfont\color{white}\MakeUppercase{#2}}\par\vspace{3pt}
    {\centering\popsemi\fontsize{7.8}{9.5}\selectfont\color{white}#3\par}
  \end{tcolorbox}}
\newcommand{\popcontenu}{%
  \par\noindent{\popxboldls\fontsize{7.5}{7.5}\selectfont\color{popmuted}CE COURS SIGNÉ CONTIENT}\par\vspace{7pt}
  \noindent\begin{minipage}[t]{0.235\linewidth}\poptile{popblue}{Cours}{complet, illustré, pas à pas}\end{minipage}\hfill
  \begin{minipage}[t]{0.235\linewidth}\poptile{poporange}{Méthodes}{et exercices résolus}\end{minipage}\hfill
  \begin{minipage}[t]{0.235\linewidth}\poptile{poppink}{Exercices}{type examen + corrigés}\end{minipage}\hfill
  \begin{minipage}[t]{0.235\linewidth}\poptile{popgreen}{Fiche bilan}{l'essentiel à retenir}\end{minipage}\par}

% Sommaire
\newtcolorbox{sommaire}{enhanced,colback=white,colframe=popink,boxrule=2.2pt,arc=10pt,
  drop shadow=popink,left=18pt,right=18pt,top=13pt,bottom=13pt,before skip=6pt,after skip=20pt,
  before upper={\poppill{Sommaire}{poppurple}{white}\par\vspace{9pt}\popsemi\fontsize{10.6}{14.5}\selectfont\color{popink}}}

% Signature de fin de cours — poussée tout en bas de la dernière page
\newcommand{\findecours}{%
  \par\vfill\nopagebreak
  \begin{center}\popsquiggle{poporange}\end{center}\vspace{4pt}
  \signatureonbuch
}
\AtBeginDocument{\def\llbracket{\mathopen{[\mkern-3.5mu[}}\def\rrbracket{\mathclose{]\mkern-3.5mu]}}} % intervalles d'entiers (glyphe ⟦ absent de la police)
% Glyphes absents de Fira Math : redéfinis à partir de glyphes disponibles
\AtBeginDocument{%
  \def\setminus{\mathbin{\backslash}}\let\smallsetminus\setminus
  \def\vdots{\mathord{\vbox{\baselineskip3.2pt\lineskiplimit0pt\kern4pt\hbox{.}\hbox{.}\hbox{.}}}}%
  \def\ddots{\mathinner{\mkern1mu\raise6.5pt\hbox{.}\mkern2mu\raise3.6pt\hbox{.}\mkern2mu\raise0.7pt\hbox{.}\mkern1mu}}%
  \def\triangle{\mathord{\scalebox{0.9}{$\symup{\Delta}$}}}%
  \def\nabla{\mathord{\raisebox{\depth}{\scalebox{1}[-1]{$\symup{\Delta}$}}}}%
  \def\checkmark{\popcheck{popdark}}%
  \def\star{\mathbin{\ast}}\def\bigstar{\mathord{\ast}}%
  \def\diamond{\mathbin{\scalebox{0.6}{$\Diamond$}}}%
  \def\bigcup{\mathop{\mathchoice{\raisebox{-0.3em}{\scalebox{1.7}{$\cup$}}}{\scalebox{1.25}{$\cup$}}{\cup}{\cup}}\displaylimits}%
  \def\bigcap{\mathop{\mathchoice{\raisebox{-0.3em}{\scalebox{1.7}{$\cap$}}}{\scalebox{1.25}{$\cap$}}{\cap}{\cap}}\displaylimits}%
  \def\lesseqqgtr{\mathrel{\lessgtr}}\def\bigoplus{\mathop{\oplus}\displaylimits}%
}
\providecommand{\ssi}{si et seulement si} % abréviation fréquente
\AtBeginDocument{\providecommand{\wideparen}{\overparen}} % arc de cercle

%% FIGFIX-BEGIN v5 — correctifs globaux des figures (docs/onbuch-generation/tools/figfix)
\usepackage{adjustbox}\usepackage{etoolbox}\tcbuselibrary{raster}
% toute figure TikZ/pgfplots plus large que la ligne est réduite à la largeur disponible
\BeforeBeginEnvironment{tikzpicture}{\begin{adjustbox}{max width=\linewidth}}
\AfterEndEnvironment{tikzpicture}{\end{adjustbox}}
% légendes pgfplots lisibles par-dessus les courbes
\pgfplotsset{every axis/.append style={legend style={fill=white,fill opacity=0.92,text opacity=1,draw=black!15,inner sep=3pt}}}
% étiquettes d'arêtes (syntaxe edge["..."]) avec fond blanc
\tikzset{every edge quotes/.append style={fill=white,inner sep=1.2pt}}
%% FIGFIX-END
\begin{document}

\pagegarde{Lesson 2 — Grammar: Modals: can/may (permission), can/could (possibility), ought to/ought not to, need(n't), don't have to, should/ought to, must + deduction}{Anglais}{1ère A}{Chapter 1 : Family and Social Life}{ENA02}{2 h}{Cette leçon de grammaire présente les principaux verbes modaux anglais : can/may pour la permission, can/could pour la possibilité, ought to, should, must pour le conseil, l'obligation et la déduction, ainsi que needn't et don't have to pour l'absence d'obligation. Observation, règles, comparaison avec le français et entraînement systématique préparent l'élève à les employer avec précision dans le thème « Family and Social Life ».
\vspace{4pt}
\begin{multicols}{2}\small
\begin{itemize}
\item À la fin de cette leçon, je sais identifier un verbe modal et reconnaître sa valeur (permission, possibilité, conseil, obligation, déduction, absence d'obligation)
\item À la fin de cette leçon, je sais conjuguer correctement un modal : sans -s à la 3e personne, suivi de la base verbale, négation et interrogation sans auxiliaire do
\item À la fin de cette leçon, je sais demander et accorder la permission avec can, could et may selon le degré de politesse
\item À la fin de cette leçon, je sais exprimer la possibilité présente et passée avec can et could
\item À la fin de cette leçon, je sais donner un conseil avec should et ought to, et un déconseil avec shouldn't et ought not to
\item À la fin de cette leçon, je sais distinguer mustn't (interdiction) de don't have to / needn't (absence d'obligation)
\end{itemize}\end{multicols}}

\begin{sommaire}
\begin{multicols}{2}
\begin{enumerate}
\item Observation : les modaux en action
\item Permission et possibilité : can, could, may
\item Conseil : should et ought to
\item Obligation, interdiction et absence d'obligation : must, mustn't, needn't, don't have to
\item Déduction : must et can't logiques
\item Comparaison français/anglais, synthèse et entraînement
\item Activité d'intégration
\item Méthodes et exercices résolus
\item Exercices
\item Corrigés des exercices
\item Fiche bilan
\end{enumerate}
\end{multicols}
\end{sommaire}

\begin{objectifs}
\begin{itemize}
\item À la fin de cette leçon, je sais identifier un verbe modal et reconnaître sa valeur (permission, possibilité, conseil, obligation, déduction, absence d'obligation)
\item À la fin de cette leçon, je sais conjuguer correctement un modal : sans -s à la 3e personne, suivi de la base verbale, négation et interrogation sans auxiliaire do
\item À la fin de cette leçon, je sais demander et accorder la permission avec can, could et may selon le degré de politesse
\item À la fin de cette leçon, je sais exprimer la possibilité présente et passée avec can et could
\item À la fin de cette leçon, je sais donner un conseil avec should et ought to, et un déconseil avec shouldn't et ought not to
\item À la fin de cette leçon, je sais distinguer mustn't (interdiction) de don't have to / needn't (absence d'obligation)
\item À la fin de cette leçon, je sais exprimer une déduction logique avec must (certitude) et can't (impossibilité)
\item À la fin de cette leçon, je sais éviter les erreurs typiques des francophones (to can, must to, ought do, need to not...)
\item À la fin de cette leçon, je sais utiliser les modaux dans des phrases et un court paragraphe sur la vie familiale et sociale
\end{itemize}
\end{objectifs}

\begin{prerequis}
\begin{itemize}
\item Le présent simple et la base verbale (infinitif sans to)
\item La formation de la négation et de l'interrogation en anglais (auxiliaire do/does)
\item Les verbes can et must vus au collège (sens de base : capacité, obligation)
\item Le prétérit simple (utile pour could comme passé de can)
\item Le vocabulaire de la famille et de la vie sociale (Lesson 1 du chapitre)
\end{itemize}
\end{prerequis}

\newpage

\coursec{Observation : les modaux en action}

\courssub{Situation d'accroche : une réponse, quatre modaux}

Pendant les grandes vacances, Amina, élève de Première à Yaoundé, veut rendre visite à sa tante à Buea. Sa mère lui répond : \eng{You may go, but you must be back before Sunday. You ought to call your aunt first, and you needn't worry about transport — your uncle can drive you.} « Tu peux y aller, mais tu dois être de retour avant dimanche. Tu devrais d'abord appeler ta tante, et tu n'as pas à t'inquiéter pour le transport — ton oncle peut te conduire. »

En une seule réponse, quatre verbes modaux expriment la \cle{permission} (\eng{may}), l'\cle{obligation} (\eng{must}), le \cle{conseil} (\eng{ought to}) et l'\cle{absence de nécessité} (\eng{needn't}), sans oublier la \cle{capacité} (\eng{can}). Le français utilise des verbes ordinaires conjugués (\emph{pouvoir, devoir, falloir}) ; l'anglais, lui, emploie de petits auxiliaires spéciaux, les \cle{modaux}, qui obéissent à des règles très strictes. Comment les reconnaître, comment les construire, et comment choisir le bon modal ? C'est toute l'ambition de cette leçon. Commençons par observer.

\courssub{Un dialogue à observer}

\begin{textebox}[Dialogue : At home before the village feast]
\textbf{Manka:} Mum, can I go to the village feast on Saturday?

\textbf{Mother:} You may go, but you ought to help your grandmother first. She cannot carry the heavy pots alone.

\textbf{Manka:} Of course, Mum. Must I wear my uniform?

\textbf{Mother:} No, you needn't wear your uniform. But you should put on clean clothes — the chief will be there.

\textbf{Manka:} Could Uncle Tabi drive me there? The bendskins are always full on Saturdays.

\textbf{Mother:} Yes, he could. But you must be back before dark, and you mustn't stay near the road at night.

\textbf{Manka:} Do I have to take my little brother with me?

\textbf{Mother:} No, you don't have to. He can stay with me. Oh, and that noise in the kitchen — it must be your father. He is always hungry before a feast!

\textbf{Manka:} Thank you, Mum. I ought to finish my homework first, oughtn't I?

\textbf{Mother:} Exactly. You should never leave your books behind, even during a feast.
\end{textebox}

\textbf{Glossaire.} \emph{feast} (n.) : fête, festin ; \emph{pot} (n.) : marmite ; \emph{bendskin} (n.) : moto-taxi (Cameroun) ; \emph{chief} (n.) : chef (du village) ; \emph{dark} (n.) : la nuit tombée ; \emph{noise} (n.) : bruit ; \emph{hungry} (adj.) : affamé.

\textbf{Questions de compréhension.}
\begin{enumerate}
\item Where does Manka want to go on Saturday?
\item What must she do before going?
\item Why does she mention Uncle Tabi?
\item Why does the mother say \eng{it must be your father}?
\end{enumerate}

\courssub{Repérage : trouvons les modaux dans le dialogue}

Relisons le dialogue et soulignons les verbes modaux. Nous trouvons :

\begin{center}
\begin{tabularx}{\linewidth}{|l|X|l|}
\hline
\rowcolor{popblueL} \textbf{Modal} & \textbf{Exemple tiré du dialogue} & \textbf{Fonction} \\
\hline
\eng{can} & \eng{Can I go to the village feast?} & permission \\
\hline
\eng{may} & \eng{You may go.} & permission \\
\hline
\eng{can} & \eng{She cannot carry the heavy pots.} & capacité (négation) \\
\hline
\eng{could} & \eng{Could Uncle Tabi drive me there?} & possibilité \\
\hline
\eng{could} & \eng{Yes, he could.} & capacité / possibilité \\
\hline
\eng{should} & \eng{You should put on clean clothes.} & conseil \\
\hline
\eng{ought to} & \eng{You ought to help your grandmother.} & conseil \\
\hline
\eng{must} & \eng{You must be back before dark.} & obligation \\
\hline
\eng{mustn't} & \eng{You mustn't stay near the road.} & interdiction \\
\hline
\eng{needn't} & \eng{You needn't wear your uniform.} & absence d'obligation \\
\hline
\eng{don't have to} & \eng{You don't have to take your brother.} & absence d'obligation \\
\hline
\eng{can} & \eng{He can stay with me.} & capacité \\
\hline
\eng{must} (déduction) & \eng{It must be your father.} & déduction logique \\
\hline
\eng{ought to} & \eng{I ought to finish my homework.} & conseil / auto-obligation morale \\
\hline
\end{tabularx}
\end{center}

\begin{exemplebox}[Trois exemples traduits]
\eng{He can speak Ewondo.} « Il sait parler ewondo. » (capacité)

\eng{You must listen to your parents.} « Tu dois écouter tes parents. » (obligation)

\eng{We needn't hurry.} « Nous n'avons pas besoin de nous dépêcher. » (absence de nécessité)
\end{exemplebox}

\courssub{Constats guidés : quatre observations grammaticales}

\textbf{Constat 1 — le modal est suivi de la base verbale.} Observons : \eng{can go}, \eng{may go}, \eng{must be}, \eng{should put on}, \eng{could drive}, \eng{cannot carry}. Après chaque modal, on trouve l'\cle{infinitif sans to}, appelé \emph{base verbale}. On dit \eng{you must be back}, jamais \eng{you must to be back}. Seule exception apparente : \eng{ought to}, où le \eng{to} fait partie du modal lui-même (\eng{ought to help}).

\textbf{Constat 2 — pas de -s à la troisième personne.} Observons : \eng{She cannot carry the heavy pots}, \eng{He can stay with me}. Le sujet est à la troisième personne du singulier (\eng{she}, \eng{he}), pourtant le modal ne prend jamais de \emph{-s}. On dit \eng{he can}, jamais \eng{he cans}.

\textbf{Constat 3 — la négation se fait par not, sans do.} Observons : \eng{cannot} (souvent écrit en un seul mot), \eng{mustn't}, \eng{needn't}. Il suffit d'ajouter \eng{not} après le modal ; l'auxiliaire \eng{do/does} est inutile et interdit : jamais \eng{she doesn't can}.

\textbf{Constat 4 — l'interrogation se fait par inversion.} Observons : \eng{Can I go?}, \eng{Could Uncle Tabi drive me?}, \eng{Must I wear my uniform?} Le modal se place devant le sujet, sans \eng{do/does} : jamais \eng{Do you can come?}.

\begin{propriete}[Le verbe modal : un auxiliaire défectif]
Un \cle{verbe modal} est un auxiliaire \emph{défectif} : il n'a ni infinitif, ni participe, ni conjugaison complète. Trois règles absolues :
\begin{enumerate}
\item il ne prend \textbf{jamais de -s} à la troisième personne du singulier ;
\item il est suivi de la \textbf{base verbale} (infinitif \textbf{sans to}) ;
\item il forme sa \textbf{négation} avec \eng{not} et son \textbf{interrogation} par \textbf{inversion}, sans \eng{do/does}.
\end{enumerate}
Exemple : \eng{She can swim.} / \eng{She cannot swim.} / \eng{Can she swim?}
\end{propriete}

\begin{popfigure}
\begin{tikzpicture}[mindmap, grow cyclic, every node/.style={concept}, root concept/.style={concept color=popblue, text=white, font=\bfseries}, level 1/.append style={level distance=3.4cm, sibling angle=60}, level 2/.append style={level distance=2.2cm}]
\node[root concept]{MODALS}
  child[concept color=popgreen] { node {permission} child { node[concept color=popgreenL, text=popdark, align=center]{can, may} } }
  child[concept color=poporange] { node {possibility} child { node[concept color=poporangeL, text=popdark, align=center]{can, could} } }
  child[concept color=poppurple] { node {advice} child { node[concept color=poppurpleL, text=popdark, align=center]{should,\\ought to} } }
  child[concept color=poppink] { node {obligation} child { node[concept color=poppinkL, text=popdark, align=center]{must,\\mustn't} } }
  child[concept color=popgold] { node {deduction} child { node[concept color=popgoldL, text=popdark, align=center]{must,\\can't} } }
  child[concept color=popmuted] { node[align=center]{no\\obligation} child { node[concept color=popcream, text=popdark, align=center]{needn't,\\don't have to} } };
\end{tikzpicture}
\legende{Carte mentale : les modaux classés par fonction.}
\end{popfigure}

\begin{popfigure}
\begin{tikzpicture}[node distance=0.5cm]
\node[draw=popblue, fill=popblueL, rounded corners, minimum width=2.6cm, minimum height=1cm, font=\bfseries] (s) {Subject};
\node[draw=poppink, fill=poppinkL, rounded corners, minimum width=2.6cm, minimum height=1cm, font=\bfseries, right=of s] (m) {Modal};
\node[draw=popgreen, fill=popgreenL, rounded corners, minimum width=3.2cm, minimum height=1cm, font=\bfseries, right=of m] (b) {Base verbale};
\node[draw=popmuted, fill=popcream, rounded corners, minimum width=3cm, minimum height=1cm, font=\bfseries, right=of b] (c) {(complement)};
\draw[-{Stealth}, thick, popdark] (s) -- (m);
\draw[-{Stealth}, thick, popdark] (m) -- (b);
\draw[-{Stealth}, thick, popdark] (b) -- (c);
\node[below=0.35cm of s, font=\itshape] {You};
\node[below=0.35cm of m, font=\itshape] {must};
\node[below=0.35cm of b, font=\itshape] {be};
\node[below=0.35cm of c, font=\itshape] {back before dark.};
\end{tikzpicture}
\legende{Schéma de structure : Sujet + modal + base verbale (+ complément).}
\end{popfigure}

\courssub{Comparaison avec le français}

\begin{center}
\begin{tabularx}{\linewidth}{|X|X|X|}
\hline
\rowcolor{popblueL} \textbf{Français} & \textbf{English} & \textbf{Remarque} \\
\hline
Il sait nager. & \eng{He can swim.} & Pas de \emph{-s} sur \eng{can}. \\
\hline
Elle doit partir. & \eng{She must go.} & Pas de \eng{to} après \eng{must}. \\
\hline
Peux-tu venir ? & \eng{Can you come?} & Inversion, sans \eng{do}. \\
\hline
Tu ne dois pas fumer. & \eng{You mustn't smoke.} & Interdiction (et non « ne pas être obligé »). \\
\hline
Tu n'es pas obligé de venir. & \eng{You needn't come.} & Absence d'obligation. \\
\hline
\end{tabularx}
\end{center}

\begin{attention}[Erreurs fréquentes des francophones]
Ne dites jamais \eng{He cans swim}, \eng{She must to go}, \eng{Do you can come?}. Le modal \textbf{remplace} l'auxiliaire \eng{do} et \textbf{interdit} le \eng{to} devant la base verbale. Le français « il peut nager » fait croire qu'il faut conjuguer \eng{can} comme un verbe ordinaire : c'est faux, \eng{can} est invariable.
\end{attention}

\begin{methode}[Comment identifier un modal]
Pour vérifier qu'un mot est bien un modal, contrôlez \textbf{trois indices} :
\begin{enumerate}
\item \textbf{Pas de -s} à la troisième personne : \eng{he must}, jamais \eng{he musts}.
\item \textbf{Base verbale sans to} juste après : \eng{must go}, \eng{can swim}.
\item \textbf{Inversion} à l'interrogation : \eng{Must he go?}, sans \eng{does}.
\end{enumerate}
Si les trois indices sont présents, c'est un modal. Attention : \eng{have to} échoue au troisième test (\eng{Does he have to go?}) : ce n'est pas un vrai modal, mais un \emph{semi-modal}.
\end{methode}

\begin{exoresolu}[Classement par fonction]
Énoncé : classez ces phrases tirées du dialogue selon leur fonction : permission, possibilité, conseil, obligation, déduction, absence d'obligation, capacité.
(a) \eng{You may go.} (b) \eng{It must be your father.} (c) \eng{You needn't wear your uniform.} (d) \eng{You should put on clean clothes.} (e) \eng{Uncle Tabi could drive you.} (f) \eng{You must be back before dark.} (g) \eng{She cannot carry the heavy pots.}
\tcblower
Solution : (a) \eng{may} exprime la \textbf{permission} accordée par la mère. (b) \eng{must} exprime ici une \textbf{déduction} logique : la mère conclut que c'est le père, à cause du bruit. (c) \eng{needn't} exprime l'\textbf{absence d'obligation} : ce n'est pas nécessaire. (d) \eng{should} exprime un \textbf{conseil}. (e) \eng{could} exprime une \textbf{possibilité} (ou une capacité future). (f) \eng{must} exprime une \textbf{obligation} ferme. (g) \eng{cannot} exprime l'\textbf{incapacité} (négation de \eng{can} capacité). Remarquez que le même modal (\eng{must}) peut changer de fonction selon le contexte : obligation en (f), déduction en (b) ; de même \eng{can} exprime la permission en (a) et la capacité en (g).
\end{exoresolu}

\begin{exercice}[À vous de jouer]
Soulignez le modal dans chaque phrase et indiquez sa fonction (permission, conseil, obligation, absence d'obligation, possibilité, déduction, capacité) :
\begin{enumerate}
\item \eng{Can I borrow your bicycle, Pa?}
\item \eng{You ought to visit your grandmother more often.}
\item \eng{We needn't buy rice; we have enough at home.}
\item \eng{That girl looks like Amina — it must be her sister.}
\item \eng{Students must respect their teachers.}
\item \eng{My little brother can already read.}
\end{enumerate}
\end{exercice}

\begin{corrige}[Exercice : À vous de jouer]
1. \eng{Can} — permission (demande polie). 2. \eng{ought to} — conseil. 3. \eng{needn't} — absence d'obligation. 4. \eng{must} — déduction logique. 5. \eng{must} — obligation. 6. \eng{can} — capacité.
\end{corrige}

\begin{aretenir}[L'essentiel de l'observation]
Les modaux anglais (\eng{can, could, may, must, should, ought to, needn't}) sont des auxiliaires défectifs : invariables, suivis de la base verbale sans \eng{to}, avec négation par \eng{not} et interrogation par inversion. Chaque modal exprime une fonction : permission, possibilité, capacité, conseil, obligation, interdiction, absence d'obligation ou déduction. La suite de la leçon étudiera chaque fonction en détail : permission et possibilité (\eng{can, could, may}), conseil (\eng{should, ought to}), obligation et absence d'obligation (\eng{must, mustn't, needn't, don't have to}), déduction (\eng{must, can't}).
\end{aretenir}

\coursec{Permission et possibilité : can, could, may}

Dans la section précédente, nous avons observé les modaux \eng{can}, \eng{could} et \eng{may} en action dans un dialogue authentique. Nous avons remarqué qu'ils expriment des idées très différentes : demander la permission, évoquer une possibilité, parler d'une capacité passée. Il est temps maintenant de systématiser ces emplois. Cette section étudie en profondeur deux grandes fonctions de ces trois auxiliaires : la \cle{permission} (demander, donner, refuser le droit de faire quelque chose) et la \cle{possibilité} (dire qu'un événement est possible).

\courssub{Observer avant de comprendre : un dialogue en classe}

Lisons d'abord ce court dialogue qui se déroule dans une classe de première à Bamenda. Repérez les modaux en gras et demandez-vous : qui parle ? à qui ? dans quel but ?

\begin{textebox}[In the classroom — Bamenda]
\textbf{Student:} Excuse me, Sir. \textbf{May} I come in? I am late because the bendskin broke down.

\textbf{Teacher:} Yes, you \textbf{may}. But next time, leave home earlier.

\textbf{Student:} Thank you, Sir. \textbf{Could} I also borrow a pen from my neighbour? Mine is finished.

\textbf{Teacher:} Of course. Ngong, \textbf{can} you lend him a pen?

\textbf{Ngong:} Sure, Sir. Here you are.

\textbf{Student:} Sir, one more thing. \textbf{Can} I go to the toilet, please?

\textbf{Teacher:} You \textbf{can} go after the exercise. By the way, it \textbf{can} rain heavily this afternoon, so nobody \textbf{may} leave the school before the bell rings.

\textbf{Student:} When I was in primary school, we \textbf{could} go home at any time!

\textbf{Teacher:} Times have changed, Ekane. Sit down and open your book.
\end{textebox}

\textbf{Glossaire :} \emph{bendskin} (n., mot camerounais : moto-taxi), \emph{to break down} (v. : tomber en panne), \emph{to borrow} (v. : emprunter), \emph{to lend} (v. : prêter), \emph{the bell} (n. : la sonnerie), \emph{neighbour} (n. : voisin / camarade de banc).

\textbf{Questions de compréhension :}
\begin{enumerate}
\item Why is Ekane late?
\item Which modal does he use to address the teacher? Why this one and not \eng{can}?
\item Find two sentences where a modal expresses a \emph{possibility}, not a permission.
\item In the last sentence of the student, what does \eng{could} refer to: present, past or future?
\end{enumerate}

Observons : le même mot \eng{can} sert à demander la permission (\eng{Can I go to the toilet?}) et à exprimer une possibilité (\eng{It can rain heavily}). Le mot \eng{could} sert à formuler une demande polie (\eng{Could I borrow a pen?}) et à parler du passé (\eng{we could go home at any time}). Le mot \eng{may} apparaît dans les échanges avec une autorité (\eng{May I come in?} / \eng{nobody may leave}). Ce n'est pas un hasard : chaque modal a son \emph{registre}, c'est-à-dire son niveau de langue.

\courssub{La permission : can, could, may}

\begin{propriete}[Permission — can (courant), could (poli), may (formel)]
Pour \textbf{demander la permission}, on utilise un modal + sujet + base verbale :
\begin{itemize}
\item \eng{Can I...?} : registre \textbf{courant}, entre amis, camarades, membres de la famille. \eng{Can I use your pen?} « Je peux prendre ton stylo ? »
\item \eng{Could I...?} : registre \textbf{poli, atténué}, pour ménager l'interlocuteur ou s'adresser à quelqu'un qu'on connaît peu. \eng{Could I leave early today?} « Est-ce que je pourrais partir plus tôt aujourd'hui ? »
\item \eng{May I...?} : registre \textbf{formel, très poli}, devant une autorité (professeur, proviseur, chef) ou dans un contexte officiel. \eng{May I speak to the headmaster?} « Puis-je parler au proviseur ? »
\end{itemize}
Pour \textbf{donner ou refuser la permission}, on répond avec \eng{can} ou \eng{may} (jamais avec \eng{could}, qui ne sert qu'à \emph{demander}) :
\begin{itemize}
\item \eng{Yes, you can.} / \eng{Yes, you may.} « Oui, tu peux. / Oui, vous pouvez. »
\item \eng{No, you can't.} (refus courant) / \eng{No, you may not.} (refus \textbf{formel}, souvent écrit sur les panneaux et règlements).
\end{itemize}
\end{propriete}

\begin{attention}[Could demande, can répond]
Erreur classique du francophone : répondre \eng{Yes, you could} à une demande polie. C'est impossible en anglais ! \eng{Could} est une forme d'\emph{atténuation} de la demande ; la réponse revient au présent :
\begin{center}
\eng{Could I use your phone? — Yes, you can.} (et jamais \emph{Yes, you could.})
\end{center}
De même, \eng{May I...?} appelle la réponse \eng{Yes, you may} ou \eng{No, you may not}.
\end{attention}

\begin{exemplebox}[Exemples traduits]
\begin{itemize}
\item \eng{Could I go out, please?} « Est-ce que je pourrais sortir, s'il vous plaît ? »
\item \eng{Can I sit here?} « Je peux m'asseoir ici ? »
\item \eng{May I come in, Sir?} « Puis-je entrer, Monsieur ? »
\item \eng{You may not enter the staffroom.} « Vous n'êtes pas autorisé à entrer dans la salle des professeurs. »
\item \eng{You can't use your phone during the exam.} « Tu ne peux pas utiliser ton téléphone pendant l'examen. »
\item \eng{Could we possibly start the meeting now?} « Pourrions-nous éventuellement commencer la réunion maintenant ? »
\end{itemize}
\end{exemplebox}

\begin{savaistu}[« May I come in? » : la formule rituelle]
Dans l'anglais britannique soigné — et dans les écoles des régions anglophones du Cameroun (Buea, Bamenda, Kumba) —, \eng{may} reste la marque de politesse attendue pour s'adresser à un supérieur. Un élève qui arrive en retard ne dira jamais \eng{Can I come in?} à son professeur : la formule rituelle est \eng{May I come in, Sir?} ou \eng{May I come in, Madam?}. Dans la conversation quotidienne en Grande-Bretagne et aux États-Unis, \eng{can} a largement remplacé \eng{may}, mais à l'écrit formel et à l'école, \eng{may} garde tout son prestige. À l'examen, montrez que vous connaissez les deux registres !
\end{savaistu}

\courssub{La possibilité : can au présent, could au passé ou pour l'hypothèse}

\begin{propriete}[Possibilité — can (présent) et could (passé ou hypothèse)]
\begin{itemize}
\item \eng{Can} exprime une \textbf{possibilité générale au présent} : quelque chose qui arrive parfois, qui est possible dans certaines conditions. \eng{It can be dangerous.} « Cela peut être dangereux. » \eng{It can rain heavily in Buea.} « Il peut pleuvoir très fort à Buea. »
\item \eng{Could} exprime une \textbf{possibilité au passé} (passé de \eng{can}) : \eng{When I was young, we could play outside till night.} « Quand j'étais jeune, nous pouvions jouer dehors jusqu'à la nuit. »
\item \eng{Could} exprime aussi une \textbf{possibilité hypothétique au présent ou au futur}, moins certaine que \eng{can} : \eng{It could rain tomorrow.} « Il pourrait pleuvoir demain. »
\end{itemize}
Attention à la nuance : \eng{It can rain here in July} (cela arrive réellement, c'est un fait général) ; \eng{It could rain this evening} (c'est une éventualité, une hypothèse pour ce soir précis).
\end{propriete}

\begin{exemplebox}[Exemples traduits]
\begin{itemize}
\item \eng{In the rainy season, roads can be blocked.} « En saison des pluies, les routes peuvent être bloquées. »
\item \eng{Harmattan winds can be very strong in the North.} « Les vents d'harmattan peuvent être très forts dans le Nord. »
\item \eng{Last year, we could visit grandma every week.} « L'année dernière, nous pouvions rendre visite à grand-mère chaque semaine. »
\item \eng{Be careful with that knife — it could be dangerous.} « Fais attention avec ce couteau — il pourrait être dangereux. »
\item \eng{Don't worry, the bus could still arrive on time.} « Ne t'inquiète pas, le bus pourrait encore arriver à l'heure. »
\end{itemize}
\end{exemplebox}

\textbf{Deux emplois de could : ne pas confondre.} Le mot \eng{could} a deux vies : c'est à la fois le \emph{passé} de \eng{can} et une forme \emph{conditionnelle d'atténuation}. Comment les distinguer ? Regardez le contexte et le temps des autres verbes :

\begin{center}
\begin{tabularx}{\linewidth}{|X|X|X|}
\hline
\rowcolor{popblueL} Phrase & Valeur de could & Indice \\
\hline
\eng{When I was ten, I could swim well.} & passé (« je savais nager ») & \eng{when I was ten} \\
\hline
\eng{Could you open the window?} & demande polie au présent & question à quelqu'un de présent \\
\hline
\eng{It could rain tomorrow.} & hypothèse future (« il pourrait ») & \eng{tomorrow} \\
\hline
\end{tabularx}
\end{center}

\courssub{Tableau récapitulatif des formes}

Rappel fondamental : après \eng{can}, \eng{could} et \eng{may}, le verbe est toujours à la \cle{base verbale} (infinitif sans \eng{to}), sans \eng{-s} à la troisième personne, sans \eng{do} à la question et à la négation.

\begin{center}
\begin{tabular}{|l|l|l|l|}
\hline
\rowcolor{popblueL} Modal & Affirmatif & Négatif & Interrogatif \\
\hline
\eng{can} & \eng{I can go.} & \eng{I can't (cannot) go.} & \eng{Can I go?} \\
\hline
\eng{could} & \eng{I could go.} & \eng{I couldn't go.} & \eng{Could I go?} \\
\hline
\eng{may} & \eng{I may go.} & \eng{I may not go.} & \eng{May I go?} \\
\hline
\end{tabular}
\end{center}

\begin{attention}[Can n'a ni futur ni infinitif]
Le français conjugue « pouvoir » à tous les temps ; l'anglais, non ! \eng{Can} n'existe qu'au présent, et \eng{could} au passé. Pour tous les autres temps, on utilise le substitut \cle{be able to} :
\begin{itemize}
\item Futur : \eng{I will be able to help you tomorrow.} « Je pourrai t'aider demain. » (jamais \emph{I will can...})
\item Infinitif : \eng{I hope to be able to come.} « J'espère pouvoir venir. » (jamais \emph{to can})
\item Present perfect : \eng{She has been able to finish.} « Elle a pu finir. »
\end{itemize}
Autre piège : ne jamais dire \emph{I can to go} — après un modal, pas de \eng{to} : \eng{I can go.}
\end{attention}

\courssub{Figures de synthèse}

\begin{popfigure}
\begin{center}
\begin{tikzpicture}
\draw[->, thick, popline] (0,0) -- (10.5,0);
\node[draw, rounded corners, fill=popgreenL, align=center, minimum width=2.6cm, minimum height=1.3cm] at (1.8,0.9) {\textbf{CAN}\\ courant\\ \footnotesize entre amis};
\node[draw, rounded corners, fill=poporangeL, align=center, minimum width=2.6cm, minimum height=1.3cm] at (5.2,0.9) {\textbf{COULD}\\ poli, atténué\\ \footnotesize pour ménager};
\node[draw, rounded corners, fill=poppurpleL, align=center, minimum width=2.6cm, minimum height=1.3cm] at (8.6,0.9) {\textbf{MAY}\\ formel\\ \footnotesize autorité, officiel};
\node[align=center, font=\footnotesize, popmuted] at (1.8,-0.6) {Can I use your pen?};
\node[align=center, font=\footnotesize, popmuted] at (5.2,-0.6) {Could I leave early?};
\node[align=center, font=\footnotesize, popmuted] at (8.6,-0.6) {May I come in, Sir?};
\end{tikzpicture}
\end{center}
\legende{L'échelle de politesse pour demander la permission : de \eng{can} (le plus courant) à \eng{may} (le plus formel).}
\end{popfigure}

\begin{popfigure}
\begin{center}
\begin{tikzpicture}
\matrix[column sep=0.5cm, row sep=0.35cm,
  every node/.style={draw, rounded corners, align=center, font=\small, minimum height=0.85cm}] {
\node[fill=popblueL, font=\bfseries] {}; &
\node[fill=popblueL, font=\bfseries, minimum width=3.4cm] {Permission}; &
\node[fill=popblueL, font=\bfseries, minimum width=3.4cm] {Possibilité}; &
\node[fill=popblueL, font=\bfseries, minimum width=3.4cm] {Passé}; \\
\node[fill=popgreenL, font=\bfseries, minimum width=1.4cm] {can}; &
\node[minimum width=3.4cm] {\footnotesize Can I borrow your phone?}; &
\node[minimum width=3.4cm] {\footnotesize It can rain heavily in Buea.}; &
\node[minimum width=3.4cm, fill=popcream, align=center]{\footnotesize impossible \\ (utiliser could)}; \\
\node[fill=poporangeL, font=\bfseries, minimum width=1.4cm] {could}; &
\node[minimum width=3.4cm] {\footnotesize Could I leave early? (poli)}; &
\node[minimum width=3.4cm] {\footnotesize It could rain tomorrow. (hypothèse)}; &
\node[minimum width=3.4cm] {\footnotesize We could play outside till night.}; \\
\node[fill=poppurpleL, font=\bfseries, minimum width=1.4cm] {may}; &
\node[minimum width=3.4cm] {\footnotesize May I come in? (formel)}; &
\node[minimum width=3.4cm, fill=popcream] {\footnotesize emploi étudié plus tard}; &
\node[minimum width=3.4cm, fill=popcream] {\footnotesize impossible}; \\
};
\end{tikzpicture}
\end{center}
\legende{Tableau à double entrée : les trois modaux selon leurs fonctions, avec un exemple par case.}
\end{popfigure}

\courssub{Comparaison avec le français et méthode de choix}

\begin{center}
\begin{tabularx}{\linewidth}{|X|X|X|}
\hline
\rowcolor{popblueL} Français & English & Remarque \\
\hline
Je peux sortir ? (familier) & \eng{Can I go out?} & registre courant \\
\hline
Pourrais-je sortir ? & \eng{Could I go out?} & conditionnel français = could poli \\
\hline
Puis-je entrer ? (soutenu) & \eng{May I come in?} & registre formel \\
\hline
Il peut pleuvoir ici en juillet. & \eng{It can rain here in July.} & fait général \\
\hline
Il pourrait pleuvoir demain. & \eng{It could rain tomorrow.} & hypothèse \\
\hline
Je pourrai venir demain. & \eng{I will be able to come tomorrow.} & futur : pas de \eng{can} ! \\
\hline
\end{tabularx}
\end{center}

\begin{methode}[Choisir entre can, could et may pour demander la permission]
\begin{enumerate}
\item \textbf{Identifiez votre interlocuteur.} À qui parlez-vous ?
\item \textbf{Ami, camarade, frère ou sœur} $\rightarrow$ utilisez \eng{can} : \eng{Can I take your bike?}
\item \textbf{Personne que vous voulez ménager}, demande délicate, inconnu $\rightarrow$ utilisez \eng{could} : \eng{Could I ask you a question?}
\item \textbf{Autorité (professeur, proviseur, fonctionnaire) ou contexte officiel} $\rightarrow$ utilisez \eng{may} : \eng{May I see the principal, please?}
\item \textbf{Vérifiez la réponse} : on répond toujours avec \eng{can} ou \eng{may}, jamais avec \eng{could}.
\end{enumerate}
\end{methode}

\begin{exoresolu}[Exercice résolu : compléter avec can, could ou may]
Énoncé — Complétez chaque phrase avec \eng{can}, \eng{can't}, \eng{could}, \eng{couldn't}, \eng{may} ou \eng{may not}, puis traduisez.
\begin{enumerate}
\item ......... I come in, Sir? I am very sorry I am late. (formel)
\item ......... you lend me your ruler, please? (à un camarade)
\item When my grandfather was young, children ......... walk five kilometres to school.
\item Take an umbrella : it ......... rain this afternoon. (hypothèse)
\item Students ......... not use dictionaries during the test. (règlement officiel)
\item Sorry, you ......... not park here ; this place is reserved. (refus courant)
\end{enumerate}
\tcblower
Solution détaillée :
\begin{enumerate}
\item \eng{May I come in, Sir?} — On s'adresse à un professeur : registre formel $\rightarrow$ \eng{may}. « Puis-je entrer, Monsieur ? »
\item \eng{Can you lend me your ruler, please?} — Entre camarades : registre courant $\rightarrow$ \eng{can}. « Peux-tu me prêter ta règle, s'il te plaît ? » (\eng{Could} serait aussi correct, plus poli.)
\item \eng{...children could walk five kilometres to school.} — Possibilité au passé (\eng{when my grandfather was young}) $\rightarrow$ \eng{could}, passé de \eng{can}. « ...les enfants pouvaient marcher cinq kilomètres jusqu'à l'école. »
\item \eng{...it could rain this afternoon.} — Hypothèse incertaine pour cet après-midi $\rightarrow$ \eng{could}. « ...il pourrait pleuvoir cet après-midi. »
\item \eng{Students may not use dictionaries...} — Interdiction officielle dans un règlement $\rightarrow$ \eng{may not}. « Les élèves ne sont pas autorisés à utiliser de dictionnaires... »
\item \eng{...you can't park here.} — Refus courant, oral $\rightarrow$ \eng{can't}. « Désolé, vous ne pouvez pas vous garer ici. »
\end{enumerate}
\end{exoresolu}

\begin{exercice}[À vous de jouer]
\begin{enumerate}
\item Traduisez en anglais : a) « Est-ce que je pourrais utiliser tes ciseaux ? » b) « Puis-je parler au directeur ? » c) « En saison sèche, les pistes peuvent être très poussiéreuses. » d) « Quand nous étions petits, nous pouvions regarder les matchs de football gratuitement. »
\item Répondez aux demandes suivantes en donnant ou en refusant la permission : a) \eng{Could I close the door?} (acceptez) b) \eng{May I leave the room?} (refusez poliment) c) \eng{Can I have some water?} (acceptez).
\item Corrigez les erreurs : a) \emph{I will can visit you next week.} b) \emph{Could I borrow your book? — Yes, you could.} c) \emph{She cans speak English very well.} d) \emph{May I to come in?}
\end{enumerate}
\end{exercice}

\begin{corrige}[Exercice « À vous de jouer »]
\begin{enumerate}
\item a) \eng{Could I use your scissors?} b) \eng{May I speak to the headmaster?} c) \eng{In the dry season, the tracks can be very dusty.} d) \eng{When we were small, we could watch football matches for free.}
\item a) \eng{Yes, you can.} b) \eng{No, you may not. / I'm sorry, you may not.} c) \eng{Yes, of course you can.}
\item a) \eng{I will be able to visit you next week.} (\eng{can} n'a pas de futur.) b) \eng{Yes, you can.} (on ne répond jamais avec \eng{could}.) c) \eng{She can speak English very well.} (pas de \eng{-s} après un modal.) d) \eng{May I come in?} (pas de \eng{to} après un modal.)
\end{enumerate}
\end{corrige}

\begin{aretenir}[L'essentiel de la section]
\begin{itemize}
\item Permission : \eng{can} (courant) $<$ \eng{could} (poli) $<$ \eng{may} (formel). Réponses : \eng{Yes, you can/may} — \eng{No, you can't/may not}.
\item Possibilité : \eng{can} au présent (fait général), \eng{could} au passé ou pour une hypothèse.
\item Après un modal : base verbale, sans \eng{to}, sans \eng{-s}, sans \eng{do}.
\item \eng{Can} n'a ni futur ni infinitif : utilisez \eng{will be able to} et \eng{to be able to}.
\end{itemize}
\end{aretenir}

\coursec{Conseil : should et ought to}

Dans la section précédente, nous avons vu que \eng{can}, \eng{could} et \eng{may} servent à demander ou à donner la \cle{permission} et à exprimer la \cle{possibilité}. Mais la vie familiale et sociale ne se résume pas à ce qui est permis ou possible : il y a aussi ce qui est \emph{souhaitable}, ce qu'il serait \emph{bien} de faire. Quand une grand-mère de Buea dit à son petit-fils : \eng{You should respect your elders}, elle ne lui donne pas un ordre, elle lui donne un \cle{conseil}. C'est exactement le rôle des deux modaux que nous allons étudier maintenant : \cle{should} et \cle{ought to}.

\courssub{Observation : les modaux du conseil en action}

\begin{textebox}[A letter to Auntie Comfort — advice column]
Dear Auntie Comfort,

I am a sixteen-year-old girl from Bamenda. My younger brother and I quarrel all the time, and my mother is always angry with us. Last week, he broke my phone and I refused to talk to him for three days. What should I do? — Worried Sister

Dear Worried Sister,

Family peace is precious, so you should apologise first, even if your brother started the quarrel. You ought to remember that he is younger than you: elders ought to show the example. You shouldn't keep silent for days; silence destroys families. You ought not to shout at him either — speak calmly and explain your feelings. Children ought to help at home and support one another, not fight. Should you offer him something? Yes, perhaps a small gesture, like helping him with his homework. If you do this, I am sure your mother will smile again.

Auntie Comfort
\end{textebox}

\textbf{Glossaire :} \emph{to quarrel} (v.) = se disputer ; \emph{angry} (adj.) = en colère ; \emph{to apologise} (v.) = s'excuser ; \emph{elders} (n.) = les aînés ; \emph{to shout} (v.) = crier ; \emph{a gesture} (n.) = un geste (symbolique).

\textbf{Questions d'observation :}
\begin{enumerate}
\item Souligne dans la lettre tous les verbes précédés de \eng{should} ou \eng{ought to}.
\item Ces phrases donnent-elles des ordres ou des conseils ? Justifie ta réponse.
\item Quelle forme suit \eng{ought} ? Et quelle forme suit \eng{should} ?
\end{enumerate}

\textbf{Constats :} \eng{should} et \eng{ought to} sont suivis de la \cle{base verbale} (l'infinitif sans \eng{to}) ; ils expriment un conseil, une recommandation, jamais un ordre strict. Remarque bien : \eng{ought} garde toujours son \eng{to}, alors que \eng{should} est directement suivi du verbe.

\courssub{Should : le modal du conseil}

\begin{propriete}[Should + base verbale = conseil ou recommandation]
\eng{Should} est un auxiliaire modal \textbf{invariable} (jamais de \emph{-s} à la troisième personne) suivi de la \textbf{base verbale} (infinitif sans \eng{to}). Il exprime un \cle{conseil}, une \cle{recommandation} ou ce qui est \emph{souhaitable}.

Forme : \textbf{sujet + should + base verbale}.

\eng{You should apologise to your sister.} « Tu devrais t'excuser auprès de ta sœur. »

\eng{He should apologise to his sister.} « Il devrait s'excuser auprès de sa sœur. » (pas de \emph{-s} sur \eng{should} !)
\end{propriete}

\begin{exemplebox}[Should en contexte familial et social]
\begin{itemize}
\item \eng{We should visit our grandparents more often.} « Nous devrions rendre visite à nos grands-parents plus souvent. »
\item \eng{You should eat more vegetables and less fried food.} « Tu devrais manger plus de légumes et moins de friture. »
\item \eng{Students should revise their lessons every evening.} « Les élèves devraient réviser leurs leçons chaque soir. »
\item \eng{She should listen to her parents' advice.} « Elle devrait écouter les conseils de ses parents. »
\end{itemize}
\end{exemplebox}

\courssub{Ought to : le quasi-synonyme de should}

\begin{propriete}[Ought to + base verbale]
\eng{Ought to} a \textbf{le même sens} que \eng{should} (conseil, recommandation), mais il est \textbf{légèrement plus formel}. Il est très fréquent pour exprimer ce qui est \emph{socialement souhaitable} ou moralement attendu.

Forme : \textbf{sujet + ought to + base verbale} (invariable, sans \emph{-s}).

\eng{You ought to respect your elders.} « Tu devrais respecter tes aînés. »

\eng{Children ought to help at home.} « Les enfants devraient aider à la maison. » (devoir social, rôle attendu dans la famille)
\end{propriete}

\begin{exemplebox}[Ought to pour le devoir social et moral]
\begin{itemize}
\item \eng{We ought to take care of our ageing parents.} « Nous devrions prendre soin de nos parents vieillissants. »
\item \eng{Young people ought to greet older people politely.} « Les jeunes devraient saluer les personnes âgées poliment. »
\item \eng{Every citizen ought to keep the neighbourhood clean.} « Chaque citoyen devrait garder le quartier propre. »
\item \eng{You ought to thank your host before leaving.} « Tu devrais remercier ton hôte avant de partir. »
\end{itemize}
\end{exemplebox}

\begin{attention}[Le piège majeur : ought to garde toujours son to]
Contrairement à tous les autres modaux (\eng{can}, \eng{must}, \eng{should}...), \eng{ought} est \textbf{toujours suivi de} \eng{to}. Mais attention : après ce \eng{to}, le verbe reste à la \textbf{base verbale}, sans second \eng{to}.

\begin{center}
\begin{tabularx}{\linewidth}{|X|X|}
\hline
\rowcolor{popblueL} Incorrect & Correct \\
\hline
You ought go. & You ought \textbf{to} go. \\
\hline
You ought to \textbf{to go}. & You ought to go. \\
\hline
He oughts to help. & He ought to help. (pas de \emph{-s}) \\
\hline
\end{tabularx}
\end{center}

Astuce de mémorisation : \eng{ought to} forme un \emph{bloc inséparable} — prononce-le « ô-teu » comme un seul mot.
\end{attention}

\courssub{Formes négative et interrogative}

\begin{propriete}[Négation et interrogation]
\textbf{Négation :} on place \eng{not} après le modal.
\begin{itemize}
\item \eng{should not} $\rightarrow$ contraction : \eng{shouldn't} (« chou-dent »)
\item \eng{ought not to} $\rightarrow$ contraction orale britannique : \eng{oughtn't to}
\end{itemize}

\eng{You shouldn't quarrel with your siblings.} « Tu ne devrais pas te disputer avec tes frères et sœurs. »

\eng{He ought not to speak to his mother like that.} « Il ne devrait pas parler à sa mère ainsi. »

\textbf{Interrogation :} inversion du sujet et de \eng{should} : \eng{Should I...?}

\eng{Should I bring a gift to the ceremony?} « Devrais-je apporter un cadeau à la cérémonie ? »

\eng{Ought to} est \textbf{rarement employé à la forme interrogative} ; on lui préfère \eng{should} : on dit \eng{Should I call her?} plutôt que \eng{Ought I to call her?} (très soutenu, vieilli).
\end{propriete}

\begin{center}
\begin{tabular}{|l|l|l|}
\hline
\rowcolor{popblueL} \textbf{Forme} & \textbf{should} & \textbf{ought to} \\
\hline
Affirmative & You should help. & You ought to help. \\
\hline
Négative & You shouldn't help. / shouldn't & You ought not to help. \\
\hline
Interrogative & Should you help? & (rare : Ought you to help?) \\
\hline
Réponse courte & Yes, I should. / No, I shouldn't. & Yes, I ought to. / No, I ought not to. \\
\hline
\end{tabular}
\end{center}

\begin{savaistu}[Oughtn't to : une contraction à connaître, pas à utiliser]
À l'oral britannique familier, \eng{ought not to} se contracte parfois en \eng{oughtn't to} (« ô-tent teu ») : \eng{You oughtn't to say that!} Cette forme sonne aujourd'hui un peu désuète. À l'écrit — et surtout à l'examen du Probatoire ou du Baccalauréat — préfère toujours la forme complète \eng{ought not to}, ou plus simplement \eng{shouldn't}, qui est le choix le plus naturel et le plus sûr.
\end{savaistu}

\courssub{Conseil ou obligation ? La nuance essentielle}

\begin{aretenir}[Should/ought to = conseil ; must = obligation]
\eng{Should} et \eng{ought to} expriment un \textbf{conseil} : on peut ne pas le suivre. \eng{Must} exprime une \textbf{obligation} : on n'a pas le choix.

\eng{You should wear a uniform at the ceremony.} « Tu devrais porter un uniforme. » (c'est conseillé, mais pas obligatoire)

\eng{You must wear a uniform at school.} « Tu dois porter l'uniforme à l'école. » (c'est le règlement)
\end{aretenir}

\begin{popfigure}
\begin{center}
\begin{tikzpicture}
\node[fill=popink, text=white, rounded corners=2pt, minimum width=9.5cm, minimum height=1.05cm, align=center, font=\bfseries] (a) at (0,0) {MUST — obligation forte : You must obey the school rules.};
\node[fill=poporange, text=white, rounded corners=2pt, minimum width=7.5cm, minimum height=1.05cm, align=center, font=\bfseries] (b) at (0,-1.5) {SHOULD / OUGHT TO — conseil : You should revise daily.};
\node[fill=popgreen, text=white, rounded corners=2pt, minimum width=5.5cm, minimum height=1.05cm, align=center, font=\bfseries] (c) at (0,-3) {COULD — simple suggestion : You could ask your uncle.};
\node[align=center, text=popmuted, font=\itshape] at (0,-4) {Force décroissante : de l'ordre vers la simple idée};
\draw[-{Stealth}, thick, popmuted] (5.6,0) -- (5.6,-3) node[midway, right, align=left, font=\itshape] {de plus en plus\\faible};
\end{tikzpicture}
\end{center}
\legende{L'entonnoir de la force modale : obligation (must), conseil (should / ought to), suggestion (could).}
\end{popfigure}

\begin{attention}[Le « devoir » français a deux visages]
En français, le verbe \emph{devoir} sert à la fois pour le conseil et pour l'obligation. En anglais, il faut \textbf{choisir} :

\begin{center}
\begin{tabularx}{\linewidth}{|X|X|X|}
\hline
\rowcolor{popblueL} Français & English & Remarque \\
\hline
Tu devrais dormir. (conseil) & You \textbf{should} sleep. & pas \eng{must} ! \\
\hline
Tu dois dormir. (ordre du médecin) & You \textbf{must} sleep. & obligation réelle \\
\hline
Tu ne devrais pas mentir. & You \textbf{shouldn't} lie. & conseil négatif \\
\hline
Tu ne dois pas mentir. (interdiction) & You \textbf{mustn't} lie. & interdiction \\
\hline
\end{tabularx}
\end{center}

Erreur typique du francophone : traduire « Tu devrais dormir » par \eng{You must sleep}, ce qui transforme un gentil conseil en ordre autoritaire. Demande-toi toujours : \emph{est-ce un conseil ou une obligation ?}
\end{attention}

\courssub{Exercice résolu et entraînement}

\begin{exoresolu}[Complète avec should, shouldn't, ought to ou ought not to]
Énoncé — Complète chaque phrase avec le modal qui convient selon le sens entre parenthèses.

1. You \ldots\ldots\ldots\ldots respect your elders. (conseil, formel)
2. You \ldots\ldots\ldots\ldots quarrel with your siblings. (conseil négatif)
3. \ldots\ldots\ldots\ldots I bring a gift to the ceremony? (question)
4. Children \ldots\ldots\ldots\ldots help their parents at home. (devoir social)
5. He is ill; he \ldots\ldots\ldots\ldots see a doctor. (conseil simple)

\tcblower
Solution détaillée :

1. \eng{You \textbf{ought to} respect your elders.} — Le contexte « formel » appelle \eng{ought to}, le modal du devoir social et moral. \eng{Should} serait aussi correct, mais moins soutenu.

2. \eng{You \textbf{shouldn't} quarrel with your siblings.} — Conseil négatif : \eng{shouldn't} est la forme la plus naturelle. \eng{Ought not to} serait possible mais plus lourd.

3. \eng{\textbf{Should} I bring a gift to the ceremony?} — À l'interrogative, on utilise \eng{should} avec inversion ; \eng{ought to} est quasi inemployé dans les questions.

4. \eng{Children \textbf{ought to} help their parents at home.} — Il s'agit d'un rôle socialement attendu : \eng{ought to} convient parfaitement.

5. \eng{He is ill; he \textbf{should} see a doctor.} — Conseil personnel simple : \eng{should} est le choix naturel. Remarque : pas de \emph{-s} à \eng{should}, même avec \eng{he}.
\end{exoresolu}

\begin{exercice}[À toi de jouer]
1. Traduis en anglais : a) « Tu devrais t'excuser auprès de ta sœur. » b) « Les enfants devraient aider à la maison. » c) « Il ne devrait pas parler à sa mère ainsi. » d) « Devrais-je apporter un cadeau ? »

2. Corrige les erreurs : a) \eng{You ought go now.} b) \eng{She shoulds study more.} c) \eng{You ought to to apologise.} d) \eng{Should I to call my grandmother?}

3. Choisis entre \eng{should} et \eng{must} : a) « Tu devrais manger moins de sucre » (conseil d'un ami). b) « Tu dois porter l'uniforme » (règlement du lycée).
\end{exercice}

\begin{corrige}[Exercice « À toi de jouer »]
1. a) \eng{You should apologise to your sister.} b) \eng{Children ought to help at home.} c) \eng{He ought not to speak to his mother like that.} (ou \eng{shouldn't speak}) d) \eng{Should I bring a gift?}

2. a) \eng{You ought \textbf{to} go now.} (il manque \eng{to}) b) \eng{She \textbf{should} study more.} (modal invariable, sans \emph{-s}) c) \eng{You ought to \textbf{apologise}.} (base verbale après \eng{ought to}, pas de second \eng{to}) d) \eng{Should I \textbf{call} my grandmother?} (base verbale directement après \eng{should})

3. a) \eng{You should eat less sugar.} (conseil) b) \eng{You must wear the uniform.} (obligation)
\end{corrige}

\begin{aretenir}[L'essentiel de la section]
\begin{itemize}
\item \eng{Should} + base verbale = conseil, recommandation : \eng{You should respect your elders.}
\item \eng{Ought to} + base verbale = même sens, plus formel, idéal pour le devoir social : \eng{Children ought to help at home.}
\item Négation : \eng{shouldn't} / \eng{ought not to} ; interrogation : \eng{Should I...?} (jamais \eng{ought to} en pratique).
\item Conseil $\neq$ obligation : « tu devrais » = \eng{should} ; « tu dois » = \eng{must}.
\item \eng{Ought} garde toujours son \eng{to}, mais le verbe qui suit reste à la base verbale.
\end{itemize}
\end{aretenir}

\coursec{Obligation, interdiction et absence d'obligation : must, mustn't, needn't, don't have to}

Dans la section précédente, nous avons vu comment exprimer le \cle{conseil} avec \eng{should} et \eng{ought to} : \eng{You should respect your parents.} (Tu devrais respecter tes parents.) Mais le conseil n'est qu'une recommandation : on peut le suivre ou non. Que se passe-t-il lorsqu'une règle \emph{s'impose} à nous ? Lorsque le proviseur fixe l'heure d'arrivée, lorsque le règlement intérieur interdit les téléphones, ou au contraire lorsque quelque chose n'est \emph{pas obligatoire} ? C'est tout l'objet de cette section : exprimer l'\cle{obligation}, l'\cle{interdiction} et l'\cle{absence d'obligation} en anglais.

\courssub{Observation : les règles d'un lycée camerounais}

\begin{textebox}[School rules — Government High School, Buea]
Dear students, welcome to a new school year. Please read the following rules carefully.

First, all students \textbf{must} wear the school uniform from Monday to Friday. You \textbf{must} arrive at school before 7.30 a.m.; the gate closes at 7.45. During lessons, you \textbf{mustn't} use your mobile phones, and you \textbf{mustn't} leave the classroom without the teacher's permission.

On Wednesdays, however, you \textbf{don't have to} wear your uniform: you may come in sports clothes because of the games afternoon. You \textbf{needn't} bring your own food either; the school canteen provides lunch for everyone.

Finally, remember that students who live far away \textbf{have to} wake up very early, so the school has organised a bus service. Last year, students \textbf{had to} pay for this service, but this year it is free.

The Principal
\end{textebox}

\textbf{Glossaire :} \emph{rules} (n., règles) ; \emph{to wear} (v., porter) ; \emph{gate} (n., portail) ; \emph{canteen} (n., cantine) ; \emph{to provide} (v., fournir) ; \emph{to wake up} (v., se réveiller) ; \emph{free} (adj., gratuit).

\textbf{Questions de compréhension :}
\begin{enumerate}
\item What must students do every morning?
\item What are students forbidden to do during lessons?
\item On which day is the uniform not compulsory?
\item Who had to pay for the bus service last year?
\end{enumerate}

\textbf{Observation grammaticale.} Relevons les formes employées :
\begin{itemize}
\item \eng{must wear}, \eng{must arrive} → une \cle{obligation} ;
\item \eng{mustn't use}, \eng{mustn't leave} → une \cle{interdiction} ;
\item \eng{don't have to wear}, \eng{needn't bring} → une \cle{absence d'obligation} (ce n'est pas nécessaire) ;
\item \eng{have to wake up}, \eng{had to pay} → une obligation extérieure, conjuguée dans le temps.
\end{itemize}
Remarquons que tous ces auxiliaires sont suivis de la \cle{base verbale} (l'infinitif sans \eng{to}) : \eng{must wear}, \eng{mustn't use}, \eng{needn't bring}.

\courssub{Must : l'obligation forte}

\begin{propriete}[Must + base verbale = obligation]
\eng{Must} est un auxiliaire modal \textbf{invariable} (pas de \eng{-s} à la 3\textsuperscript{e} personne, pas d'infinitif, pas de participe). Il exprime une obligation forte, imposée par le locuteur, par une règle ou par la loi.
\begin{itemize}
\item Forme : \textbf{sujet + must + base verbale}.
\item Négation : \textbf{must not / mustn't} (attention : sens d'interdiction, voir plus bas).
\item Interrogation (rare) : \eng{Must I come?} (Dois-je venir ?)
\end{itemize}
\eng{Students must arrive at 7.30.} (Les élèves doivent arriver à 7 h 30.)\\
\eng{You must respect your teachers.} (Tu dois respecter tes professeurs.)\\
\eng{She must finish her homework before dinner.} (Elle doit finir ses devoirs avant le dîner.) — pas de \eng{musts} !
\end{propriete}

\begin{attention}[Must n'a ni passé ni futur]
\eng{Must} ne se conjugue pas : on ne peut pas dire \eng{musted} ni \eng{will must}. Pour exprimer l'obligation au passé ou au futur, on utilise \eng{have to} : \eng{had to} (passé), \eng{will have to} (futur). Voir la frise des temps plus bas.
\end{attention}

\courssub{Mustn't : l'interdiction}

\begin{propriete}[Mustn't = il est interdit de]
La négation de \eng{must} ne signifie pas « ne pas être obligé » mais « \textbf{il est interdit de} », « \textbf{surtout ne pas} ». C'est une défense absolue.
\begin{itemize}
\item \eng{Students mustn't use phones in class.} (Il est interdit aux élèves d'utiliser leur téléphone en classe.)
\item \eng{You mustn't cheat in an exam.} (Tu ne dois pas tricher à un examen — c'est interdit.)
\item \eng{You mustn't insult anybody.} (Tu ne dois insulter personne.)
\item \eng{Visitors mustn't enter the laboratory.} (Les visiteurs ne doivent pas entrer dans le laboratoire.)
\end{itemize}
\end{propriete}

\courssub{Don't have to et needn't : l'absence d'obligation}

\begin{propriete}[Don't have to / needn't = ce n'est pas nécessaire]
Pour dire qu'une action \textbf{n'est pas obligatoire} (on peut la faire ou non, au choix), l'anglais utilise :
\begin{itemize}
\item \textbf{don't have to / doesn't have to + base verbale} : \eng{You don't have to wear a tie.} (Tu n'es pas obligé de porter une cravate.)
\item \textbf{needn't + base verbale} (need modal) : \eng{You needn't wear a tie.} (même sens)
\end{itemize}
\eng{You needn't bring food; everything is provided.} (Tu n'as pas besoin d'apporter à manger ; tout est fourni.)\\
\eng{You don't have to come early.} (Tu n'es pas obligé de venir tôt.)\\
\eng{She doesn't have to work on Saturdays.} (Elle n'est pas obligée de travailler le samedi.)
\end{propriete}

\begin{attention}[L'erreur n°1 des francophones : mustn't ≠ pas obligé]
En français, « tu ne dois pas venir » est ambigu : cela peut vouloir dire « tu n'es pas obligé de venir » ou « il est interdit que tu viennes ». En anglais, la distinction est \textbf{obligatoire} :
\begin{itemize}
\item « Tu n'es pas obligé de venir. » = \eng{You don't have to come.} ou \eng{You needn't come.}
\item « Il est interdit que tu viennes. » = \eng{You mustn't come.}
\end{itemize}
Traduire « tu n'es pas obligé » par \eng{mustn't} transforme une simple dispense en interdiction sévère !
\end{attention}

\begin{methode}[Interdit ou simplement pas nécessaire ?]
Avant de traduire une phrase négative avec « devoir », pose-toi une seule question :
\begin{enumerate}
\item L'action est-elle \textbf{interdite} (défendue, punie) ? → utilise \textbf{mustn't}. \eng{You mustn't smoke here.} (Il est interdit de fumer ici.)
\item L'action est-elle simplement \textbf{pas nécessaire} (libre choix) ? → utilise \textbf{don't have to} ou \textbf{needn't}. \eng{You needn't cook tonight; I have prepared dinner.} (Tu n'as pas besoin de cuisiner ce soir ; j'ai préparé le dîner.)
\end{enumerate}
\end{methode}

\begin{popfigure}
\begin{center}
\begin{tikzpicture}[
  box/.style={draw=popline, rounded corners=4pt, fill=#1, text width=4.6cm, align=center, inner sep=6pt, font=\small},
  head/.style={font=\bfseries\small, align=center}
]
\node[head, text=popdark] (t1) at (-3.2,3.4) {Obligatoire};
\node[head, text=popdark] (t2) at (3.2,3.4) {Interdit};
\node[box=popblueL, align=center] (m) at (-3.2,2.2){\textbf{must}\\ \eng{You must wear your uniform.}\\ (Tu dois porter ton uniforme.)};
\node[box=poppinkL, align=center] (mn) at (3.2,2.2){\textbf{mustn't}\\ \eng{You mustn't cheat.}\\ (Il est interdit de tricher.)};
\node[head, text=popdark] (t3) at (0,-0.2) {Pas obligatoire (libre choix)};
\node[box=popgreenL, align=center] (dht) at (-3.2,-1.4){\textbf{don't have to}\\ \eng{You don't have to come early.}\\ (Tu n'es pas obligé de venir tôt.)};
\node[box=popgreenL, align=center] (nd) at (3.2,-1.4){\textbf{needn't}\\ \eng{You needn't bring food.}\\ (Tu n'as pas besoin d'apporter à manger.)};
\draw[-{Stealth}, very thick, poporange] (mn.south) -- (nd.north) node[midway, right, align=left, font=\footnotesize\bfseries, text=poporange] {ATTENTION\\ sens opposés !};
\draw[-{Stealth}, thick, popmuted] (m.south) -- (dht.north) node[midway, left, font=\footnotesize, text=popmuted] {négation de l'obligation};
\end{tikzpicture}
\end{center}
\legende{Obligatoire (\eng{must}), interdit (\eng{mustn't}) et pas obligatoire (\eng{don't have to} / \eng{needn't}) : attention, \eng{mustn't} et \eng{don't have to} ont des sens opposés.}
\end{popfigure}

\courssub{Have to : le semi-modal de l'obligation extérieure}

\eng{Have to} exprime également l'obligation, mais il se comporte comme un \textbf{verbe ordinaire} : il prend l'auxiliaire \eng{do} aux formes négative et interrogative, un \eng{-s} à la 3\textsuperscript{e} personne, et il possède un passé et un futur. C'est pourquoi on l'appelle \cle{semi-modal}.

\begin{propriete}[Conjugaison de have to]
\begin{center}
\begin{tabular}{|l|l|l|}
\hline
\rowcolor{popblueL} Temps & Forme & Exemple \\
\hline
Présent (I, you, we, they) & have to & \eng{I have to go home now.} \\
\hline
Présent (he, she, it) & has to & \eng{She has to look after her little brother.} \\
\hline
Passé (toutes personnes) & had to & \eng{We had to wait for the bus.} \\
\hline
Futur & will have to & \eng{You will have to work harder.} \\
\hline
Négation présent & don't / doesn't have to & \eng{He doesn't have to pay.} \\
\hline
Interrogation & Do/Does ... have to ? & \eng{Do you have to leave?} \\
\hline
\end{tabular}
\end{center}
Traductions : \eng{She has to look after her little brother.} (Elle doit garder son petit frère.) — \eng{We had to wait for the bus.} (Nous avons dû attendre le bus.) — \eng{You will have to work harder.} (Tu devras travailler davantage.)
\end{propriete}

\begin{aretenir}[Must ou have to ?]
La différence est souvent légère : \eng{must} exprime une obligation ressentie ou imposée par le locuteur (\eng{You must listen to me!} — c'est moi qui l'exige), tandis que \eng{have to} exprime une obligation extérieure (règle, loi, circonstances) : \eng{In Cameroon, drivers have to drive on the right.} (Au Cameroun, les conducteurs doivent rouler à droite.) Au passé et au futur, seul \eng{have to} est possible.
\end{aretenir}

\begin{popfigure}
\begin{center}
\begin{tikzpicture}
\draw[-{Stealth}, thick, popline] (-5.2,0) -- (5.2,0);
\foreach \x/\lab in {-3.4/{Passé}, 0/{Présent}, 3.4/{Futur}} {
  \fill[popblue] (\x,0) circle (2.5pt);
  \node[above, font=\small\bfseries, text=popdark] at (\x,0.05) {\lab};
}
\node[below, align=center, font=\small, text=popink] at (-3.4,-0.1) {\textbf{had to}\\ \eng{I had to walk to school.}\\ (J'ai dû aller à l'école à pied.)};
\node[below, align=center, font=\small, text=popink] at (0,-0.1) {\textbf{have / has to}\\ \eng{She has to help her mother.}\\ (Elle doit aider sa mère.)};
\node[below, align=center, font=\small, text=popink] at (3.4,-0.1) {\textbf{will have to}\\ \eng{They will have to repeat the class.}\\ (Ils devront redoubler.)};
\end{tikzpicture}
\end{center}
\legende{Frise des temps de \eng{have to} : \eng{had to} (passé), \eng{have/has to} (présent), \eng{will have to} (futur). \eng{Must} n'existe qu'au présent.}
\end{popfigure}

\courssub{Need : verbe modal ou verbe ordinaire ?}

\eng{Need} a une double vie en anglais, ce qui déroute beaucoup les francophones.

\begin{propriete}[Les deux emplois de need]
\textbf{1. Need modal} (rare, surtout à la forme négative et interrogative) : \eng{needn't + base verbale}, sans \eng{to}, sans \eng{-s}, sans auxiliaire \eng{do}.
\begin{itemize}
\item \eng{You needn't bring your books tomorrow.} (Tu n'as pas besoin d'apporter tes livres demain.)
\item \eng{Need I say more?} (Ai-je besoin d'en dire plus ?)
\end{itemize}
\textbf{2. Need verbe ordinaire} : \eng{need to + base verbale}, conjugué normalement (\eng{needs}, \eng{needed}), avec \eng{do} à la négation et à l'interrogation.
\begin{itemize}
\item \eng{You need to work harder.} (Tu as besoin de travailler davantage.)
\item \eng{She needs to see a doctor.} (Elle a besoin de voir un médecin.)
\item \eng{Do you need to leave now?} (As-tu besoin de partir maintenant ?)
\end{itemize}
\end{propriete}

\begin{attention}[Need modal : presque uniquement à la forme négative]
Le \eng{need} modal ne s'emploie guère qu'à la forme \textbf{négative} (\eng{needn't}) et parfois \textbf{interrogative} (\eng{Need I...?}). À l'affirmatif, on utilise toujours le verbe ordinaire : on dit \eng{You need to work harder.} et jamais \eng{You need work harder.} Attention aussi : \eng{needn't} est suivi de la base verbale \textbf{sans} \eng{to} (\eng{You needn't to come.} est faux).
\end{attention}

\courssub{Tableau récapitulatif et comparaison avec le français}

\begin{center}
\begin{tabularx}{\linewidth}{|l|X|X|}
\hline
\rowcolor{popblueL} Forme & Sens & Exemple et traduction \\
\hline
must + BV & obligation forte & \eng{You must wear your uniform.} (Tu dois porter ton uniforme.) \\
\hline
mustn't + BV & interdiction & \eng{You mustn't cheat.} (Il est interdit de tricher.) \\
\hline
don't have to + BV & absence d'obligation & \eng{You don't have to come early.} (Tu n'es pas obligé de venir tôt.) \\
\hline
needn't + BV & absence d'obligation & \eng{You needn't cook tonight.} (Tu n'as pas besoin de cuisiner ce soir.) \\
\hline
have/has to + BV & obligation extérieure & \eng{She has to look after her brother.} (Elle doit garder son frère.) \\
\hline
had to + BV & obligation au passé & \eng{We had to wait.} (Nous avons dû attendre.) \\
\hline
\end{tabularx}
\end{center}

\begin{center}
\begin{tabularx}{\linewidth}{|X|X|X|}
\hline
\rowcolor{popblueL} Français & English & Remarque \\
\hline
Tu dois partir. & You must leave. / You have to leave. & must = locuteur ; have to = circonstance \\
\hline
Tu ne dois pas fumer. (interdit) & You mustn't smoke. & interdiction absolue \\
\hline
Tu n'es pas obligé de partir. & You don't have to leave. / You needn't leave. & jamais \eng{mustn't} ici ! \\
\hline
Elle doit garder son frère. & She has to look after her brother. & \eng{has to} à la 3\textsuperscript{e} personne \\
\hline
Nous avons dû attendre. & We had to wait. & passé : \eng{had to}, jamais \eng{musted} \\
\hline
\end{tabularx}
\end{center}

\courssub{Entraînement}

\begin{exoresolu}[Choisir le bon modal]
Complète avec \eng{must}, \eng{mustn't}, \eng{don't have to} ou \eng{needn't}.
\begin{enumerate}
\item You \ldots\ touch that wire; it is dangerous.
\item You \ldots\ buy a ticket; the concert is free.
\item All passengers \ldots\ show their identity cards at the airport.
\item We \ldots\ hurry; the train leaves in two hours.
\item Students \ldots\ be late for the examination.
\end{enumerate}
\tcblower
\textbf{Solution détaillée :}
\begin{enumerate}
\item \eng{mustn't} — c'est \emph{dangereux}, donc c'est une \textbf{interdiction} : \eng{You mustn't touch that wire.} (Il est interdit de toucher ce fil.)
\item \eng{don't have to} / \eng{needn't} — le concert est \emph{gratuit}, donc ce n'est \textbf{pas nécessaire} : \eng{You don't have to buy a ticket.} (Tu n'es pas obligé d'acheter un billet.)
\item \eng{must} — règle administrative, \textbf{obligation} : \eng{All passengers must show their identity cards.} (Tous les passagers doivent présenter leur pièce d'identité.)
\item \eng{don't have to} / \eng{needn't} — le train part dans deux heures, il n'y a \textbf{pas nécessité} de se dépêcher : \eng{We needn't hurry.} (Nous n'avons pas besoin de nous dépêcher.)
\item \eng{mustn't} — être en retard à un examen est \textbf{interdit} : \eng{Students mustn't be late.} (Il est interdit aux élèves d'être en retard.)
\end{enumerate}
Méthode : pour chaque phrase, on s'est demandé « interdit ou pas nécessaire ? » avant de choisir la forme négative.
\end{exoresolu}

\begin{exercice}[À toi de jouer]
\begin{enumerate}
\item Traduis en anglais : a) Tu n'es pas obligé de répondre à cette question. b) Il est interdit de parler pendant l'examen. c) Ma sœur doit aider notre mère au marché. d) L'année dernière, nous avons dû porter l'uniforme le samedi.
\item Transforme à la forme négative correcte (interdiction ou absence d'obligation selon le sens) : \eng{You must bring your dictionary.} (sens : ce n'est pas nécessaire) puis \eng{You can use your phone.} (sens : c'est interdit).
\item Complète avec la bonne forme de \eng{have to} : a) Yesterday, she \ldots\ stay at home. b) Tomorrow, we \ldots\ take an exam. c) \ldots\ your brother \ldots\ work on Sundays?
\end{enumerate}
\end{exercice}

\begin{corrige}[Exercice « À toi de jouer »]
\begin{enumerate}
\item a) \eng{You don't have to answer this question.} / \eng{You needn't answer this question.} b) \eng{You mustn't talk during the exam.} c) \eng{My sister has to help our mother at the market.} d) \eng{Last year, we had to wear the uniform on Saturdays.}
\item \eng{You don't have to bring your dictionary.} / \eng{You needn't bring your dictionary.} (pas nécessaire) ; \eng{You mustn't use your phone.} (interdit).
\item a) \eng{had to} ; b) \eng{will have to} ; c) \eng{Does your brother have to work on Sundays?}
\end{enumerate}
\end{corrige}

\begin{aretenir}[L'essentiel de la section]
\begin{itemize}
\item \eng{Must} + base verbale = obligation forte ; invariable, sans passé ni futur.
\item \eng{Mustn't} = \textbf{interdiction} (« il est interdit de »).
\item \eng{Don't have to} et \eng{needn't} = \textbf{absence d'obligation} (« ce n'est pas nécessaire ») — jamais \eng{mustn't} pour ce sens !
\item \eng{Have to} se conjugue : \eng{has to}, \eng{had to}, \eng{will have to}, avec l'auxiliaire \eng{do}.
\item \eng{Need} modal : presque uniquement négatif (\eng{needn't} + base verbale sans \eng{to}) ; à l'affirmatif, verbe ordinaire \eng{need to do}.
\end{itemize}
\end{aretenir}

\coursec{Déduction : must et can't logiques}

Dans la section précédente, nous avons vu \eng{must} comme marqueur de l'\emph{obligation} : \eng{Students must wear their uniform.} (Les élèves doivent porter leur uniforme.) Mais le même mot \eng{must} peut exprimer une idée très différente : une \emph{certitude logique}, une conclusion que l'on tire à partir d'indices. C'est ce deuxième emploi, essentiel au Bac, que nous allons maintenant étudier, avec son contraire \eng{can't}.

\courssub{Observation : un dialogue en famille}

\begin{textebox}[At home in Yaoundé --- a family dialogue]
\textbf{Saturday evening. Mrs Ngo Bell and her son Junior are at home.}

--- Junior, listen! Someone is knocking at the door. It \textbf{must be} your father. He always comes back at this time.

--- No, Mum, it \textbf{can't be} Dad. His car is not in front of the house, and he took the car this morning.

--- You are right. Then who is it? It \textbf{could be} Auntie Rose. She said she might visit us this weekend.

--- Look through the window! The visitor is wearing a nurse's uniform. It \textbf{must be} Cousin Sandrine. She works at the Central Hospital and she finishes her shift at six.

--- Yes! And she is carrying a big bag. She \textbf{must have} brought us some mangoes from her farm!

--- Wait... the bag is very small, Junior. It \textbf{can't be} full of mangoes!

--- Well, it \textbf{could be} groundnuts then. Open the door, Mum, I am hungry!
\end{textebox}

\begin{definition}[Glossaire]
\begin{itemize}
\item \eng{to knock at the door} (v.) : frapper à la porte
\item \eng{auntie} (n.) : tatie, tante (familier)
\item \eng{a shift} (n.) : un poste, un horaire de travail (infirmière, etc.)
\item \eng{groundnuts} (n.) : arachides
\item \eng{to look through the window} (v.) : regarder par la fenêtre
\end{itemize}
\end{definition}

\textbf{Questions de compréhension :}
\begin{enumerate}
\item Why does Junior think the visitor is not his father?
\item Who is the visitor? What clues help the family to identify her?
\item Find in the dialogue two deductions with \eng{must}, two with \eng{can't} and two with \eng{could}.
\end{enumerate}

\courssub{La règle : must et can't de déduction}

Dans ce dialogue, personne n'oblige personne à rien. Les modaux servent à \emph{conclure} à partir d'indices observables : la voiture absente, l'uniforme d'infirmière, la taille du sac. C'est la \cle{déduction logique}.

\begin{propriete}[Must de déduction : « il est certain que »]
\textbf{Forme :} sujet + \cle{must} + base verbale (sans \eng{to}).\par
\textbf{Emploi :} le locuteur est presque sûr (90--100 \%) que quelque chose est vrai, parce qu'il observe des indices. En français, on traduit par « il doit être... », « c'est sûrement... », « il est certain que... ».\par
\eng{The lights are on; someone must be at home.} (Les lumières sont allumées ; quelqu'un doit être à la maison.)\par
\eng{She is crying; she must be sad.} (Elle pleure ; elle doit être triste.)
\end{propriete}

\begin{propriete}[La déduction négative se dit can't, jamais mustn't]
\textbf{Forme :} sujet + \cle{can't} + base verbale.\par
\textbf{Emploi :} le locuteur est sûr que quelque chose est \emph{impossible} (0 \%), toujours à partir d'indices. En français : « il ne peut pas être... », « ce n'est pas possible que... ».\par
\eng{He can't be at school; it's Sunday.} (Il ne peut pas être à l'école ; c'est dimanche.)\par
\eng{He can't be a doctor; he is only seventeen.} (Il ne peut pas être médecin ; il n'a que dix-sept ans.)\par
\textbf{Exception capitale :} \eng{mustn't} n'exprime \emph{jamais} la déduction négative ; il exprime l'\emph{interdiction} (\eng{You mustn't smoke here.} = Il est interdit de fumer ici.).
\end{propriete}

\begin{exemplebox}[Exemples traduits]
\begin{itemize}
\item \eng{Listen! That must be the school bell.} (Écoute ! Ce doit être la cloche de l'école.)
\item \eng{This can't be Mary's bag; hers is red.} (Ça ne peut pas être le sac de Mary ; le sien est rouge.)
\item \eng{They could be at the funeral in the village.} (Ils sont peut-être aux funérailles au village.)
\item \eng{Grandmother must be tired; she has worked on the farm all day.} (Grand-mère doit être fatiguée ; elle a travaillé au champ toute la journée.)
\item \eng{That man can't be a teacher; he doesn't know the name of the school.} (Cet homme ne peut pas être enseignant ; il ne connaît pas le nom de l'école.)
\item \eng{The baby is crying. She may be hungry.} (Le bébé pleure. Elle a peut-être faim.)
\end{itemize}
\end{exemplebox}

\courssub{Les degrés de certitude : must, could/may, can't}

Entre la certitude (\eng{must}) et l'impossibilité (\eng{can't}), il existe un degré intermédiaire : la \emph{possibilité}, exprimée par \cle{could} ou \cle{may} (environ 50 \% de certitude). C'est un complément très utile pour le Bac.

\begin{center}
\begin{tabularx}{\linewidth}{|l|X|X|}\hline
\rowcolor{popblueL} \textbf{Modal} & \textbf{Degré de certitude} & \textbf{Traduction française} \\ \hline
\eng{must} & 90--100 \% sûr & il doit être, c'est sûrement \\ \hline
\eng{could / may} & environ 50 \% & il se peut que, peut-être \\ \hline
\eng{can't} & 0 \% (impossible) & il ne peut pas être, impossible \\ \hline
\end{tabularx}
\end{center}

\begin{popfigure}
\begin{center}
\begin{tikzpicture}
\draw[-{Stealth}, very thick, popdark] (0,0) -- (11,0);
\foreach \x/\p in {0/0, 2.75/25, 5.5/50, 8.25/75, 11/100}
  \draw[popdark] (\x,0.12) -- (\x,-0.12) node[below, font=\small] {\p\,\%};
\node[above, align=center, font=\bfseries, text=popgreen] at (10.3,0.35) {must\\ (certain)};
\node[above, align=center, font=\bfseries, text=poporange] at (5.5,0.35) {could / may\\ (possible)};
\node[above, align=center, font=\bfseries, text=poppink] at (0.8,0.35) {can't\\ (impossible)};
\fill[popgreen] (10.3,0) circle (0.14);
\fill[poporange] (5.5,0) circle (0.14);
\fill[poppink] (0.8,0) circle (0.14);
\end{tikzpicture}
\end{center}
\legende{L'échelle de certitude : de l'impossibilité (\eng{can't}) à la certitude (\eng{must}), en passant par la possibilité (\eng{could/may}).}
\end{popfigure}

\courssub{La démarche du détective : observer, évaluer, conclure}

Une déduction repose toujours sur un \cle{indice observable} dans la situation : un bruit, une lumière, une absence, un vêtement. Sans indice, pas de déduction valable ! Voici la démarche en trois étapes :

\begin{popfigure}
\begin{center}
\begin{tikzpicture}[node distance=0.5cm and 1.1cm]
\node[draw, rounded corners, fill=popblueL, align=center, font=\small\bfseries, text width=3.2cm] (obs) {1. Observer un indice\\ \emph{The lights are on.}};
\node[draw, rounded corners, fill=poporangeL, align=center, font=\small\bfseries, text width=3.2cm, right=of obs] (choix) {2. Évaluer ma certitude\\ \emph{How sure am I?}};
\node[draw, rounded corners, fill=popgreenL, align=center, font=\small\bfseries, text width=3.6cm, right=of choix] (concl) {3. Choisir le modal\\ \emph{must / could / can't}};
\draw[-{Stealth}, very thick, popdark] (obs) -- (choix);
\draw[-{Stealth}, very thick, popdark] (choix) -- (concl);
\node[draw, rounded corners, fill=popcream, align=center, font=\small, text width=9.5cm, below=0.8cm of choix] (ex) {\eng{Someone must be at home.} (sûr) --- \eng{Someone could be at home.} (possible) --- \eng{Nobody can be at home.} (impossible)};
\draw[-{Stealth}, thick, popmuted] (concl.south) -- (ex.east);
\end{tikzpicture}
\end{center}
\legende{Organigramme de la déduction : de l'indice au bon modal.}
\end{popfigure}

\courssub{Must d'obligation ou must de déduction ? Le contexte décide}

Le même mot \eng{must} a deux sens très différents. Comment les distinguer ? En regardant le \emph{contexte} et le \emph{sens de la phrase}.

\begin{center}
\begin{tabularx}{\linewidth}{|X|X|X|}\hline
\rowcolor{popblueL} \textbf{Phrase} & \textbf{Sens} & \textbf{Indice de reconnaissance} \\ \hline
\eng{You must do your homework.} & obligation (tu dois) & ordre, règle, devoir \\ \hline
\eng{You must be tired after the journey.} & déduction (tu dois être fatigué) & conclusion tirée d'un indice \\ \hline
\eng{Students must arrive at 7.30.} & obligation & règlement scolaire \\ \hline
\eng{It must be 7.30; the bell is ringing.} & déduction & indice : la cloche sonne \\ \hline
\end{tabularx}
\end{center}

\begin{aretenir}[L'astuce de traduction]
Pour savoir quel \eng{must} on a sous les yeux, traduisez en français : si « il est obligé de » fonctionne, c'est l'\emph{obligation} ; si « il est certain que / sûrement » fonctionne, c'est la \emph{déduction}. Notez aussi que la déduction porte souvent sur le \emph{présent} avec \eng{must be} + adjectif ou \eng{must be} + nom.
\end{aretenir}

\begin{attention}[Le piège n° 1 des francophones]
Pour la déduction négative, le français « il ne doit pas être là » se traduit par \eng{He can't be there.} — \textbf{jamais} \eng{He mustn't be there.}, qui signifierait « il ne doit pas être là » au sens d'une \emph{interdiction} (« il n'a pas le droit d'être là »).\par
Autre erreur fréquente : ne mettez jamais \eng{to} après un modal : \eng{She must \textbf{be} sad}, et non \eng{She must to be sad}. Et le modal ne prend jamais de \eng{-s} à la troisième personne : \eng{He must be}, et non \eng{He musts be}.
\end{attention}

\begin{exoresolu}[Déductions en famille]
\textbf{Énoncé.} Complétez chaque phrase avec \eng{must}, \eng{can't} ou \eng{could}, selon le degré de certitude indiqué par l'indice.\par
1. Amina is wearing her best dress and new shoes. She \ldots\ be going to a wedding. (sûr à 95 \%)\par
2. That \ldots\ be Paul's bicycle; his is blue and this one is black. (impossible)\par
3. Father is not answering his phone. He \ldots\ be in a meeting, or perhaps he is driving. (possible)\par
4. The whole family is speaking English fluently. They \ldots\ come from the North-West or South-West region. (sûr à 95 \%)\par
5. My sister \ldots\ be at school; the schools are closed for the holidays. (impossible)
\tcblower
\textbf{Solution détaillée.}\par
1. \eng{She \textbf{must} be going to a wedding.} — L'indice (belle robe, chaussures neuves) rend la conclusion quasi certaine : déduction affirmative avec \eng{must}.\par
2. \eng{That \textbf{can't} be Paul's bicycle.} — L'indice (la couleur différente) rend l'hypothèse impossible : déduction négative avec \eng{can't} (jamais \eng{mustn't}).\par
3. \eng{He \textbf{could} be in a meeting.} — Le locuteur hésite entre deux hypothèses : possibilité à 50 \%, donc \eng{could} (ou \eng{may}).\par
4. \eng{They \textbf{must} come from the North-West or South-West region.} — Parler couramment anglais est un indice fort : déduction affirmative avec \eng{must}.\par
5. \eng{My sister \textbf{can't} be at school.} — Les écoles sont fermées : impossibilité logique, donc \eng{can't}.
\end{exoresolu}

\begin{exercice}[À vous de déduire]
Observez chaque situation et écrivez une déduction avec \eng{must}, \eng{can't} ou \eng{could}.\par
1. Your brother is sleeping and snoring loudly. (He / be / very tired)\par
2. Your neighbour's daughter is carrying a stethoscope. (She / be / a nurse or a doctor)\par
3. It is 2 a.m. and the phone is ringing. (It / be / good news)\par
4. Your friend is limping after the football match. (His leg / hurt)\par
5. You don't know where your uncle is; he mentioned Douala and Bafoussam. (He / be / in Douala)
\end{exercice}

\begin{corrige}[Exercice « À vous de déduire »]
1. \eng{He must be very tired.} (certitude : il ronfle, il dort profondément)\par
2. \eng{She could be a nurse or a doctor.} (possibilité : le stéthoscope est un indice, pas une preuve absolue)\par
3. \eng{It can't be good news.} (impossibilité logique : on n'annonce pas de bonnes nouvelles à 2 h du matin)\par
4. \eng{His leg must hurt.} (certitude : il boite)\par
5. \eng{He could be in Douala.} (possibilité : deux hypothèses restent ouvertes)
\end{corrige}

\begin{savaistu}[Must et can't au Bac]
Au baccalauréat, la déduction avec \eng{must} et \eng{can't} apparaît souvent dans les questions de compréhension de type \eng{What can we deduce from the text?} ou \eng{What does the speaker imply?}. Savoir repérer ces modaux dans un texte aide à interpréter l'attitude du locuteur : \eng{must} signale qu'il est convaincu, \eng{can't} qu'il rejette une hypothèse, \eng{could/may} qu'il hésite. Dans les épreuves de grammaire, on vous demandera souvent de choisir entre \eng{mustn't} et \eng{can't} : rappelez-vous que seul \eng{can't} exprime l'impossibilité logique !
\end{savaistu}

\coursec{Comparaison français/anglais, synthèse et entraînement}

Nous savons maintenant exprimer la permission, la possibilité, le conseil, l'obligation et même la déduction logique avec \eng{must} et \eng{can't}. Il reste une dernière étape, peut-être la plus importante pour un élève francophone : comparer systématiquement l'anglais et le français, rassembler toutes les règles dans une synthèse claire, puis s'entraîner jusqu'à ce que le choix du bon modal devienne un réflexe. C'est l'objectif de cette dernière partie.

\courssub{Les équivalents français des modaux anglais}

\courssub{Observer avant de traduire}

Lisons d'abord ce court dialogue entre une mère et son fils à Yaoundé, et essayons de repérer à quel verbe français correspond chaque modal anglais.

\begin{textebox}[At home in Yaoundé — a dialogue]
\textbf{Mother:} Junior, you must finish your homework before you watch television.\\
\textbf{Junior:} Must I do it now? Can't I rest a little first?\\
\textbf{Mother:} No, you can't. And you should help your sister with her English afterwards.\\
\textbf{Junior:} May I go out with my friends after that?\\
\textbf{Mother:} Yes, you may, but you ought to be back before seven o'clock. You needn't cook tonight — I will do it myself.\\
\textbf{Junior:} Thank you, Mum! You must be the best mother in Yaoundé!\\
\textbf{Mother:} And you can't be serious — you just want to go out!
\end{textebox}

\begin{propriete}[Équivalences français/anglais des modaux]
\begin{itemize}
\item \emph{je peux} (permission) = \cle{can} / \cle{may} : \eng{May I come in?} « Puis-je entrer ? »
\item \emph{tu devrais} (conseil) = \cle{should} / \cle{ought to} : \eng{You should rest.} « Tu devrais te reposer. »
\item \emph{tu dois} (obligation) = \cle{must} / \cle{have to} : \eng{You must obey your parents.} « Tu dois obéir à tes parents. »
\item \emph{il est interdit de} = \cle{mustn't} : \eng{You mustn't lie.} « Il est interdit de mentir. »
\item \emph{tu n'es pas obligé de} = \cle{don't have to} / \cle{needn't} : \eng{You needn't come.} « Tu n'es pas obligé de venir. »
\item \emph{il doit être} (certitude, déduction) = \cle{must be} : \eng{He must be tired.} « Il doit être fatigué. »
\item \emph{il ne peut pas être} (déduction négative) = \cle{can't be} : \eng{She can't be forty!} « Elle ne peut pas avoir quarante ans ! »
\end{itemize}
\end{propriete}

\begin{center}
\begin{tabularx}{\linewidth}{|X|X|X|}
\hline
\rowcolor{popblueL} Français & English & Remarque \\
\hline
pouvoir (permission) & can / may & \eng{may} est plus poli et plus formel \\
\hline
pouvoir (possibilité) & can / could & \eng{could} = possibilité moins sûre \\
\hline
devoir (conseil) : tu devrais & should / ought to & conditionnel français = modal anglais \\
\hline
devoir (obligation) : tu dois & must / have to & \eng{have to} = obligation extérieure \\
\hline
falloir : il faut & must / have to / it is necessary to & « il faut » n'a pas de sujet personnel \\
\hline
il ne faut pas / interdit & mustn't & interdiction forte \\
\hline
ne pas être obligé de & don't have to / needn't & pas d'interdiction, simple inutilité \\
\hline
devoir (certitude) : il doit être & must be & déduction logique \\
\hline
il ne peut pas être (certitude) & can't be & déduction négative \\
\hline
\end{tabularx}
\end{center}

\begin{attention}[Le piège du conditionnel français]
En français, le conseil s'exprime au conditionnel : « tu \emph{devrais} ». En anglais, on n'utilise jamais \eng{would} pour le conseil : on dit \eng{You should work harder.} « Tu devrais travailler davantage. » De même, « il \emph{faudrait} » se traduit par \eng{we should} ou \eng{we ought to}, jamais par un conditionnel.
\end{attention}

\courssub{Synthèse : les quatre propriétés communes à tous les modaux}

\begin{propriete}[Les quatre règles d'or des auxiliaires modaux]
\begin{enumerate}
\item \cle{Pas de -s} à la troisième personne du singulier : \eng{He can swim.} (jamais \emph{he cans}).
\item \cle{Pas de to} devant l'infinitif qui suit : \eng{You must go.} (jamais \emph{must to go}). \textbf{Seule exception :} \cle{ought to}.
\item \cle{Pas de do} pour la négation et l'interrogation : \eng{Can you help me?} / \eng{She must not cheat.} (jamais \emph{Do you can?}).
\item \cle{Pas de futur propre} : les modaux n'ont qu'une forme ; pour exprimer le futur, on utilise des substituts (voir ci-dessous).
\end{enumerate}
Structure générale : \textbf{Sujet + modal (+ not) + base verbale} ; interrogation : \textbf{Modal + sujet + base verbale ?}
\end{propriete}

\begin{exemplebox}[Exemples traduits]
\eng{You should help your mother.} « Tu devrais aider ta mère. »\\
\eng{Must I sweep the yard? No, you needn't.} « Dois-je balayer la cour ? Non, ce n'est pas la peine. »\\
\eng{He must be tired after the journey from Douala.} « Il doit être fatigué après le voyage depuis Douala. »\\
\eng{She ought not to speak to her grandmother like that.} « Elle ne devrait pas parler ainsi à sa grand-mère. »
\end{exemplebox}

\courssub{Les substituts pour les temps manquants (complément Bac)}

Les modaux n'ayant qu'une forme, l'anglais utilise des verbes de remplacement pour exprimer le passé, le futur ou d'autres temps. Cette notion est régulièrement exigée au baccalauréat.

\begin{center}
\begin{tabularx}{\linewidth}{|X|X|X|}
\hline
\rowcolor{popblueL} Modal & Substitut & Exemple au futur / au passé \\
\hline
can (capacité) & be able to & \eng{I will be able to drive next year.} « Je saurai conduire l'an prochain. » \\
\hline
must (obligation) & have to & \eng{We had to wear uniforms last year.} « Nous devions porter l'uniforme l'an dernier. » \\
\hline
may (permission) & be allowed to & \eng{The students were allowed to leave early.} « Les élèves ont été autorisés à partir tôt. » \\
\hline
\end{tabularx}
\end{center}

\begin{exemplebox}[Exemples traduits]
\eng{When I am eighteen, I will be able to vote.} « Quand j'aurai dix-huit ans, je pourrai voter. »\\
\eng{My father had to work on Sundays when he was young.} « Mon père devait travailler le dimanche quand il était jeune. »\\
\eng{Will we be allowed to use our phones during the break?} « Serons-nous autorisés à utiliser nos téléphones pendant la récréation ? »
\end{exemplebox}

\courssub{Erreurs fréquentes des francophones : récapitulatif corrigé}

\begin{attention}[Les six pièges à connaître par cœur]
\begin{enumerate}
\item Jamais de \cle{to} après un modal, sauf \eng{ought to} : \emph{I must to go} ✗ → \eng{I must go.} ✓
\item Jamais de \cle{-s} : \emph{He cans swim} ✗ → \eng{He can swim.} ✓
\item Jamais de \cle{do} : \emph{Do you can dance?} ✗ → \eng{Can you dance?} ✓
\item \eng{mustn't} ≠ « pas obligé » : \eng{You mustn't smoke.} « Il est interdit de fumer » ; \eng{You don't have to smoke.} « Tu n'es pas obligé de fumer. »
\item Déduction négative = \cle{can't}, jamais \emph{mustn't} : \eng{He can't be ill; I saw him playing football.} « Il ne peut pas être malade. »
\item \eng{can} n'a pas de futur : \emph{I will can} ✗ → \eng{I will be able to.} ✓
\end{enumerate}
\end{attention}

\begin{center}
\begin{tabularx}{\linewidth}{|X|X|}
\hline
\rowcolor{popblueL} Phrase fautive de l'élève francophone & Correction \\
\hline
\emph{I must to clean my room.} & \eng{I must clean my room.} \\
\hline
\emph{She cans speak three languages.} & \eng{She can speak three languages.} \\
\hline
\emph{Do you can help me?} & \eng{Can you help me?} \\
\hline
\emph{You mustn't come if you are busy.} (= pas obligé) & \eng{You needn't come if you are busy.} \\
\hline
\emph{He ought do his duty.} & \eng{He ought to do his duty.} \\
\hline
\emph{Tomorrow I will can visit my uncle.} & \eng{Tomorrow I will be able to visit my uncle.} \\
\hline
\end{tabularx}
\end{center}

\courssub{Tableau récapitulatif final et carte mentale}

\begin{popfigure}
\begin{center}
\begin{tikzpicture}[scale=0.92, transform shape]
\node[draw=popblue, fill=popblueL, rounded corners, font=\bfseries\small, align=center] (f1) at (0,3.4) {Permission\\ \eng{can / may}\\ \eng{You may go out.}};
\node[draw=popgreen, fill=popgreenL, rounded corners, font=\bfseries\small, align=center] (f2) at (5.2,3.4) {Possibilité\\ \eng{can / could}\\ \eng{It could rain.}};
\node[draw=poporange, fill=poporangeL, rounded corners, font=\bfseries\small, align=center] (f3) at (10.4,3.4) {Conseil\\ \eng{should / ought to}\\ \eng{You should rest.}};
\node[draw=poppurple, fill=poppurpleL, rounded corners, font=\bfseries\small, align=center] (f4) at (0,0) {Obligation\\ \eng{must / have to}\\ \eng{I must obey.}};
\node[draw=poppink, fill=poppinkL, rounded corners, font=\bfseries\small, align=center] (f5) at (5.2,0) {Interdiction\\ \eng{mustn't}\\ \eng{You mustn't cheat.}};
\node[draw=popgold, fill=popgoldL, rounded corners, font=\bfseries\small, align=center] (f6) at (10.4,0) {Absence d'obligation\\ \eng{needn't / don't have to}\\ \eng{You needn't wait.}};
\node[draw=popdark, fill=popcream, rounded corners, font=\bfseries\small, align=center] (f7) at (5.2,-3.2) {Déduction\\ \eng{must be / can't be}\\ \eng{She must be tired.}};
\end{tikzpicture}
\end{center}
\legende{Tableau récapitulatif : fonction, modal et exemple.}
\end{popfigure}

\begin{popfigure}
\begin{center}
\begin{tikzpicture}[mindmap, grow cyclic, every node/.style={concept, font=\small}, concept color=popblue!40, level 1/.append style={level distance=3.6cm, sibling angle=60, concept color=poporange!40}, level 2/.append style={level distance=2.6cm, concept color=popgreen!40, font=\scriptsize}]
\node[font=\bfseries] {MODALS}
  child { node {Permission} child { node[align=center] {can / may\\ not: can't} } }
  child { node {Possibility} child { node[align=center] {can / could\\ not: can't} } }
  child { node {Advice} child { node[align=center] {should / ought to\\ not: shouldn't} } }
  child { node {Obligation} child { node[align=center] {must / have to\\ not: mustn't = interdit} } }
  child { node {No obligation} child { node[align=center] {needn't /\\ don't have to} } }
  child { node {Deduction} child { node[align=center] {must be /\\ can't be} } };
\end{tikzpicture}
\end{center}
\legende{Carte mentale de synthèse : les six fonctions des modaux et leurs formes négatives.}
\end{popfigure}

\courssub{Méthode : réussir les exercices de transformation}

\begin{methode}[Transformer une phrase avec un modal en quatre étapes]
\begin{enumerate}
\item \textbf{Repérer la fonction demandée} dans la consigne : permission, conseil, obligation, interdiction, absence d'obligation ou déduction.
\item \textbf{Choisir le modal} correspondant : \eng{may} pour la permission polie, \eng{should} pour le conseil, \eng{must} pour l'obligation, \eng{mustn't} pour l'interdiction, \eng{needn't} pour l'absence d'obligation, \eng{must be} pour la certitude.
\item \textbf{Appliquer la structure} : Sujet + modal (+ \eng{not}) + base verbale ; à l'interrogatif, placer le modal devant le sujet.
\item \textbf{Vérifier} : pas de \eng{-s}, pas de \eng{to} (sauf \eng{ought to}), pas de \eng{do}, et le sens correspond bien à la fonction demandée.
\end{enumerate}
\end{methode}

\begin{exoresolu}[Transformation guidée]
Énoncé : réécris la phrase \eng{It is forbidden to make noise in the library} en utilisant le modal \eng{mustn't}, puis mets la phrase \eng{You should visit your grandparents} à la forme interrogative.
\tcblower
Solution détaillée :
\begin{enumerate}
\item Fonction : interdiction → modal \eng{mustn't}. Structure : Sujet + mustn't + base verbale. On obtient : \eng{You mustn't make noise in the library.} « Il est interdit de faire du bruit à la bibliothèque. » Vérification : pas de \eng{to}, pas de \eng{-s}.
\item Forme interrogative : on place le modal devant le sujet, sans \eng{do} : \eng{Should you visit your grandparents?} « Devrais-tu rendre visite à tes grands-parents ? » (ou, avec \eng{I} : \eng{Should I visit my grandparents?}).
\end{enumerate}
\end{exoresolu}

\courssub{Entraînement}

\begin{exercice}[Exercice 1 — Phrases à trous]
Complète avec \eng{can}, \eng{may}, \eng{must}, \eng{mustn't}, \eng{should}, \eng{ought to}, \eng{needn't} ou \eng{can't}.
\begin{enumerate}
\item You \ldots\ldots cross the road without looking; it is dangerous.
\item \ldots\ldots I borrow your pen, please?
\item You look exhausted; you \ldots\ldots go to bed earlier.
\item That boy \ldots\ldots be Peter's brother: they look exactly the same!
\item You \ldots\ldots bring any food; the school provides lunch.
\item We \ldots\ldots to respect our elders. (utilise \eng{ought})
\end{enumerate}
\end{exercice}

\begin{exercice}[Exercice 2 — Choix du modal selon la fonction]
Pour chaque situation, choisis le modal correct entre parenthèses.
\begin{enumerate}
\item Permission polie : \eng{\ldots\ldots I use your phone?} (\eng{must} / \eng{may})
\item Interdiction : \eng{You \ldots\ldots play with matches.} (\eng{mustn't} / \eng{needn't})
\item Absence d'obligation : \eng{You \ldots\ldots wash the dishes; I will do it.} (\eng{mustn't} / \eng{don't have to})
\item Déduction négative : \eng{She \ldots\ldots be at home; her bicycle is not there.} (\eng{mustn't} / \eng{can't})
\item Conseil : \eng{You \ldots\ldots apologise to your sister.} (\eng{should} / \eng{can})
\end{enumerate}
\end{exercice}

\begin{exercice}[Exercice 3 — Correction d'erreurs]
Chaque phrase contient une erreur. Corrige-la.
\begin{enumerate}
\item \emph{He cans play the guitar very well.}
\item \emph{I must to visit my aunt in Bafoussam.}
\item \emph{Do you can speak English?}
\item \emph{She ought apologise to her friend.}
\item \emph{You mustn't come tomorrow if you are tired.} (sens voulu : pas obligé)
\item \emph{Next year, I will can travel alone.}
\end{enumerate}
\end{exercice}

\begin{exercice}[Exercice 4 — Transformations]
\begin{enumerate}
\item Mets à la forme négative : \eng{You should tell lies.}
\item Mets à la forme interrogative : \eng{She can swim.}
\item Reformule avec \eng{must} : \eng{It is necessary to respect the school rules.}
\item Reformule avec \eng{needn't} : \eng{It is not necessary for you to bring your dictionary.}
\item Reformule avec \eng{can't} (déduction) : \eng{I am sure that man is not a teacher; he is too young.}
\end{enumerate}
\end{exercice}

\begin{corrige}[Exercices 1 à 4]
\textbf{Exercice 1 :} 1. \eng{mustn't} (interdiction, danger) ; 2. \eng{May} (permission polie) ; 3. \eng{should} (conseil) ; 4. \eng{must} (déduction positive) ; 5. \eng{needn't} (absence d'obligation) ; 6. \eng{ought} (\eng{We ought to respect our elders.}).

\textbf{Exercice 2 :} 1. \eng{may} ; 2. \eng{mustn't} ; 3. \eng{don't have to} ; 4. \eng{can't} ; 5. \eng{should}.

\textbf{Exercice 3 :} 1. \eng{He can play the guitar very well.} (pas de \eng{-s}) ; 2. \eng{I must visit my aunt in Bafoussam.} (pas de \eng{to}) ; 3. \eng{Can you speak English?} (pas de \eng{do}) ; 4. \eng{She ought to apologise to her friend.} (\eng{to} obligatoire après \eng{ought}) ; 5. \eng{You needn't come tomorrow if you are tired.} (\eng{mustn't} = interdiction) ; 6. \eng{Next year, I will be able to travel alone.} (\eng{can} n'a pas de futur).

\textbf{Exercice 4 :} 1. \eng{You should not (shouldn't) tell lies.} ; 2. \eng{Can she swim?} ; 3. \eng{You must respect the school rules.} ; 4. \eng{You needn't bring your dictionary.} ; 5. \eng{That man can't be a teacher; he is too young.}
\end{corrige}

\courssub{Micro-production : les règles de vie dans ma famille}

\begin{experience}[Writing — « Rules in my family »]
Rédige un court paragraphe (six à huit phrases) intitulé \eng{Rules in my family} en utilisant \textbf{au moins cinq modaux différents} parmi : \eng{can}, \eng{may}, \eng{must}, \eng{mustn't}, \eng{should}, \eng{ought to}, \eng{needn't}, \eng{don't have to}, \eng{can't}. Souligne chaque modal utilisé. Exemple de début : \eng{In my family, we must respect one another. We mustn't shout at home, but we can speak freely about our problems\ldots}
\end{experience}

\begin{exemplebox}[Modèle de paragraphe]
\eng{In my family, we \underline{must} respect one another. We \underline{mustn't} come home late, because our parents worry about us. We \underline{should} always tell the truth, and we \underline{ought to} help with the housework. On Saturdays, we \underline{don't have to} wake up early, so we \underline{can} rest a little. My little brother \underline{needn't} wash the dishes yet, because he is too young. When my father comes home smiling, we know he \underline{must} be in a good mood!} « Dans ma famille, nous devons nous respecter les uns les autres. Il est interdit de rentrer tard, car nos parents s'inquiètent. Nous devrions toujours dire la vérité et nous devrions aider aux tâches ménagères. Le samedi, nous ne sommes pas obligés de nous lever tôt, donc nous pouvons nous reposer un peu. Mon petit frère n'est pas encore obligé de laver la vaisselle, car il est trop jeune. Quand mon père rentre en souriant, nous savons qu'il doit être de bonne humeur ! »
\end{exemplebox}

\begin{savaistu}[Les modaux au baccalauréat A]
Dans les épreuves d'anglais du Bac A au Cameroun, les auxiliaires modaux sont fréquemment testés dans la section « Grammar » par des items à choix multiples. Les deux confusions préférées des concepteurs d'épreuves sont précisément \eng{mustn't} contre \eng{don't have to} (interdiction contre absence d'obligation) et \eng{should} contre \eng{ought to} (avec le piège du \eng{to} oublié ou ajouté au mauvais endroit). Un entraînement régulier sur ces deux points peut rapporter des points précieux le jour de l'examen.
\end{savaistu}

\begin{aretenir}[L'essentiel de la leçon]
\begin{itemize}
\item Chaque fonction a son modal : permission \eng{can/may}, possibilité \eng{can/could}, conseil \eng{should/ought to}, obligation \eng{must/have to}, interdiction \eng{mustn't}, absence d'obligation \eng{needn't/don't have to}, déduction \eng{must be/can't be}.
\item Quatre règles d'or : pas de \eng{-s}, pas de \eng{to} (sauf \eng{ought to}), pas de \eng{do}, pas de futur propre.
\item Pour les temps manquants : \eng{be able to} (can), \eng{have to} (must), \eng{be allowed to} (may).
\item \eng{mustn't} signifie « interdit », jamais « pas obligé » ; la déduction négative se dit \eng{can't}.
\end{itemize}
\end{aretenir}

\newpage

\coursec{Activité d'intégration}

\courssub{Objectifs de l'activité}

À l'issue de cette activité d'intégration, l'élève sera capable de :

\begin{itemize}
\item \cle{mobiliser} l'ensemble des auxiliaires modaux étudiés dans la leçon (\eng{can, may, could, should, ought to, must, have to, mustn't, needn't, don't have to}) dans une production écrite authentique ;
\item \cle{distinguer} les nuances de sens : permission, conseil, obligation, interdiction, absence d'obligation et déduction logique ;
\item \cle{produire} un texte collectif court, structuré et correct (une charte de vie familiale) ;
\item \cle{s'exprimer oralement} devant la classe et poser des questions avec les modaux (\eng{Can we...? Do we have to...? Should we...?}) ;
\item \cle{évaluer} la production des autres groupes selon des critères de correction grammaticale et de réalisme.
\end{itemize}

\courssub{Situation}

\textbf{Contexte.} La famille Nkeng habite à Buea, dans la région du Sud-Ouest. Le père est enseignant, la mère tient un petit commerce au marché central, et ils ont quatre enfants : deux au lycée et deux au primaire. Après plusieurs disputes à propos des tâches ménagères, de l'argent de poche et de l'utilisation du téléphone, les parents décident d'organiser un \emph{family meeting} (conseil de famille) pour rédiger ensemble une \emph{Family Rules Charter}, une charte de vie que chacun devra respecter.

Voici l'annonce que le père, Mr Nkeng, a écrite au tableau de la maison :

\begin{textebox}[Mr Nkeng's announcement]
Dear family,

This Saturday at 5 p.m., we shall hold our first family meeting. We must talk about our life together: housework, pocket money, phones, visitors and respect. Everybody can give ideas, and nobody has to stay silent! Together, we shall write our Family Rules Charter: ten to twelve rules that we must all follow. You should think about your ideas before Saturday. You needn't write a long speech — just bring your suggestions.

See you on Saturday!

Dad
\end{textebox}

\begin{definition}[Glossaire]
\begin{itemize}
\item \eng{housework} (nom) : les tâches ménagères ;
\item \eng{pocket money} (nom) : l'argent de poche ;
\item \eng{a rule} (nom) : une règle ;
\item \eng{to follow a rule} (verbe) : respecter une règle ;
\item \eng{a suggestion} (nom) : une suggestion, une proposition.
\end{itemize}
\end{definition}

\textbf{Votre mission.} Par groupes de quatre, vous êtes les membres de la famille Nkeng (ou de votre propre famille, si vous préférez). Rédigez la charte de vie de la famille en 10 à 12 règles, puis présentez-la à la classe sous forme d'affiche.

\courssub{Consignes}

\textbf{Tâche 1 — Compréhension du document (5 min).} Lisez l'annonce de Mr Nkeng et répondez aux questions :
\begin{enumerate}
\item When is the family meeting?
\item What topics must the family talk about? Name at least three.
\item Find two modal verbs in the text and say what they express (obligation, permission, advice or absence of obligation).
\end{enumerate}

\textbf{Tâche 2 — Remue-méninges lexical (10 min).} En groupe, listez en anglais au moins huit domaines de la vie familiale à régler. Utilisez le lexique suivant et ajoutez vos propres idées : \eng{housework, meals, pocket money, mobile phones, homework, visitors, bedtime, respect, noise, TV, the family shop, church or mosque}.

\textbf{Tâche 3 — Rédaction de la charte (25 min).} Rédigez votre \emph{Family Rules Charter} en 10 à 12 règles numérotées. Votre charte doit contenir \textbf{obligatoirement} :
\begin{enumerate}
\item \textbf{deux permissions} avec \eng{can} ou \eng{may} ;
\item \textbf{deux conseils} avec \eng{should} ou \eng{ought to} ;
\item \textbf{deux obligations} avec \eng{must} ou \eng{have to} ;
\item \textbf{une interdiction} avec \eng{mustn't} ;
\item \textbf{deux absences d'obligation} avec \eng{needn't} ou \eng{don't have to} ;
\item \textbf{une déduction logique} avec \eng{must} ou \eng{can't} (par exemple : \eng{If the door is locked, Dad must be at the shop.}).
\end{enumerate}
Soulignez chaque modal et indiquez entre parenthèses sa fonction : \eng{(permission)}, \eng{(advice)}, \eng{(obligation)}, \eng{(prohibition)}, \eng{(no obligation)}, \eng{(deduction)}.

\textbf{Tâche 4 — Relecture collective (10 min).} Échangez votre brouillon avec un autre groupe. Vérifiez :
\begin{enumerate}
\item la base verbale après chaque modal (jamais de \eng{to} après \eng{must} ou \eng{can}) ;
\item la forme négative correcte (\eng{mustn't}, \eng{needn't}, \eng{don't have to}) ;
\item la présence des huit catégories de modaux exigées.
\end{enumerate}

\textbf{Tâche 5 — Présentation orale et débat (20 min).} Chaque groupe présente son affiche à la classe. Les autres groupes posent au moins trois questions avec des modaux : \eng{Can we invite friends on Sundays? Do we have to wake up at 5 a.m.? Should children cook alone? Why mustn't we use phones at dinner?}

\textbf{Tâche 6 — Vote final (5 min).} La classe vote pour la charte la plus réaliste \emph{et} la mieux construite grammaticalement. Justifiez votre vote en une phrase : \eng{I vote for Group 3 because their charter is realistic and all the modals are correct.}

\courssub{Ressources à mobiliser}

\begin{aretenir}[Rappel des structures]
\begin{center}
\begin{tabularx}{\linewidth}{|l|X|X|}
\hline
\rowcolor{popblueL} \textbf{Fonction} & \textbf{Modaux} & \textbf{Exemple} \\
\hline
Permission & \eng{can / may} & \eng{Children can watch TV after homework.} \\
\hline
Conseil & \eng{should / ought to} & \eng{We ought to respect our parents.} \\
\hline
Obligation & \eng{must / have to} & \eng{Everybody must wash the dishes.} \\
\hline
Interdiction & \eng{mustn't} & \eng{You mustn't shout in the house.} \\
\hline
Absence d'obligation & \eng{needn't / don't have to} & \eng{You needn't cook on Sundays.} \\
\hline
Déduction & \eng{must / can't} & \eng{The lights are off — they must be asleep.} \\
\hline
\end{tabularx}
\end{center}
\end{aretenir}

\begin{attention}[Pièges à éviter]
\begin{itemize}
\item Après un modal, toujours la \cle{base verbale} : \eng{We must to help} est faux ; écrivez \eng{We must help.}
\item \eng{Ought to} garde son \eng{to} : \eng{We ought to listen.} Mais à la forme négative : \eng{We ought not to shout.}
\item \eng{Mustn't} (interdiction) ne veut pas dire \eng{don't have to} (absence d'obligation) : \eng{You mustn't smoke} = il est interdit de fumer ; \eng{You don't have to smoke} = tu n'es pas obligé de fumer.
\item Pas de \eng{-s} à la troisième personne : \eng{She must works} est faux ; écrivez \eng{She must work.}
\end{itemize}
\end{attention}

\textbf{Lexique utile pour la charte :} \eng{to tidy one's room} (ranger sa chambre), \eng{to do the dishes} (faire la vaisselle), \eng{to sweep} (balayer), \eng{to fetch water} (aller chercher de l'eau), \eng{to switch off} (éteindre), \eng{to share} (partager), \eng{bedtime} (l'heure du coucher), \eng{to be polite} (être poli), \eng{to help Mum at the market} (aider maman au marché).

\courssub{Proposition de production et corrigé commenté}

\begin{exoresolu}[Modèle de charte — The Nkeng Family Rules Charter]
\textbf{Énoncé.} Rédigez la charte de vie de la famille Nkeng en 10 à 12 règles respectant la répartition imposée des modaux.
\tcblower
\textbf{Production modèle :}

\eng{\textbf{The Nkeng Family Rules Charter}}

\eng{1. Everybody must greet the parents in the morning. (obligation)}

\eng{2. Children have to do their homework before watching TV. (obligation)}

\eng{3. You mustn't use your phone during meals. (prohibition)}

\eng{4. We should eat dinner together every evening. (advice)}

\eng{5. The older children ought to help the younger ones with their schoolwork. (advice)}

\eng{6. Children can play outside after finishing their housework. (permission)}

\eng{7. On Saturdays, we may invite our friends if we ask Mum first. (permission)}

\eng{8. You needn't wake up at 5 a.m. on Sundays — you can rest! (no obligation)}

\eng{9. The little ones don't have to cook; the big children will do it. (no obligation)}

\eng{10. Everybody must switch off the lights before bedtime. (obligation)}

\eng{11. If the shop keys are not on the table, Mum must be at the market. (deduction)}

\eng{12. We should always say “please” and “thank you”. (advice)}

\medskip
\textbf{Commentaire des choix de langue :}
\begin{itemize}
\item Règle 1 : \eng{must} exprime une obligation forte posée par la famille ; base verbale \eng{greet} sans \eng{to}.
\item Règle 2 : \eng{have to} marque une obligation extérieure (imposée par la scolarité) ; attention à la question correspondante : \eng{Do children have to...?}
\item Règle 3 : \eng{mustn't} = interdiction absolue, à ne pas confondre avec \eng{don't have to}.
\item Règles 4, 5 et 12 : \eng{should} et \eng{ought to} sont interchangeables pour le conseil ; \eng{ought to} garde toujours son \eng{to}.
\item Règles 6 et 7 : \eng{can} (permission courante) et \eng{may} (permission plus formelle, avec condition : \eng{if we ask Mum first}).
\item Règles 8 et 9 : \eng{needn't} et \eng{don't have to} expriment tous deux l'absence d'obligation ; varier les deux formes enrichit le texte.
\item Règle 11 : déduction logique avec \eng{must} — on conclut à partir d'un indice (les clés absentes). La déduction négative serait \eng{can't} : \eng{Mum can't be at home.}
\item Aucun modal ne prend de \eng{-s} à la troisième personne, et chaque modal est suivi directement de la base verbale.
\end{itemize}
\end{exoresolu}

\courssub{Bilan de l'activité}

Cette activité a permis de réinvestir dans une tâche authentique et collective l'ensemble des modaux étudiés dans la leçon. Vous avez :

\begin{itemize}
\item \cle{compris} un document authentique annonçant un conseil de famille et repéré les modaux en contexte ;
\item \cle{produit} un texte normatif réel (une charte) en choisissant chaque modal en fonction de son sens : permission (\eng{can/may}), conseil (\eng{should/ought to}), obligation (\eng{must/have to}), interdiction (\eng{mustn't}), absence d'obligation (\eng{needn't/don't have to}) et déduction (\eng{must/can't}) ;
\item \cle{interagi oralement} en présentant votre affiche et en posant des questions correctement formées avec les modaux ;
\item \cle{évalué} la production d'autrui selon des critères grammaticaux précis, ce qui renforce votre propre vigilance sur les erreurs fréquentes des francophones (base verbale sans \eng{to}, absence de \eng{-s}, distinction \eng{mustn't} / \eng{don't have to}).
\end{itemize}

\begin{methode}[Pour aller plus loin à la maison]
Rédigez individuellement une mini-charte de cinq règles pour votre propre famille ou pour votre dortoir au lycée, en utilisant cinq modaux différents. Relisez-vous avec la grille suivante : modal bien choisi pour le sens voulu ? base verbale correcte ? forme négative correcte ? Présentez votre texte à un camarade qui vous posera deux questions avec \eng{Do we have to...?} et \eng{Can we...?}
\end{methode}

\newpage

\coursec{Méthodes et exercices résolus}

\coursec{Lesson 2 — Grammar: Modals — Permission, Possibility, Advice, Obligation, Deduction}

\courssub{Observation : les modaux en action}

\begin{textebox}[Dialogue familial à Yaoundé — Extrait d'une conversation authentique]
“Good morning, Dad. \textbf{May} I go to the youth centre with Ngozi this afternoon? We \textbf{have to} finish our Geography project.”
“You \textbf{can} go, but you \textbf{must} be back before 6 p.m. It \textbf{could} rain heavily later, so take an umbrella.”
“Okay. \textbf{Should} I cook the rice before I leave?”
“No, you \textbf{needn't} cook today. Your aunt is coming over and she \textbf{will} prepare the meal. But you \textbf{mustn't} forget to sweep the veranda.”
“Don't worry, Dad. I \textbf{ought to} do it right now. Look, the sky is dark. It \textbf{must} be the rainy season starting early!”
\end{textebox}

\begin{exemplebox}[Identification des modaux dans le texte]
\begin{itemize}
\item \eng{May I go?} — Permission formelle (demande).
\item \eng{have to finish} — Obligation externe (semi-modal, conjugué).
\item \eng{can go} — Permission accordée.
\item \eng{must be back} — Obligation forte / règle parentale.
\item \eng{could rain} — Possibilité hypothétique (incertitude).
\item \eng{Should I cook?} — Demande de conseil.
\item \eng{needn't cook} — Absence d'obligation (dispense).
\item \eng{mustn't forget} — Interdiction formelle.
\item \eng{ought to do} — Conseil / devoir moral (synonyme de \eng{should}).
\item \eng{must be} — Déduction logique (certitude positive).
\end{itemize}
\end{exemplebox}

\begin{aretenir}[Ce qu'est un verbe modal]
Un \cle{verbe modal} (ou auxiliaire modal) est un verbe \emph{défectif} : il n'a ni infinitif (\eng{*to can}), ni participe (\eng{*canned}), ni \eng{-s} à la 3\textsuperscript{e} personne du singulier. Il est toujours suivi de la \cle{base verbale} (infinitif sans \eng{to}), sauf \eng{ought to}. Les principaux : \cle{can, could, may, might, must, shall, should, will, would, ought to}. Les \emph{semi-modaux} (\eng{have to, need to, be able to}) se conjuguent normalement.
\end{aretenir}

\courssub{Permission et possibilité : can, could, may}

\begin{propriete}[Formes : Permission et Possibilité]
\begin{center}
\begin{tabularx}{\linewidth}{|X|X|X|}
\hline
\rowcolor{popblueL} \textbf{Fonction} & \textbf{Modal} & \textbf{Nuance / Registre} \\ \hline
\multicolumn{3}{|c|}{\textbf{PERMISSION}} \\ \hline
Demander (informel) & \eng{Can I / we \ldots?} & Courant, entre amis, famille. \eng{Can I borrow your pen?} \\ \hline
Demander (politesse standard) & \eng{Could I / we \ldots?} & Plus poli, recommandé avec un adulte / inconnu. \eng{Could I leave early?} \\ \hline
Demander (formel / écrit) & \eng{May I \ldots?} & Très formel, respectueux. \eng{May I speak to the Principal?} \\ \hline
Accorder & \eng{Yes, you can / may.} & \eng{May} reste formel à l'oral. \\ \hline
Refuser & \eng{No, you can't / may not.} & \eng{May not} est formel ; \eng{can't} est standard. \\ \hline
\multicolumn{3}{|c|}{\textbf{POSSIBILITÉ}} \\ \hline
Possibilité générale (fait) & \eng{can} & “It \textbf{can} get very hot in Maroua in March.” (C'est possible / ça arrive.) \\ \hline
Possibilité hypothétique (présent/futur) & \eng{could / may / might} & “It \textbf{could} rain this evening.” (Peut-être, moins certain que \eng{can}.) \\ \hline
Possibilité passée (non réalisée) & \eng{could have + V-en} & “You \textbf{could have} won.” (Tu aurais pu gagner, mais tu ne l'as pas fait.) \\ \hline
\end{tabularx}
\end{center}
\end{propriete}

\begin{methode}[Choisir le bon modal pour la permission / possibilité]
\begin{enumerate}
\item \textbf{Identifie le registre} : Familial/Amis $\rightarrow$ \eng{Can} ; Politesse courante $\rightarrow$ \eng{Could} ; Formel/Administratif $\rightarrow$ \eng{May}.
\item \textbf{Distinction \eng{can} / \eng{could} (possibilité)} : \eng{can} = possibilité théorique, générale (“Les bendskins \textbf{can} be dangerous”). \eng{could} = hypothèse ponctuelle, incertaine (“This bendskin \textbf{could} break down”).
\item \textbf{Structure invariante} : Sujet + Modal + Base Verbale (sans \eng{to}). Négation : Modal + \eng{not} (\eng{can't, couldn't, may not}). Question : Inversion (Modal + Sujet + BV).
\end{enumerate}
\end{methode}

\begin{exoresolu}[Permission et possibilité en contexte camerounais]
\textbf{Énoncé.} Complete with the appropriate modal (\emph{can, could, may, can't, couldn't, may not}) and translate.\\
1. \textbf{(Formal request to a teacher)} \ldots\ I submit my assignment tomorrow, Sir?\\
2. \textbf{(Father to son)} You \ldots\ use the car tonight, but you \ldots\ drive fast. The roads are wet.\\
3. \textbf{(General truth)} In Buea, it \ldots\ rain at any time, even in dry season.\\
4. \textbf{(Hypothesis)} Look at those clouds! It \ldots\ rain before we reach the market.\\
5. \textbf{(Past missed opportunity)} You \ldots\ have met the Minister yesterday; he was at the ceremony.
\tcblower
\textbf{Solution détaillée.}
\begin{enumerate}
\item \eng{May I submit my assignment tomorrow, Sir?} — Demande formelle à un supérieur. « Puis-je rendre mon devoir demain, Monsieur ? »
\item \eng{You \textbf{can} use the car tonight, but you \textbf{can't} drive fast.} — Permission accordée (\eng{can}) + Interdiction (\eng{can't}). « Tu peux prendre la voiture ce soir, mais tu ne dois pas rouler vite. »
\item \eng{In Buea, it \textbf{can} rain at any time.} — Possibilité générale, fait connu. « À Buea, il peut pleuvoir à tout moment. »
\item \eng{Look at those clouds! It \textbf{could} rain before we reach the market.} — Hypothèse basée sur un indice visuel (nuages), incertitude. « Regarde ces nuages ! Il pourrait pleuvoir avant qu'on n'arrive au marché. »
\item \eng{You \textbf{could have} met the Minister yesterday.} — Possibilité passée non réalisée (regret / constat). « Tu aurais pu rencontrer le Ministre hier. »
\end{enumerate}
\end{exoresolu}

\courssub{Conseil : should et ought to}

\begin{propriete}[Formes : Should / Ought to]
\begin{center}
\begin{tabular}{|l|l|l|}
\hline
\rowcolor{popgreenL} \textbf{Forme} & \textbf{Should} & \textbf{Ought to} \\ \hline
Affirmative & \eng{You should rest.} & \eng{You ought to rest.} \\ \hline
Négative & \eng{You shouldn't rest.} (\eng{should not}) & \eng{You ought not to rest.} (formel) / \eng{You oughtn't to rest.} (rare) \\ \hline
Interrogative & \eng{Should I rest?} & \eng{Ought I to rest?} (très formel, rare à l'oral) \\ \hline
\multicolumn{3}{|l|}{\textbf{Traduction :} « Tu devrais te reposer. / Tu ne devrais pas te reposer. / Devrais-je me reposer ? »} \\ \hline
\end{tabular}
\end{center}
\end{propriete}

\begin{attention}[Pièges \eng{should} vs \eng{ought to}]
\begin{enumerate}
\item \textbf{Le \eng{to} obligatoire} : \eng{ought} \textbf{est toujours suivi de \eng{to}} + base verbale. \emph{Faux :} “You ought rest.” \textbf{Juste :} “You \textbf{ought to} rest.”
\item \textbf{Négation} : \eng{ought not to} (formel) ou \eng{oughtn't to} (britannique, rare). Préfère \eng{shouldn't} à l'oral et à l'écrit courant.
\item \textbf{Question} : \eng{Should I\ldots?} est la norme. \eng{Ought I to\ldots?} sonne vieilli / littéraire.
\item \textbf{Nuance} : \eng{ought to} insiste davantage sur le \emph{devoir moral}, l'obligation sociale ou la correction (ex. respect des aînés au Cameroun). \eng{should} est plus neutre, plus fréquent.
\end{enumerate}
\end{attention}

\begin{exoresolu}[Conseils pour la vie scolaire et familiale]
\textbf{Énoncé.} Your friend wants to succeed in the \eng{Probatoire} but spends nights on TikTok. Write 4 sentences of advice using \emph{should}, \emph{shouldn't}, \emph{ought to}, \emph{ought not to}.
\tcblower
\textbf{Solution détaillée.}
\begin{enumerate}
\item \eng{You \textbf{should} organise a revision timetable immediately.} (Conseil positif, neutre).
\item \eng{You \textbf{shouldn't} stay up late scrolling on your phone.} (Conseil négatif, contraction naturelle).
\item \eng{You \textbf{ought to} ask your teachers for past papers.} (Devoir moral / démarche recommandée formellement).
\item \eng{You \textbf{ought not to} neglect your health; sleep is essential for memory.} (Devoir moral fort, forme complète \eng{ought not to} pour l'insistance).
\end{enumerate}
\textbf{Traduction :} 1. Tu devrais organiser un planning de révision immédiatement. 2. Tu ne devrais pas veiller tard à scroller sur ton téléphone. 3. Tu devrais (il convient que tu) demandes des annales à tes professeurs. 4. Tu ne devrais pas négliger ta santé ; le sommeil est essentiel pour la mémoire.
\end{exoresolu}

\courssub{Obligation, interdiction, absence d'obligation : must, mustn't, needn't, don't have to}

\begin{propriete}[Le triangle de l'obligation : Distinctions capitales]
\begin{center}
\begin{tabularx}{\linewidth}{|X|X|X|X|}
\hline
\rowcolor{poporangeL} \textbf{Notion} & \textbf{Modal / Structure} & \textbf{Sens français} & \textbf{Exemple (Contexte Cameroun)} \\ \hline
\textbf{Obligation forte} (règle, loi, opinion du locuteur) & \eng{MUST} + BV & Il faut / Je dois (impératif) & \eng{Students \textbf{must} wear uniforms in public schools.} \\ \hline
\textbf{Interdiction} & \eng{MUSTN'T} + BV & Il est interdit de / Ne pas … & \eng{You \textbf{mustn't} use phones in class.} \\ \hline
\textbf{Absence d'obligation} (dispense, choix) & \eng{NEEDN'T} + BV \newline \eng{DON'T HAVE TO} + BV & Pas besoin de / Pas obligé de & \eng{We \textbf{needn't} pay fees in public primary schools.} \newline \eng{You \textbf{don't have to} come if you are sick.} \\ \hline
\textbf{Obligation externe / Passé / Futur} & \eng{HAVE TO} (conjugué) & Devoir (circonstances) & \eng{I \textbf{had to} walk to school yesterday.} \newline \eng{You \textbf{will have to} work hard.} \\ \hline
\end{tabularx}
\end{center}
\end{propriete}

\begin{attention}[Les 3 confusions fatales au Bac]
\begin{enumerate}
\item \textbf{\eng{MUSTN'T} $\neq$ \eng{DON'T HAVE TO}}.
\begin{itemize}
\item \eng{You mustn't cross the red light.} = \textbf{Interdiction} (Danger, loi). « Tu ne dois pas / Il est interdit de… »
\item \eng{You don't have to cross the red light.} = \textbf{Absence d'obligation} (Tu as le choix, ce n'est pas nécessaire). « Tu n'es pas obligé de… »
\end{itemize}
\item \textbf{\eng{NEEDN'T} vs \eng{DON'T NEED TO}}.
\begin{itemize}
\item \eng{Needn't} = Modal (pas de \eng{to}, pas de \eng{do}, 3\textsuperscript{e} pers. sans \eng{s}). “He \textbf{needn't} come.”
\item \eng{Don't need to} = Verbe plein (auxiliaire \eng{do}, \eng{to} + BV). “He \textbf{doesn't need to} come.”
\item \textbf{Au passé} : Seul \eng{didn't need to} / \eng{didn't have to} est possible (\eng{needn't} n'a pas de passé).
\end{itemize}
\item \textbf{\eng{MUST} n'a pas de passé ni de futur}. On le remplace par \eng{HAD TO} (passé) et \eng{WILL HAVE TO} (futur). \emph{Jamais} \eng{musted} ou \eng{will must}.
\end{enumerate}
\end{attention}

\begin{exoresolu}[Obligation, interdiction ou dispense ?]
\textbf{Énoncé.} Choose the correct form and justify the choice.\\
1. Candidates \ldots\ (must / mustn't / don't have to) bring their ID cards to the exam hall.\\
2. Since the teacher is absent, we \ldots\ (must / mustn't / don't have to) do the homework for tomorrow.\\
3. You \ldots\ (mustn't / don't have to) shout at your elders in Bamenda culture.\\
4. I \ldots\ (must / had to / will have to) renew my passport last year.
\tcblower
\textbf{Solution détaillée.}
\begin{enumerate}
\item \eng{Candidates \textbf{must} bring their ID cards.} — Obligation stricte (règle d'examen). « Les candidats doivent apporter leur carte d'identité. »
\item \eng{Since the teacher is absent, we \textbf{don't have to} do the homework for tomorrow.} — Absence d'obligation (dispense due aux circonstances). \eng{Mustn't} signifierait « interdiction de le faire », ce qui est illogique. « Comme le prof est absent, nous n'avons pas à faire le devoir pour demain. »
\item \eng{You \textbf{mustn't} shout at your elders in Bamenda culture.} — Interdiction culturelle forte, tabou. « Il est interdit / Tu ne dois pas crier sur tes aînés… »
\item \eng{I \textbf{had to} renew my passport last year.} — Passé : \eng{must} $\rightarrow$ \eng{had to}. « J'ai dû renouveler mon passeport l'an dernier. »
\end{enumerate}
\end{exoresolu}

\courssub{Déduction : must et can't logiques}

\begin{propriete}[Déduction au présent et au passé]
\begin{center}
\begin{tabularx}{\linewidth}{|X|X|X|}
\hline
\rowcolor{poppurpleL} \textbf{Degré de certitude} & \textbf{Structure} & \textbf{Exemple + Traduction} \\ \hline
\textbf{Certitude positive} (Indice visible $\rightarrow$ Conclusion sûre) & \eng{MUST} + \textbf{Base Verbale} (Présent) & \eng{His bag is here. He \textbf{must be} in the library.} \\
& & (« Son sac est là. Il \textbf{doit être} à la bibliothèque. ») \\ \hline
& \eng{MUST} + \textbf{HAVE + V-en} (Passé) & \eng{The gate is open. The watchman \textbf{must have forgotten} to lock it.} \\
& & (« Le portail est ouvert. Le gardien \textbf{a dû oublier} de le fermer. ») \\ \hline
\textbf{Certitude négative} (Indice $\rightarrow$ Impossibilité logique) & \eng{CAN'T} + \textbf{Base Verbale} (Présent) & \eng{She \textbf{can't be} 18; she was born in 2010.} \\
& & (« Elle \textbf{ne peut pas avoir} 18 ans ; elle est née en 2010. ») \\ \hline
& \eng{CAN'T} + \textbf{HAVE + V-en} (Passé) & \eng{He \textbf{can't have taken} the taxi; he has no money.} \\
& & (« Il \textbf{n'a pas pu prendre} le taxi ; il n'a pas d'argent. ») \\ \hline
\end{tabularx}
\end{center}
\end{propriete}

\begin{methode}[Réussir un exercice de déduction]
\begin{enumerate}
\item \textbf{Repère l'indice} (fait observable : lumière allumée, uniforme, note, absence, météo).
\item \textbf{Détermine la polarité} : L'indice \emph{confirme} l'hypothèse $\rightarrow$ \eng{MUST}. L'indice la \emph{contredit} / la rend impossible $\rightarrow$ \eng{CAN'T}.
\item \textbf{Choisis le temps} : Situation présente $\rightarrow$ \eng{MUST / CAN'T} + BV. Situation passée (cause antérieure) $\rightarrow$ \eng{MUST / CAN'T} + \eng{HAVE} + \textbf{Participe Passé}.
\item \textbf{Interdit} : \eng{MUSTN'T} pour la déduction négative (c'est une interdiction, pas une déduction).
\end{enumerate}
\end{methode}

\begin{exoresolu}[Déductions logiques : indices du quotidien]
\textbf{Énoncé.} Make logical deductions (\emph{must / can't} + appropriate infinitive).\\
1. The classroom is empty and the lights are off. The students \ldots\ (be) in class.\\
2. Look at that huge house and the 4x4 parked outside. The owner \ldots\ (be) very rich.\\
3. My neighbour played music all night. He \ldots\ (sleep) well this morning.\\
4. The cocoa harvest was excellent this year. Farmers \ldots\ (earn) good money.\\
5. You have been running for an hour. You \ldots\ (be) tired.
\tcblower
\textbf{Solution détaillée.}
\begin{enumerate}
\item \eng{The students \textbf{can't be} in class.} — Indice : salle vide, lumières éteintes $\rightarrow$ Impossibilité qu'ils soient là. Présent : \eng{can't be}.
\item \eng{The owner \textbf{must be} very rich.} — Indices : grande maison, 4x4 $\rightarrow$ Certitude positive. Présent : \eng{must be}.
\item \eng{He \textbf{can't have slept} well this morning.} — Cause passée (bruit toute la nuit) $\rightarrow$ Conséquence passée négative. \eng{can't have} + V-en.
\item \eng{Farmers \textbf{must have earned} good money.} — Cause passée (récolte excellente) $\rightarrow$ Conséquence passée positive. \eng{must have} + V-en.
\item \eng{You \textbf{must be} tired.} — Indice physique (courir 1h) $\rightarrow$ Certitude positive état présent. \eng{must be}.
\end{enumerate}
\end{exoresolu}

\courssub{Comparaison français/anglais, synthèse et entraînement}

\begin{propriete}[Tableau récapitulatif : Modaux vs Français — Pièges à éviter]
\begin{center}
\begin{tabularx}{\linewidth}{|X|X|X|}
\hline
\rowcolor{poppinkL} \textbf{Français} & \textbf{Anglais correct} & \textbf{Erreur fréquente (Francophonisme)} \\ \hline
« Il doit venir » (obligation) & \eng{He \textbf{must} come.} & \eng{*He \textbf{musts} come.} (pas de \eng{-s}) \\ \hline
« Il doit venir » (déduction) & \eng{He \textbf{must} come.} (contextuel) & — \\ \hline
« Il ne doit pas venir » (interdiction) & \eng{He \textbf{mustn't} come.} & \eng{*He \textbf{don't have to} come.} (sens inverse !) \\ \hline
« Il n'est pas obligé de venir » & \eng{He \textbf{needn't} come.} / \eng{He \textbf{doesn't have to} come.} & \eng{*He \textbf{mustn't} come.} (interdit au lieu de dispense) \\ \hline
« Il ne peut pas être là » (déduction) & \eng{He \textbf{can't} be there.} & \eng{*He \textbf{mustn't} be there.} \\ \hline
« Il a dû partir » (déduction passée) & \eng{He \textbf{must have} left.} & \eng{*He \textbf{must} left.} (manque \eng{have}) \\ \hline
« Il a pu partir » (permission passée) & \eng{He \textbf{was allowed to} leave.} & \eng{*He \textbf{could} leave.} (\eng{could} = capacité/générique, pas permission accordée ponctuelle) \\ \hline
« Tu devrais y aller » & \eng{You \textbf{should} go.} / \eng{You \textbf{ought to} go.} & \eng{*You \textbf{should to} go.} (parasite \eng{to}) \\ \hline
« Peux-tu m'aider ? » (politesse) & \eng{\textbf{Could} you help me?} & \eng{* \textbf{Can} you help me?} (trop familier pour un adulte/inconnu) \\ \hline
\end{tabularx}
\end{center}
\end{propriete}

\begin{savaistu}[Culture : La politesse au Cameroun anglophone]
Dans les régions du Nord-Ouest et du Sud-Ouest (ex. Buea, Bamenda), l'usage de \eng{May I\ldots?}, \eng{Could you\ldots?}, \eng{Would you mind\ldots?} est la norme absolue face aux aînés, enseignants, autorités. L'usage direct de \eng{Can I\ldots?} ou \eng{Give me\ldots} peut être perçu comme un manque de respect (\eng{rude}). En zone francophone (Yaoundé, Douala), le tutoiement ou le \eng{Can I} passe plus facilement entre jeunes, mais le \eng{vous} et le conditionnel (\eng{Pourriez-vous}) restent de rigueur en situation formelle. Maîtriser cette alternance codique est un atout social et professionnel.
\end{savaistu}

\begin{methode}[Transformer et rédiger avec les modaux — Méthode Bac]
\begin{enumerate}
\item \textbf{Négation} : Ajouter \eng{not} après le modal (sans \eng{do}). \eng{She can swim} $\rightarrow$ \eng{She \textbf{cannot} / \textbf{can't} swim.}
\item \textbf{Interrogation} : Inversion Sujet/Modal (sans \eng{do}). \eng{You may go} $\rightarrow$ \eng{\textbf{May} you go?} (rare) / \eng{\textbf{May} I go?}
\item \textbf{Reformulation (Paraphrase)} :
\begin{itemize}
\item \eng{should} $\leftrightarrow$ \eng{ought to} (conserver \eng{to}).
\item \eng{must} (obligation) $\leftrightarrow$ \eng{have to} (conjugable).
\item \eng{needn't} $\leftrightarrow$ \eng{don't have to} / \eng{don't need to}.
\item \eng{can't} (déduction) $\leftrightarrow$ \eng{It is impossible that} + subjonctif/indicatif.
\end{itemize}
\item \textbf{Production écrite (Paragraph)} : Sujet imposé (Règles familiales, Règlement scolaire, Conseils santé). Contraintes : 5 phrases minimum, 4 modaux différents (\eng{must, mustn't, should, can't, needn't, don't have to}), connecteurs logiques (\eng{however, therefore, besides, finally}).
\end{enumerate}
\end{methode}

\begin{exoresolu}[Transformations et mini-composition type Bac]
\textbf{Énoncé.} (a) Transform as indicated.\\
1. He can speak French fluently. (Negative)\\
2. You should respect the traffic lights. (Question with \emph{I})\\
3. We needn't buy tickets; entry is free. (Use \emph{don't have to})\\
4. The baby is crying. He must be hungry. (Negative deduction)\\
5. Students must wear uniforms. (Past tense)\\
(b) Write a paragraph (50--60 words) on: “Rules in my secondary school”. Use at least 5 different modals.
\tcblower
\textbf{Solution détaillée.}\\
\textbf{(a) Transformations :}
\begin{enumerate}
\item \eng{He \textbf{cannot} / \textbf{can't} speak French fluently.}
\item \eng{\textbf{Should} I respect the traffic lights?} (Adaptation du possessif si nécessaire : \eng{my}).
\item \eng{We \textbf{don't have to} buy tickets; entry is free.} (Attention : \eng{do} + \eng{have to}).
\item \eng{The baby is crying. He \textbf{can't be} hungry.} (Déduction négative $\rightarrow$ \eng{can't}, jamais \eng{mustn't}).
\item \eng{Students \textbf{had to} wear uniforms.} (\eng{Must} n'a pas de passé $\rightarrow$ \eng{had to}).
\end{enumerate}
\textbf{(b) Modèle de rédaction :}\\
“In my secondary school, discipline is strict. All students \textbf{must} wear the official uniform and \textbf{mustn't} use mobile phones in class. We \textbf{ought to} arrive before 7:30 a.m., otherwise we \textbf{can't} enter the gate. However, on Wednesdays, we \textbf{don't have to} attend afternoon classes. I think these rules \textbf{should} be respected because they prepare us for professional life.”
\textbf{Traduction :} « Dans mon lycée, la discipline est stricte. Tous les élèves doivent porter l'uniforme officiel et il est interdit d'utiliser les portables en classe. Nous devrions arriver avant 7h30, sinon nous ne pouvons pas entrer. Cependant, le mercredi, nous n'avons pas cours l'après-midi. Je pense que ces règles devraient être respectées car elles nous préparent à la vie professionnelle. »
\textbf{Justification :} 6 modaux (\eng{must, mustn't, ought to, can't, don't have to, should}), structure \eng{modal + BV} respectée, connecteurs (\eng{however, otherwise, because}), vocabulaire précis (\eng{official uniform, attend, prepare for}).
\end{exoresolu}

\begin{exercice}[Entraînement personnel — À rendre pour correction]
\textbf{Section A : Choix multiple.} Choose the best option.\\
1. “\ldots\ I speak to the Principal, please?” (Can / May / Must) — Formal context.\\
2. You \ldots\ drive without a licence. It's the law. (mustn't / don't have to / needn't)\\
3. “Where is my father?” “He \ldots\ be in the garden; I saw his hat there.” (must / can't / should)\\
4. You \ldots\ have told him the truth; he is your best friend. (should / must / could)\\
5. When I was young, I \ldots\ walk 10 km to school every day. (must / had to / have to)\\
\\
\textbf{Section B : Complétion.} Fill in the gaps with a suitable modal (\emph{can, could, may, must, mustn't, should, ought to, needn't, don't have to, can't}) + correct verb form.\\
1. \textbf{(Polite request)} \ldots\ I borrow your dictionary for a moment?\\
2. \textbf{(Prohibition)} Passengers \ldots\ cross the railway line.\\
3. \textbf{(Absence of obligation)} Tomorrow is Sunday, so we \ldots\ wake up early.\\
4. \textbf{(Advice)} You \ldots\ see a dentist; your toothache is getting worse.\\
5. \textbf{(Deduction present)} The lights are off. They \ldots\ (be) at home.\\
6. \textbf{(Deduction past)} He \ldots\ (lose) his way; he usually knows this road perfectly.\\
7. \textbf{(Past obligation)} My father \ldots\ (pay) a fine for speeding last month.\\
8. \textbf{(Possibility)} Be careful! The floor \ldots\ (be) slippery after the rain.\\
\\
\textbf{Section C : Expression écrite.} Write a short letter (80--100 words) to your younger brother who just entered Form 1. Give him advice on school rules, homework, behaviour with teachers, and free time. Use at least 6 different modal forms.
\end{exercice}

\begin{corrige}[Exercice Lesson 2]
\textbf{Section A :}
\begin{enumerate}
\item \eng{May} (Formel).
\item \eng{mustn't} (Interdiction légale).
\item \eng{must} (Déduction positive : indice = chapeau vu).
\item \eng{should} (Conseil moral / regret passé : \eng{should have told} — \textit{Note: l'énoncé dit "You ... have told", la réponse attendue est "should" pour former "should have told"}). \textbf{Correction :} L'énoncé "You \ldots\ have told him the truth" attend \eng{should} (ou \eng{ought to}) pour former \eng{should have told} (reproche/conseil passé). \eng{Must have told} serait déduction. \eng{Could have told} possibilité passée.
\item \eng{had to} (Obligation passée).
\end{enumerate}
\textbf{Section B :}
\begin{enumerate}
\item \eng{May I borrow your dictionary for a moment?} (ou \eng{Could I}).
\item \eng{Passengers \textbf{mustn't} cross the railway line.}
\item \eng{Tomorrow is Sunday, so we \textbf{don't have to} / \textbf{needn't} wake up early.}
\item \eng{You \textbf{should} / \textbf{ought to} see a dentist\ldots}
\item \eng{The lights are off. They \textbf{can't be} at home.}
\item \eng{He \textbf{can't have lost} his way\ldots} (Impossibilité logique passée).
\item \eng{My father \textbf{had to pay} a fine\ldots} (Obligation passée).
\item \eng{Be careful! The floor \textbf{could be} / \textbf{may be} slippery\ldots} (Possibilité/hypothèse).
\end{enumerate}
\textbf{Section C :} \textit{Production libre. Vérification attendue : structure modal + BV, variété des modaux (must, mustn't, should, ought to, can't, don't have to, needn't, have to), cohérence du ton (fraternel, conseillant), connecteurs, orthographe.}
\end{corrige}

\begin{attention}[Check-list finale avant l'évaluation]
\begin{itemize}
\item[$\square$] Je ne mets \textbf{jamais} de \eng{-s} à la 3\textsuperscript{e} personne du modal (\eng{He must}, pas \eng{He musts}).
\item[$\square$] Je ne mets \textbf{jamais} \eng{to} après un modal (sauf \eng{ought TO}, \eng{have TO}, \eng{need TO}).
\item[$\square$] Je distingue \textbf{\eng{mustn't} (interdit)} de \textbf{\eng{don't have to} (pas obligé)}.
\item[$\square$] Pour la déduction \textbf{négative}, j'utilise \textbf{\eng{can't}} (pas \eng{mustn't}).
\item[$\square$] Au \textbf{passé}, \eng{must} devient \eng{had to} (obligation) ou \eng{must have + V-en} (déduction).
\item[$\square$] \eng{Needn't} (modal) n'a pas de passé $\rightarrow$ \eng{didn't need to} / \eng{didn't have to}.
\item[$\square$] \eng{Could} = politesse (demande) OU hypothèse (possibilité) OU capacité passée. Le contexte décide.
\end{itemize}
\end{attention}

\newpage

\coursec{Exercices}

\courssub{Je vérifie mes connaissances}

\begin{exercice}[QCM : le bon modal]
Pour chaque phrase, choisis la réponse correcte (a, b, c ou d).
\begin{enumerate}
\item You \cle{\dots} smoke in the hospital. It is forbidden.\\
(a) don't have to \quad (b) mustn't \quad (c) needn't \quad (d) couldn't
\item \cle{\dots} I use your phone, please? Mine is dead.\\
(a) Must \quad (b) Ought \quad (c) May \quad (d) Should
\item You look exhausted. You \cle{\dots} go to bed earlier.\\
(a) mustn't \quad (b) should \quad (c) can't \quad (d) may
\item It's Saturday tomorrow, so we \cle{\dots} wear our school uniform.\\
(a) mustn't \quad (b) can't \quad (c) don't have to \quad (d) ought not to
\item Listen! Someone is knocking at the door. It \cle{\dots} be the postman; he always comes at this time.\\
(a) may not \quad (b) must \quad (c) shouldn't \quad (d) needn't
\end{enumerate}
\end{exercice}

\begin{exercice}[Vrai ou faux ? Justifie]
Dis si chaque affirmation est vraie (V) ou fausse (F), puis justifie ta réponse par la règle de grammaire.
\begin{enumerate}
\item \eng{He cans swim very well.} est une phrase correcte.
\item Après un modal, le verbe qui suit est toujours à la base verbale, sans \eng{to} (sauf pour \eng{ought to}).
\item \eng{Mustn't} et \eng{don't have to} ont exactement le même sens.
\item \eng{May} est plus formel que \eng{can} pour demander la permission.
\item \eng{Must} peut exprimer une déduction logique : \eng{She must be tired.}
\end{enumerate}
\end{exercice}

\begin{exercice}[Vocabulaire : permission, conseil ou obligation ?]
Classe les modaux suivants dans le tableau selon leur fonction : \eng{can}, \eng{may}, \eng{should}, \eng{ought to}, \eng{must}, \eng{mustn't}, \eng{needn't}, \eng{don't have to}.
\begin{center}
\begin{tabularx}{\linewidth}{|X|X|X|}
\hline
\rowcolor{popblueL} Permission / Possibilité & Conseil & Obligation / Interdiction / Absence d'obligation \\
\hline
 & & \\
\hline
\end{tabularx}
\end{center}
\end{exercice}

\begin{exercice}[Conjugaison : la bonne forme du verbe]
Complète chaque phrase avec la forme correcte du verbe entre parenthèses.
\begin{enumerate}
\item You must \dots (to be) polite to your elders.
\item She can \dots (to speak) three languages.
\item They ought to \dots (to help) their mother in the kitchen.
\item He needn't \dots (to come) if he is busy.
\item We should \dots (to respect) the school rules.
\end{enumerate}
\end{exercice}

\courssub{Je m'entraîne}

\begin{methode}[Bien choisir son modal]
\begin{enumerate}
\item Repère la \textbf{fonction} demandée par le contexte : permission, possibilité, conseil, obligation, interdiction, absence d'obligation ou déduction.
\item Choisis le modal correspondant : \eng{can/may} (permission), \eng{can/could} (possibilité), \eng{should/ought to} (conseil), \eng{must} (obligation ou déduction), \eng{mustn't} (interdiction), \eng{needn't/don't have to} (absence d'obligation), \eng{can't} (impossibilité ou déduction négative).
\item Vérifie la \textbf{construction} : modal + base verbale, sans \eng{to} (sauf \eng{ought to}), sans \eng{-s}, sans \eng{do} à la négation et à l'interrogation.
\end{enumerate}
\end{methode}

\begin{exercice}[Phrases à trous : la vie scolaire et familiale]
Complète les dix phrases suivantes avec le modal qui convient : \eng{can}, \eng{could}, \eng{may}, \eng{should}, \eng{ought to}, \eng{must}, \eng{mustn't}, \eng{needn't}, \eng{don't have to}, \eng{can't}.
\begin{enumerate}
\item In my family, children \dots greet visitors when they arrive. It is a rule.
\item \dots I go out with my friends tonight, Mum? I have finished my homework.
\item You \dots talk loudly in the library. Other students are reading.
\item My little brother \dots go to school on Sundays, so he stays at home.
\item You have a bad cough. You \dots see a doctor.
\item When I was five, I \dots not read, but now I read very well.
\item You \dots worry about the washing-up; I will do it myself.
\item We \dots to respect our parents; it is our duty.
\item The lights are on in Mr Etoundi's house. He \dots be at home.
\item That \dots be Sarah! She is in Buea this week.
\end{enumerate}
\end{exercice}

\begin{exercice}[Corrige les erreurs]
Chaque phrase contient une faute typique des francophones. Réécris-la correctement.
\begin{enumerate}
\item He cans swim very fast.
\item You must to go to school every day.
\item You ought help your sister with her homework.
\item Do you can come to my party on Saturday?
\item She doesn't must pay; the entrance is free.
\item They should to study harder for the exam.
\item May you to lend me your dictionary?
\item I needn't to bring my books today.
\end{enumerate}
\end{exercice}

\begin{attention}[Les pièges à éviter]
\begin{itemize}
\item Jamais de \eng{-s} : \eng{he cans} est impossible ; dites \eng{he can}.
\item Jamais de \eng{to} après \eng{can}, \eng{must}, \eng{should}, \eng{may}, \eng{needn't} ; seul \eng{ought} prend \eng{to}.
\item Jamais d'auxiliaire \eng{do} avec un modal : \eng{Do you can...?} est faux ; dites \eng{Can you...?}.
\item \eng{Mustn't} n'est pas la négation de l'obligation au sens « pas obligé » : pour « tu n'es pas obligé », utilise \eng{you don't have to} ou \eng{you needn't}.
\end{itemize}
\end{attention}

\begin{exercice}[Transformation : négation et interrogation]
Mets chaque phrase à la forme négative, puis à la forme interrogative.
\begin{enumerate}
\item She can cook ndolé very well.
\item They must leave early tomorrow.
\item He ought to apologise to his father.
\item You may enter the classroom now.
\end{enumerate}
\end{exercice}

\begin{exercice}[Reformulation avec le modal imposé]
Réécris chaque phrase en utilisant le modal indiqué entre parenthèses, sans changer le sens.
\begin{enumerate}
\item It is not necessary to bring a gift. (\eng{needn't})
\item I am sure he is tired after the long journey. (\eng{must})
\item It is forbidden to use phones during the exam. (\eng{mustn't})
\item It would be a good idea to visit your grandmother. (\eng{should})
\item Is it possible for me to ask a question? (\eng{may})
\item It is impossible that this boy is eighteen; he looks so young! (\eng{can't})
\end{enumerate}
\end{exercice}

\courssub{Je me prépare au Bac}

\begin{textebox}[Family rules in a changing world]
In many Cameroonian families, parents set clear rules for their children. “You must greet your elders, and you mustn't answer back,” Mrs Ngono tells her three children every morning before school. Her eldest son, Junior, is sixteen. He thinks some rules are too strict. “I should be free to choose my friends,” he says. “I ought to have more freedom now that I am older.”

His mother disagrees, but she listens. “You may go out on Saturdays,” she replies, “but you must be back before seven o'clock. You needn't cook dinner tonight — your sister will do it — but you ought to help her with the dishes.”

Junior smiles. He knows that rules are not only restrictions; they are also a way of showing love. A child who has no rules may feel free, but he may also feel lost. As his grandmother often says, “A house without rules must be a house without peace.”
\end{textebox}

\begin{exercice}[Compréhension et langue — type Bac]
Lis le texte \emph{Family rules in a changing world} ci-dessus, puis réponds aux questions.

\textbf{A. Comprehension questions} — Answer in complete sentences.
\begin{enumerate}
\item What rule does Mrs Ngono repeat to her children every morning?
\item Why does Junior think he deserves more freedom?
\item What time must Junior be back home on Saturdays?
\item According to the grandmother, what happens in a house without rules?
\end{enumerate}

\textbf{B. Language work}
\begin{enumerate}
\item Find in the text: (a) one modal expressing obligation, (b) one modal expressing permission, (c) one modal expressing advice, (d) one modal expressing absence of obligation.
\item Turn into the negative form: \eng{You must be back before seven o'clock.}
\item Rewrite using \eng{should}: \eng{You ought to help her with the dishes.}
\end{enumerate}
\end{exercice}

\begin{exercice}[Traduction guidée — type Bac]
Traduis les phrases suivantes en anglais, en utilisant le modal approprié.
\begin{enumerate}
\item Tu devrais écouter les conseils de tes parents.
\item Tu n'es pas obligé de m'accompagner au marché si tu es fatigué.
\item Il est interdit de fumer dans les transports en commun.
\item Regarde ses notes ! Il doit être très intelligent.
\item Puis-je quitter la salle, s'il vous plaît ?
\end{enumerate}
\end{exercice}

\begin{exercice}[Composition — type Bac]
\textbf{Topic:} \eng{Rules in my family}

Write a paragraph of \textbf{80--100 words} describing the rules in your family. You must use \textbf{at least five different modal verbs} (for example: \eng{must}, \eng{mustn't}, \eng{can}, \eng{should}, \eng{ought to}, \eng{needn't}, \eng{don't have to}, \eng{may}). Say whether you find these rules fair or too strict, and justify your opinion.

\emph{Conseil de méthode : présente d'abord une règle d'obligation, puis une interdiction, une permission, un conseil et une absence d'obligation ; termine par ton avis personnel.}
\end{exercice}



\newpage

\coursec{Corrigés des exercices}

\begin{corrige}[Exercice 1 — QCM : le bon modal]
\begin{enumerate}
\item \textbf{(b) \cle{mustn't}} : le contexte \eng{It is forbidden} (« c'est interdit ») impose le modal d'interdiction. \eng{You mustn't smoke in the hospital.} « Il est interdit de fumer à l'hôpital. »
\item \textbf{(c) \cle{May}} : demande de permission polie et formelle. \eng{May I use your phone, please?} « Puis-je utiliser votre téléphone, s'il vous plaît ? » \eng{Can} serait possible mais n'est pas proposé ; \eng{must}, \eng{ought} et \eng{should} ne demandent pas la permission.
\item \textbf{(b) \cle{should}} : conseil fondé sur un constat (\eng{You look exhausted}). \eng{You should go to bed earlier.} « Tu devrais te coucher plus tôt. »
\item \textbf{(c) \cle{don't have to}} : absence d'obligation (c'est samedi, pas d'école). \eng{We don't have to wear our school uniform.} « Nous ne sommes pas obligés de porter l'uniforme. » Attention : \eng{mustn't} signifierait « interdit », ce qui est absurde ici.
\item \textbf{(b) \cle{must}} : déduction logique fondée sur une habitude (\eng{he always comes at this time}). \eng{It must be the postman.} « Ce doit être le facteur. »
\end{enumerate}
\textbf{Point de vigilance :} toujours lire la phrase entière avant de choisir : le deuxième segment (\eng{It is forbidden}, \eng{Mine is dead}, \eng{he always comes at this time}) donne l'indice décisif.
\end{corrige}

\begin{corrige}[Exercice 2 — Vrai ou faux ? Justifie]
\begin{enumerate}
\item \textbf{Faux.} Les modaux sont \textbf{invariables} : jamais de \eng{-s} à la 3\up{e} personne du singulier. Forme correcte : \eng{He can swim very well.} « Il sait nager très bien. »
\item \textbf{Vrai.} La construction est : modal + \textbf{base verbale} (infinitif sans \eng{to}) : \eng{She must go.}, \eng{They can come.} Seul \eng{ought} est suivi de \eng{to} : \eng{He ought to go.}
\item \textbf{Faux.} \eng{Mustn't} = \textbf{interdiction} (« il ne faut pas, c'est défendu ») ; \eng{don't have to} = \textbf{absence d'obligation} (« tu n'es pas obligé, ce n'est pas nécessaire »). Comparer : \eng{You mustn't tell him} (c'est interdit) / \eng{You don't have to tell him} (ce n'est pas nécessaire).
\item \textbf{Vrai.} Échelle de politesse : \eng{May I...?} (soutenu, formel) $>$ \eng{Could I...?} (poli, atténué) $>$ \eng{Can I...?} (courant, familier).
\item \textbf{Vrai.} \eng{Must} exprime ici une \textbf{déduction logique} (certitude) : \eng{She must be tired.} « Elle doit être fatiguée » (= je suis presque sûr qu'elle est fatiguée). La déduction négative se fait avec \eng{can't} : \eng{She can't be tired.} « Il est impossible qu'elle soit fatiguée. »
\end{enumerate}
\end{corrige}

\begin{corrige}[Exercice 3 — Vocabulaire : permission, conseil ou obligation ?]
\begin{center}
\begin{tabularx}{\linewidth}{|X|X|X|}
\hline
\rowcolor{popblueL} Permission / Possibilité & Conseil & Obligation / Interdiction / Absence d'obligation \\
\hline
\eng{can}, \eng{may} & \eng{should}, \eng{ought to} & \eng{must} (obligation), \eng{mustn't} (interdiction), \eng{needn't} et \eng{don't have to} (absence d'obligation) \\
\hline
\end{tabularx}
\end{center}
\textbf{Justification :} \eng{can} et \eng{may} servent à demander ou donner la permission (\eng{May I come in?}) et à exprimer la possibilité (\eng{It can rain in July}). \eng{Should} et \eng{ought to} donnent un conseil ou expriment le devoir moral. \eng{Must} impose une obligation, \eng{mustn't} une interdiction, et \eng{needn't} / \eng{don't have to} signalent que l'action n'est pas nécessaire.
\end{corrige}

\begin{corrige}[Exercice 4 — Conjugaison : la bonne forme du verbe]
\textbf{Règle appliquée :} après un modal, on utilise toujours la \textbf{base verbale} (l'infinitif sans \eng{to}) ; \eng{ought} est le seul modal suivi de \eng{to} + base verbale.
\begin{enumerate}
\item You must \textbf{be} polite to your elders. « Tu dois être poli avec tes aînés. »
\item She can \textbf{speak} three languages. « Elle sait parler trois langues. »
\item They ought to \textbf{help} their mother in the kitchen. « Ils devraient aider leur mère à la cuisine. »
\item He needn't \textbf{come} if he is busy. « Il n'est pas obligé de venir s'il est occupé. »
\item We should \textbf{respect} the school rules. « Nous devrions respecter le règlement de l'école. »
\end{enumerate}
\textbf{Point de vigilance :} ne jamais écrire \eng{must to be}, \eng{can to speak} ou \eng{should to respect} ; et ne jamais oublier le \eng{to} de \eng{ought to}.
\end{corrige}

\begin{corrige}[Exercice 5 — Phrases à trous : la vie scolaire et familiale]
\begin{enumerate}
\item \textbf{must} : règle familiale (obligation). \eng{In my family, children must greet visitors when they arrive.} « Chez moi, les enfants doivent saluer les visiteurs. »
\item \textbf{Can} ou \textbf{May} : demande de permission ; \eng{May} est plus poli. \eng{May I go out with my friends tonight, Mum?} « Maman, puis-je sortir avec mes amis ce soir ? »
\item \textbf{mustn't} : interdiction (on dérange les lecteurs). \eng{You mustn't talk loudly in the library.}
\item \textbf{doesn't have to} : absence d'obligation ; attention à la 3\up{e} personne : \eng{doesn't}, pas \eng{don't}. \eng{My little brother doesn't have to go to school on Sundays.}
\item \textbf{should} : conseil (santé). \eng{You should see a doctor.} « Tu devrais voir un médecin. »
\item \textbf{could} : capacité au passé. \eng{When I was five, I could not read.} « À cinq ans, je ne savais pas lire. »
\item \textbf{needn't} : inutilité. \eng{You needn't worry about the washing-up.} « Tu n'as pas à te soucier de la vaisselle. »
\item \textbf{ought to} : devoir moral ; \eng{ought to} + base verbale. \eng{We ought to respect our parents.}
\item \textbf{must} : déduction logique (les lumières sont allumées). \eng{He must be at home.} « Il doit être à la maison. »
\item \textbf{can't} : déduction négative (elle est à Buea, donc impossible). \eng{That can't be Sarah!} « Ce ne peut pas être Sarah ! »
\end{enumerate}
\end{corrige}

\begin{corrige}[Exercice 6 — Corrige les erreurs]
\begin{enumerate}
\item \eng{He can swim very fast.} — les modaux ne prennent jamais de \eng{-s}.
\item \eng{You must go to school every day.} — jamais de \eng{to} après \eng{must}.
\item \eng{You ought to help your sister with her homework.} — \eng{ought} exige \eng{to} devant la base verbale.
\item \eng{Can you come to my party on Saturday?} — on inverse le modal et le sujet ; jamais d'auxiliaire \eng{do} avec un modal.
\item \eng{She doesn't have to pay; the entrance is free.} — \eng{must} n'a pas de négation en \eng{doesn't must} ; l'absence d'obligation se dit \eng{doesn't have to} (ou \eng{needn't}).
\item \eng{They should study harder for the exam.} — jamais de \eng{to} après \eng{should}.
\item \eng{Could you lend me your dictionary?} (ou \eng{May I borrow your dictionary?}) — le modal est suivi directement de la base verbale, sans \eng{to} ; \eng{May you lend...} n'est pas idiomatique pour une demande adressée à l'interlocuteur.
\item \eng{I needn't bring my books today.} — \eng{needn't} modal est suivi directement de la base verbale, sans \eng{to}.
\end{enumerate}
\textbf{Point de vigilance :} ces huit fautes viennent toutes du français (« il peut\textbf{s} », « tu dois \emph{de} partir », « est-ce que tu peux ? »). En anglais, le modal se suffit à lui-même : pas de \eng{-s}, pas de \eng{to} (sauf \eng{ought to}), pas de \eng{do}.
\end{corrige}

\begin{corrige}[Exercice 7 — Transformation : négation et interrogation]
\textbf{Rappel :} la négation se forme avec modal + \eng{not} ; l'interrogation par inversion modal + sujet. Jamais de \eng{do}.
\begin{enumerate}
\item Négation : \eng{She cannot (can't) cook ndolé very well.} — Interrogation : \eng{Can she cook ndolé very well?}
\item Négation : \eng{They must not (mustn't) leave early tomorrow.} (sens : interdiction) — Interrogation : \eng{Must they leave early tomorrow?}
\item Négation : \eng{He ought not to apologise to his father.} — Interrogation : \eng{Ought he to apologise to his father?} (forme soutenue ; à l'oral on préfère \eng{should}).
\item Négation : \eng{You may not enter the classroom now.} — Interrogation : \eng{May you enter the classroom now?} (grammatical ; à la 1\up{re} personne, \eng{May I enter the classroom now?} est la forme naturelle pour demander la permission).
\end{enumerate}
\textbf{Point de vigilance :} pour \eng{ought to}, la négation porte sur \eng{ought} : \eng{ought \textbf{not} to}, et non \eng{ought to not}.
\end{corrige}

\begin{corrige}[Exercice 8 — Reformulation avec le modal imposé]
\begin{enumerate}
\item \eng{You needn't bring a gift.} « Tu n'es pas obligé d'apporter un cadeau. »
\item \eng{He must be tired after the long journey.} « Il doit être fatigué après ce long voyage » (déduction).
\item \eng{You mustn't use phones during the exam.} « Il est interdit d'utiliser les téléphones pendant l'examen. »
\item \eng{You should visit your grandmother.} « Tu devrais rendre visite à ta grand-mère. »
\item \eng{May I ask a question?} « Puis-je poser une question ? »
\item \eng{This boy can't be eighteen; he looks so young!} « Ce garçon ne peut pas avoir dix-huit ans ; il a l'air si jeune ! »
\end{enumerate}
\textbf{Point de vigilance :} vérifier que le sens est conservé : \eng{needn't} garde l'idée de « pas nécessaire » ; \eng{must} de certitude ; \eng{mustn't} d'interdiction ; \eng{can't} d'impossibilité.
\end{corrige}

\begin{corrige}[Exercice 9 — Compréhension et langue — type Bac]
\textbf{A. Comprehension}
\begin{enumerate}
\item \eng{She repeats that they must greet their elders and they mustn't answer back.}
\item \eng{He thinks he deserves more freedom because he is older now (he is sixteen).}
\item \eng{He must be back before seven o'clock.}
\item \eng{According to the grandmother, a house without rules must be a house without peace.}
\end{enumerate}
\textbf{B. Language work}
\begin{enumerate}
\item (a) obligation : \eng{must} ; (b) permission : \eng{may} ; (c) conseil : \eng{should} ou \eng{ought to} ; (d) absence d'obligation : \eng{needn't}.
\item \eng{You mustn't be back before seven o'clock.} (transformation mécanique demandée ; noter que le sens change : l'obligation devient interdiction).
\item \eng{You should help her with the dishes.}
\end{enumerate}
\textbf{Barème indicatif (sur 10) :} A : 4 $\times$ 1 pt (phrase complète et information exacte) ; B1 : 4 $\times$ 0,5 pt ; B2 : 1 pt ; B3 : 1 pt ; qualité de la langue : 2 pts.
\textbf{Points de vigilance :} répondre par des phrases complètes (sujet + verbe) ; ne pas recopier le texte sans le comprendre ; citer les modaux exactement comme ils apparaissent.
\end{corrige}

\begin{corrige}[Exercice 10 — Traduction guidée — type Bac]
\begin{enumerate}
\item \eng{You should listen to your parents' advice.} (ou \eng{You ought to listen to...}) — attention : \eng{advice} est indénombrable, jamais \eng{advices}.
\item \eng{You don't have to come with me to the market if you are tired.} (ou \eng{You needn't come with me...}).
\item \eng{You mustn't smoke on public transport.} (ou \eng{Smoking is forbidden on public transport.}).
\item \eng{Look at his marks! He must be very intelligent.} — déduction logique ; \eng{must be} + adjectif.
\item \eng{May I leave the room, please?} — demande de permission formelle.
\end{enumerate}
\textbf{Barème indicatif (sur 10) :} 2 pts par phrase : 1 pt pour le choix du modal, 0,5 pt pour la construction (modal + base verbale), 0,5 pt pour le reste de la phrase.
\textbf{Points de vigilance :} « tu n'es pas obligé » se traduit par \eng{don't have to} ou \eng{needn't}, jamais par \eng{mustn't} ; « il doit être intelligent » (déduction) = \eng{must be}, pas \eng{has to be}.
\end{corrige}

\begin{corrige}[Exercice 11 — Composition — type Bac]
\textbf{Modèle de production (à adapter) :}

\eng{In my family, there are a few clear rules. First, we must greet our parents and our elders every morning; it is a sign of respect. Then, we mustn't go out on school nights, because we need time to study and to sleep. However, on Saturdays, I can visit my friends if I have finished my homework. My mother says I should always tell her where I am going, and I think she is right. Luckily, I don't have to cook during the week, but I ought to help with the dishes at weekends. In my opinion, these rules are fair: they teach us discipline, but they also leave us some freedom.}

(Environ 100 mots ; six modaux différents : \eng{must}, \eng{mustn't}, \eng{can}, \eng{should}, \eng{don't have to}, \eng{ought to}.)

\textbf{Barème indicatif (sur 10) :} respect de la consigne (80--100 mots, au moins cinq modaux différents) : 3 pts ; correction grammaticale des modaux (modal + base verbale, sans \eng{-s} ni \eng{to}) : 3 pts ; avis personnel justifié : 2 pts ; richesse du vocabulaire et orthographe : 2 pts.

\textbf{Points de vigilance :}
\begin{itemize}
\item Compter ses mots et ses modaux avant de rendre la copie.
\item Vérifier chaque modal : pas de \eng{to} (sauf \eng{ought to}), pas de \eng{-s}, pas de \eng{do}.
\item Varier les fonctions : une obligation, une interdiction, une permission, un conseil, une absence d'obligation.
\item Conclure par une opinion claire introduite par \eng{In my opinion...} ou \eng{I think that...}.
\end{itemize}
\end{corrige}

\newpage

\coursec{Fiche bilan}

\begin{popfigure}
\begin{tikzpicture}[mindmap, grow cyclic, every node/.style={concept, text=white, align=center}, root concept/.append style={concept color=popblue, font=\bfseries, minimum size=3.2cm}, level 1 concept/.append style={level distance=4.2cm, sibling angle=51.4, font=\small}, level 2 concept/.append style={level distance=2.6cm, font=\scriptsize}]
\node[root concept, align=center]{MODALS\\ can / may / must\\ should / ought to\\ need(n't)}
  child[concept color=popgreen]{ node[align=center]{Permission\\ \eng{can / may / could}} 
    child[concept color=popgreen!70]{ node{\eng{May I go out?}} }
    child[concept color=popgreen!70]{ node{\eng{You can leave.}} }
  }
  child[concept color=poporange]{ node[align=center]{Possibilité\\ \eng{can / could}}
    child[concept color=poporange!70]{ node{\eng{It could rain.}} }
  }
  child[concept color=poppurple]{ node[align=center]{Conseil\\ \eng{should / ought to}}
    child[concept color=poppurple!70]{ node{\eng{You should rest.}} }
    child[concept color=poppurple!70]{ node{\eng{You ought not to smoke.}} }
  }
  child[concept color=poppink]{ node[align=center]{Obligation\\ \eng{must / mustn't}}
    child[concept color=poppink!70]{ node{\eng{You mustn't cheat.}} }
  }
  child[concept color=popgold]{ node[align=center]{Pas d'obligation\\ \eng{needn't / don't have to}}
    child[concept color=popgold!70]{ node{\eng{You needn't come.}} }
  }
  child[concept color=popblue!80]{ node[align=center]{Déduction\\ \eng{must / can't}}
    child[concept color=popblue!60]{ node{\eng{She must be tired.}} }
    child[concept color=popblue!60]{ node{\eng{He can't be at home.}} }
  }
  child[concept color=popdark]{ node[align=center]{Forme\\ modal + BV\\ pas de \eng{to}\\ pas de \eng{-s}} };
\end{tikzpicture}
\legende{Carte mentale : les modaux de la leçon — emplois et forme.}
\end{popfigure}

\begin{bilan}
\textbf{Lexique clé (anglais / français)}
\begin{itemize}
  \item \eng{permission} : la permission ; \eng{to allow} : permettre ; \eng{to forbid} : interdire
  \item \eng{advice} (invariable !) : un conseil / des conseils ; \eng{a piece of advice} : un conseil
  \item \eng{duty} : le devoir ; \eng{compulsory} : obligatoire ; \eng{optional} : facultatif
  \item \eng{to deduce} : déduire ; \eng{obviously} : évidemment ; \eng{probably} : probablement
  \item \eng{to need} : avoir besoin ; \eng{necessary} : nécessaire
\end{itemize}

\textbf{Règles à connaître par cœur}
\begin{center}
\begin{tabularx}{\linewidth}{|l|X|X|}\hline
\rowcolor{popblueL} Modal & Emploi & Exemple \\ \hline
\eng{can / may} & permission (\eng{may} = poli, formel) & \eng{May I use your phone?} \\ \hline
\eng{can / could} & possibilité (\eng{could} = moins sûr) & \eng{It could rain this afternoon.} \\ \hline
\eng{should / ought to} & conseil & \eng{You should help your parents.} \\ \hline
\eng{ought not to} & conseil négatif & \eng{You ought not to be late.} \\ \hline
\eng{must} & obligation forte / déduction positive & \eng{You must wear a uniform.} \\ \hline
\eng{mustn't} & interdiction & \eng{You mustn't talk during the test.} \\ \hline
\eng{needn't / don't have to} & absence d'obligation & \eng{You needn't bring food.} \\ \hline
\eng{can't} & déduction négative (impossibilité logique) & \eng{He can't be in Douala — I saw him here!} \\ \hline
\end{tabularx}
\end{center}

\textbf{Forme (identique pour tous les modaux)}
\begin{itemize}
  \item modal + \cle{base verbale} (sans \eng{to}) : \eng{She can swim.} — jamais \eng{to swim}.
  \item Pas de \eng{-s} à la 3\textsuperscript{e} personne : \eng{He must go.} — jamais \eng{musts}.
  \item Négation : modal + \eng{not} (\eng{can't, mustn't, shouldn't, needn't}).
  \item Question : inversion sans auxiliaire : \eng{Can you help me?} — jamais \eng{Do you can...?}.
  \item Exception : \eng{ought \textbf{to}} garde son \eng{to} ; \eng{don't have to} se conjugue avec \eng{do}.
\end{itemize}

\textbf{Méthode express pour choisir le bon modal}
\begin{enumerate}
  \item Je demande la permission ? $\rightarrow$ \eng{Can I...?} / \eng{May I...?}
  \item J'exprime une possibilité ? $\rightarrow$ \eng{can / could}
  \item Je donne un conseil ? $\rightarrow$ \eng{should / ought to}
  \item J'exprime une obligation ? $\rightarrow$ \eng{must} ; une interdiction ? $\rightarrow$ \eng{mustn't}
  \item « Ce n'est pas obligatoire » ? $\rightarrow$ \eng{needn't / don't have to}
  \item Je déduis logiquement ? $\rightarrow$ \eng{must} (certain que oui) / \eng{can't} (certain que non)
\end{enumerate}
\end{bilan}

\begin{aretenir}[Checklist avant le Bac]
\begin{itemize}
  \item $\square$ Je sais qu'après un modal, le verbe est à la base verbale, sans \eng{to}.
  \item $\square$ Je n'ajoute jamais de \eng{-s} à un modal à la 3\textsuperscript{e} personne du singulier.
  \item $\square$ Je forme la question par inversion : \eng{Must we leave now?}
  \item $\square$ Je distingue \eng{mustn't} (interdiction) et \eng{don't have to} (pas obligé).
  \item $\square$ J'utilise \eng{may} pour une permission polie ou formelle.
  \item $\square$ Je sais que \eng{could} exprime une possibilité moins certaine que \eng{can}.
  \item $\square$ Je n'oublie pas le \eng{to} dans \eng{ought to} et sa négation \eng{ought not to}.
  \item $\square$ Pour une déduction positive, j'utilise \eng{must} : \eng{The lights are on — they must be at home.}
  \item $\square$ Pour une déduction négative, j'utilise \eng{can't} (et non \eng{mustn't}) : \eng{She can't be serious!}
  \item $\square$ Je sais que \eng{advice} est invariable : jamais \eng{advices}.
  \item $\square$ Je conjugue \eng{have to} avec l'auxiliaire \eng{do} : \eng{Do you have to work? You don't have to.}
  \item $\square$ Je peux employer ces modaux dans un court texte sur la vie familiale et sociale.
\end{itemize}
\end{aretenir}

\findecours

\end{document}
